\chapter{Canonical and counter-aware recursors}\label{ctr:chap}
\label{prov:ch:ctr}

The common core takes the successor $\Suc$ and the iterator $\It{n}{a}{u}{h}$ as term
constructors and states both computation equations as rules
(\cref{base:sec:terms,base:sec:kernel}). This chapter treats the term languages in which the
successor is an instance of the recursor, or the body of the recursor sees its counter, or
both, and in which part of the computation is moved into a canonical grammar of terms with a
normalizing substitution. \Cref{ctr:sec:terms} fixes the term languages and their normal forms.
\Cref{ctr:sec:counter} shows what a counter-aware body provides: a decoder that computes, a
selector, and a compiler for primitive recursion that uses no pairing. \Cref{ctr:sec:bridge}
gives the translations to and from iterator terms and transfers the bi-interpretation of
\cref{base:thm:bridge}. \Cref{ctr:sec:constfree} removes the constant $0$ from the register
language and defines zero by an equation. \Cref{ctr:sec:guarded} restricts the bodies of a
register with one bound variable so that two unconditional equations replace its graph
axiom, and proves the resulting native calculus equivalent to $\Nat$.
\Cref{ctr:sec:six,ctr:sec:five,ctr:sec:ce,ctr:sec:three} treat the
proof calculi stated on these syntaxes. \Cref{ctr:sec:binding} treats the concrete syntax in
which the recursor and the paired binder bind the same number of variables.
\Cref{ctr:sec:bounds} collects the boundary results.

\section{Term languages and normal forms}\label{ctr:sec:terms}

This section draws on \pkg{register\_graph}, \pkg{canonical\_recursors},
\pkg{clausal\_canonical}, \pkg{counter\_evaluation}, \pkg{uniform\_binding},
\pkg{uniform\_recursion} and \pkg{single\_constructor}.

\subsection{Three recursors without a successor constructor}

\begin{definition}[Term languages]\label{ctr:def:languages}
Each of the following grammars has one numerical sort, variables, the constant $0$ and one
binding constructor.
\begin{enumerate}
\item The \emph{seeded recursor}. Terms are $t ::= x \mid 0 \mid R(n,a;u.h)$; the variable $u$
is bound in $h$ only. Put $\Suc a \defeq R(0,a;u.0)$ and $1 \defeq \Suc 0$. The computation
equations are
\begin{align}
 R(0,a;u.h) &= \Suc a, \label{ctr:eq:r1zero}\\
 R(\Suc n,a;u.h) &= h\subst{R(n,a;u.h)}{u}. \label{ctr:eq:r1succ}
\end{align}
\item The \emph{counter-aware seeded recursor}. Terms are
$t ::= x \mid 0 \mid R(n,a;i,u.h)$; the variables $i,u$ are bound in $h$ only. The first
receives the old counter and the second the preceding state. Put
$\Suc a \defeq R(0,a;i,u.0)$ and $1 \defeq \Suc 0$. The computation equations are
\begin{align}
 R(0,a;i,u.h) &= \Suc a, \label{ctr:eq:r2zero}\\
 R(\Suc n,a;i,u.h) &= h[n/i,\,R(n,a;i,u.h)/u], \label{ctr:eq:r2succ}
\end{align}
with simultaneous capture-avoiding substitution on the right.
\item The \emph{register iterator}. Terms are $t ::= x \mid 0 \mid J(n;i,r.h)$; the variables
$i,r$ are bound in $h$ only. Put $\Suc a \defeq J(a;i,r.0)$ and $1\defeq \Suc 0$. The term
$J(n;i,r.h)$ executes $n$ instructions on a register that starts at $1$. The body receives
the old counter $i$ and the register $r$ and returns an instruction: the value $0$ increments
the register, and a value $\Suc t$ loads $t$. The axioms are stated in \cref{ctr:def:TJ}.
\end{enumerate}
The number of bound variables distinguishes the first two constructors, so both are written
$R$. In every language the variables of a body other than the bound ones are passive
numerical parameters, and substitution is capture avoiding.
\end{definition}

In all three languages the successor is a defined operation, not a constructor. The
initializer of the seeded recursors is an argument, but the iteration starts at its
successor; this is what allows the successor to be an instance of the recursor. The register
iterator has no initializer argument at all.

The binding notation is a notation for finite descriptions in a first-order language, as in
\cref{base:sec:terms}. Each fixed body, built from earlier descriptions, determines one
function symbol whose arguments are the counter, the initializer where there is one, and the
passive parameters of the body. No variable ranges over bodies, and no evaluator for
descriptions is postulated inside any theory of this chapter.

\subsection{Zero-body canonical terms}

The equations \eqref{ctr:eq:r1zero} and \eqref{ctr:eq:r2zero} say that a recursor with
counter $0$ does not depend on its body. A canonical grammar records this by one spelling.

\begin{definition}[Zero-body normalization]\label{ctr:def:nu}
For raw terms of either seeded language define $\nu$ by $\nu(x)=x$, $\nu(0)=0$ and
\[
 \nu\bigl(R(n,a;\vec u.h)\bigr)=
 \begin{cases}
  R(0,\nu a;\vec u.0), & n\equiv 0,\\
  R(\nu n,\nu a;\vec u.\nu h), & n\not\equiv 0,
 \end{cases}
\]
where $\vec u$ is $u$ or $i,u$ and $\equiv$ is literal identity of syntax. It is not a test
for provable equality with $0$. A term is \emph{zero-body canonical} if it is a fixed point of
$\nu$, that is, if every subterm with literal counter $0$ has literal body $0$. The map $\nu$
extends to formulas, clauses and contexts by normalizing all terms. For zero-body canonical
$T$ and a substitution $\sigma$ the \emph{normalized substitution} is
\[
 T\langle\sigma\rangle \defeq \nu\bigl(T[\sigma]\bigr).
\]
\end{definition}

The map $\nu$ is a terminating traversal of a finite tree. It preserves the root constructor
of a term and of a formula, so
\begin{equation}\label{ctr:eq:nuroot}
 \nu(t)\equiv 0 \iff t\equiv 0,\qquad \nu(\Suc a)=\Suc(\nu a).
\end{equation}
It is idempotent, commutes with renaming of free and bound variables, and introduces no free
variable. It performs no successor computation.

\begin{lemma}[Substitution coherence]\label{ctr:lem:coherence}
For every raw term or formula $T$ and every capture-free simultaneous substitution $\sigma$,
\[
 \nu\bigl(T[\sigma]\bigr)=\nu\bigl(\nu(T)[\nu\circ\sigma]\bigr).
\]
Consequently $(T\langle\sigma\rangle)\langle\tau\rangle=T\langle\sigma\tau\rangle$ for
zero-body canonical $T$, where $(\sigma\tau)(x)=\sigma(x)\langle\tau\rangle$.
\end{lemma}
\begin{proof}
Induct on $T$ with bound names chosen fresh for $\sigma$. For a variable both sides are
$\nu(\sigma(x))$ by idempotence, and the case $0$ is trivial. Let $T=R(n,a;\vec u.h)$.

If $n\equiv 0$, then $n[\sigma]\equiv 0$ and both sides erase the body: the left side is
$R(0,\nu(a[\sigma]);\vec u.0)$, the right side is
$R(0,\nu(\nu(a)[\nu\circ\sigma]);\vec u.0)$, and the induction hypothesis for $a$ applies.

If $n\not\equiv 0$ and $n[\sigma]\equiv 0$, then $n$ is a variable $x$ with
$\sigma(x)\equiv 0$, because a substitution instance of a recursor node is a recursor node.
The left side erases the body. On the right, $\nu(T)=R(x,\nu a;\vec u.\nu h)$, and after the
substitution $\nu\circ\sigma$ the counter is $\nu(\sigma(x))\equiv 0$ by
\eqref{ctr:eq:nuroot}, so the outer $\nu$ erases the body as well. The initializers agree
by the induction hypothesis.

If $n[\sigma]\not\equiv 0$, then $n\not\equiv 0$, and by \eqref{ctr:eq:nuroot} neither
$\nu(n[\sigma])$ nor $\nu(n)[\nu\circ\sigma]$ is literally $0$. Both sides retain the node,
and the induction hypotheses for the counter, the initializer and the body give the
equality; for the body the substitution is extended by the identity on $\vec u$.

Formula constructors bind variables but introduce no further case. For the consequence,
apply the identity twice: $\nu(\nu(T[\sigma])[\tau])=\nu(T[\sigma][\tau])$ since $\tau$ is
canonical, and $\nu(T[\sigma\tau])=\nu(T[\sigma][\tau])$ since $T$ is canonical and
$\nu(\sigma(x)[\tau])=(\sigma\tau)(x)$.
\end{proof}

A variable replaced by $0$ may expose a zero counter. Normalizing only the original syntax,
or testing the counter before substitution, is therefore not sufficient, and rejecting
noncanonical input is not sufficient either: instances generated by a rule have to be
normalized. In a display that belongs to a canonical calculus, an expression
$R(n,a;\vec u.h)$ denotes its normalization; the convention matters only when
$n\equiv 0$, where the written body is discarded.

\begin{remark}
\Cref{ctr:lem:coherence} is the substitution-coherence lemma of \pkg{canonical\_recursors}
for the seeded recursor. \pkg{counter\_evaluation} and \pkg{uniform\_binding} state it for
the counter-aware recursor in the form $\nu(t\sigma)=\nu(\nu(t)\,\nu(\sigma))$ with the same
proof, and \pkg{clausal\_canonical} states it for the opening operation of a de Bruijn
representation (\cref{ctr:lem:opening}). The proof above covers both binding arities.
\end{remark}

\subsection{Full normal forms}

The counter-aware language also admits a grammar in which both computation equations are
built into the syntax.

\begin{definition}[Normal terms]\label{ctr:def:normal}
In the counter-aware language, a \emph{canonical successor} is a term $R(0,a;i,u.0)$ with $a$
normal. A \emph{neutral} term is a variable or a term $R(c,a;i,u.h)$ with $c$ neutral and
$a,h$ normal. A term is \emph{normal} if it is $0$, a canonical successor or neutral. Thus a
retained recursor node has a counter that is neither $0$ nor a successor.
\end{definition}

Every normal term has the form $\Suc^k(c_0)$ with $c_0\equiv 0$ or $c_0$ neutral; $k$ is the
length of its successor prefix. Every normal term is zero-body canonical.

\begin{definition}[Normalization by evaluation]\label{ctr:def:eval}
An environment $\rho$ maps variables to normal terms. Define $\operatorname{ev}(t,\rho)$ by
recursion on the raw term $t$: $\operatorname{ev}(x,\rho)=\rho(x)$,
$\operatorname{ev}(0,\rho)=0$, and
\[
 \operatorname{ev}\bigl(R(n,a;i,u.h),\rho\bigr)
 =\operatorname{run}\bigl(\operatorname{ev}(n,\rho),\operatorname{ev}(a,\rho),H\bigr),
 \qquad H(d,w)=\operatorname{ev}\bigl(h,\rho[i\mapsto d,u\mapsto w]\bigr),
\]
where, for normal $c,v$,
\begin{equation}\label{ctr:eq:run}
 \operatorname{run}(c,v,H)=
 \begin{cases}
  \Suc v, & c\equiv 0,\\
  H\bigl(d,\operatorname{run}(d,v,H)\bigr), & c\equiv \Suc d,\\
  R\bigl(c,v;i,u.H(i,u)\bigr), & c \text{ neutral},
 \end{cases}
\end{equation}
with $i,u$ fresh in the third clause. The notation $H(i,u)$ is evaluation of the body in the
environment extended by the two fresh variables; it introduces no function object into the
term language. Put $\nf(t)=\operatorname{ev}(t,\mathrm{id})$, and for formulas, clauses and
substitutions normalize all terms.
\end{definition}

\begin{lemma}[Normalization and substitution]\label{ctr:lem:nf}
Evaluation is total on raw terms and normal environments, and its values are normal.
Moreover:
\begin{enumerate}
\item $\nf(t)=t$ for every normal $t$; hence $\nf$ is idempotent.
\item $\nf$ respects \eqref{ctr:eq:r2zero} and \eqref{ctr:eq:r2succ}: both sides of every
instance have the same normal form, and so do two terms that differ by an instance inside
any context, including under a binder.
\item For every raw $t$ and substitution $\sigma$,
\begin{equation}\label{ctr:eq:SC}
 \nf(t[\sigma])=\nf\bigl(\nf(t)[\nf\circ\sigma]\bigr).
\end{equation}
\end{enumerate}
The same statements hold for formulas and clauses; no formula binder is erased.
\end{lemma}
\begin{proof}
% ed.: the package single_constructor states the inductions for items 2 and 3 in outline; they are written out here.
\emph{Totality.} Induct on $t$. At a recursor node the values $c,v$ are normal and the body
evaluator $H$ is total with normal values, by the induction hypothesis for the proper
subterms. For fixed total $H$, induct on the successor prefix of $c$ in
\eqref{ctr:eq:run}: the first and third clauses return a normal term, and in the second
clause the recursive call has a shorter prefix and $H$ is total.

\emph{Item 1.} Induct on a normal term. A canonical successor evaluates to
$\operatorname{run}(0,a,H)=\Suc a$. A neutral $R(c,a;i,u.h)$ evaluates by the third clause to
$R(c,a;i,u.H(i,u))$, and $H(i,u)=h$ by the induction hypothesis.

\emph{Two environment laws.} For a substitution $\sigma$ and environments $\rho,\tau$ write
$(\rho\circ\sigma)(x)=\operatorname{ev}(\sigma(x),\rho)$ and
$(\tau\bullet\rho)(x)=\operatorname{ev}(\rho(x),\tau)$. Then
\begin{align}
 \operatorname{ev}(t[\sigma],\rho)&=\operatorname{ev}(t,\rho\circ\sigma),
 \label{ctr:eq:envsub}\\
 \operatorname{ev}\bigl(\operatorname{ev}(t,\rho),\tau\bigr)&=\operatorname{ev}(t,\tau\bullet\rho).
 \label{ctr:eq:envcomp}
\end{align}
Both are proved by induction on $t$ with bound names fresh. For \eqref{ctr:eq:envsub} the
recursor case holds because the body evaluator of $h[\sigma]$ in $\rho$ is the body evaluator
of $h$ in $\rho\circ\sigma$, by the induction hypothesis. For \eqref{ctr:eq:envcomp} let
$H$ be the body evaluator of $h$ in $\rho$ and $H^\tau$ the one in $\tau\bullet\rho$; the
induction hypothesis for $h$ gives
$\operatorname{ev}(H(d,w),\tau)=H^\tau(\operatorname{ev}(d,\tau),\operatorname{ev}(w,\tau))$.
It suffices to show
\[
 \operatorname{ev}\bigl(\operatorname{run}(c,v,H),\tau\bigr)
 =\operatorname{run}\bigl(\operatorname{ev}(c,\tau),\operatorname{ev}(v,\tau),H^\tau\bigr)
\]
by induction on the successor prefix of $c$. If $c\equiv 0$ both sides are
$\Suc(\operatorname{ev}(v,\tau))$. If $c\equiv\Suc d$ the left side is
$H^\tau(\operatorname{ev}(d,\tau),\operatorname{ev}(\operatorname{run}(d,v,H),\tau))$, which
by the hypothesis for $d$ is the second clause of \eqref{ctr:eq:run} at the counter
$\Suc(\operatorname{ev}(d,\tau))=\operatorname{ev}(\Suc d,\tau)$. If $c$ is neutral, the left
side evaluates the residual node $R(c,v;i,u.H(i,u))$ in $\tau$; its body evaluator is
$(d,w)\mapsto\operatorname{ev}(H(i,u),\tau[i\mapsto d,u\mapsto w])=H^\tau(d,w)$, so the left
side is $\operatorname{run}(\operatorname{ev}(c,\tau),\operatorname{ev}(v,\tau),H^\tau)$
whichever clause applies to $\operatorname{ev}(c,\tau)$. This is the case in which a
substitution exposes a new redex.

\emph{Item 2.} In every environment,
$\operatorname{ev}(R(0,a;i,u.h),\rho)=\Suc(\operatorname{ev}(a,\rho))$ independently of $h$,
and $\operatorname{ev}(R(\Suc n,a;i,u.h),\rho)=H(c,\operatorname{run}(c,v,H))$ with
$c=\operatorname{ev}(n,\rho)$, which by \eqref{ctr:eq:envsub} is the value of
$h[n/i,R(n,a;i,u.h)/u]$ in $\rho$. Evaluation is compositional: the value of a term in all
environments is determined by the values of its immediate subterms in all environments.
Hence replacing a subterm by one with the same values does not change the value.

\emph{Item 3.} By \eqref{ctr:eq:envsub}, $\nf(t[\sigma])=\operatorname{ev}(t,\nf\circ\sigma)$.
By \eqref{ctr:eq:envsub}, item~1 and \eqref{ctr:eq:envcomp},
$\nf(\nf(t)[\nf\circ\sigma])=\operatorname{ev}(\operatorname{ev}(t,\mathrm{id}),\nf\circ\sigma)
=\operatorname{ev}(t,\nf\circ\sigma)$.
\end{proof}

\begin{remark}\label{ctr:rem:nfvariants}
\pkg{single\_constructor} defines the normalizer by evaluation, as above.
\pkg{uniform\_recursion} obtains the same normal forms from rewriting: it reads
$R(n,a;i,u.h)$ as the base-type recursor of G\"odel's system $T$ with initial value
$\Suc a$ and update $h(i,u)$, derives strong normalization from the reducibility argument
for that recursor (a numerical expression is computable when it is strongly normalizing, an
update is computable when it maps computable arguments to computable expressions, and the
recursor case is an induction on the successor spine of a normal counter), observes that the
zero and successor rules have disjoint head patterns and commute with reductions in their
arguments and with erasure of an unused body, and concludes unique normal forms from
termination and local confluence. The two definitions agree, since by item~2 of
\cref{ctr:lem:nf} every reduction step preserves $\nf$ and by item~1 the normal terms are
exactly the irreducible ones.
\end{remark}

\begin{remark}[No uniform primitive recursive normalizer]\label{ctr:rem:fuel}
Each fixed term denotes a primitive recursive function of its parameters. It does not follow
that normalization of arbitrary term descriptions is primitive recursive uniformly; a
universal evaluator of that kind leads to the usual diagonal contradiction. The normalizer is
an external total recursive operation on finite syntax, and no assertion that all
descriptions normalize is an axiom of any theory here. \pkg{uniform\_recursion} supplies in
addition a checker with an explicit budget. Bounded simulation gives a primitive recursive
certificate predicate with the same set of provable formulas, since every finite successful
run has a sufficient budget.
\end{remark}

\subsection{The status of the successor and of each computation}

The languages of this chapter differ from the iterator terms of \cref{base:sec:terms}, and
from one another, in which of three items is a constructor or a rule and which is carried by
the syntax. The following table summarizes the presentations treated below.
\begin{center}
\small
\begin{tabular}{@{}p{0.27\linewidth}p{0.19\linewidth}p{0.22\linewidth}p{0.24\linewidth}@{}}
\toprule
Terms & Successor & Zero computation & Successor computation\\
\midrule
$0,\Suc,I$ (\cref{base:sec:terms}) & constructor & rule & rule\\
raw seeded $R$, either arity & $R(0,a;\vec u.0)$ & rule & rule\\
zero-body canonical $R$, one binder (\cref{ctr:sec:six,ctr:sec:five}) & $R(0,a;u.0)$ &
identity of spelling & the one computation rule\\
zero-body canonical $R$, two binders (\cref{ctr:sec:ce}) & $R(0,a;i,u.0)$ &
identity of spelling & derived from reflexivity and counter evaluation\\
normal $R$, two binders (\cref{ctr:sec:three}) & $R(0,a;i,u.0)$ & identity of spelling &
identity of spelling\\
register $J$ (\cref{ctr:sec:counter}) & $J(a;i,r.0)$ & axiom $J(0;i,r.h)=1$ &
graph biconditional\\
guarded scalar $J$ (\cref{ctr:sec:guarded}) & $J(a;v.0)$ & rule
$J(0;u.h)=\Suc(h\subst{0}{u})$ & rule at successor-headed bodies, derived at state-free
bodies\\
\bottomrule
\end{tabular}
\end{center}
A computation that is an identity of spelling has not left the trusted specification: the
normalizer is applied at every substitution and comparison, and it is counted in every
checker of \cref{ctr:sec:sizes}. Its soundness in ordinary arithmetic is
\cref{ctr:lem:normsound}; it rests on an extensionality property of recursor bodies
(\cref{ctr:lem:bodyext}) and not on syntactic idempotence alone.

\section{What the counter gives}\label{ctr:sec:counter}

This section draws on \pkg{counter\_evaluation}, \pkg{uniform\_binding},
\pkg{uniform\_recursion}, \pkg{single\_constructor}, \pkg{register\_graph} and, for the
contrast with a body that does not see its counter, \pkg{canonical\_recursors}.

\subsection{A decoder that computes}

In the counter-aware language define, with all auxiliary bound variables fresh,
\begin{equation}\label{ctr:eq:decoder}
\begin{aligned}
 P(n)&\defeq R(n,0;i,u.i),\\
 C(n,a,b)&\defeq P\bigl(R(n,a;i,u.\Suc b)\bigr),\\
 \operatorname{Rec}(n,a;i,u.h)&\defeq
   P\bigl(R(n,a;i,v.\Suc(h[P(v)/u]))\bigr),\\
 p(n)&\defeq C(n,0,P(n)),\qquad \operatorname{sg}(n)\defeq C(n,0,1).
\end{aligned}
\end{equation}

\begin{lemma}[Decoder, selector and recursion]\label{ctr:lem:decoder}
The following equations are derivable from \eqref{ctr:eq:r2zero}, \eqref{ctr:eq:r2succ},
reflexivity and congruence, for arbitrary terms and passive parameters:
\begin{align*}
 P(0)&=1, & P(\Suc n)&=n,\\
 C(0,a,b)&=a, & C(\Suc n,a,b)&=b,\\
 \operatorname{Rec}(0,a;i,u.h)&=a, &
 \operatorname{Rec}(\Suc n,a;i,u.h)&=h[n/i,\operatorname{Rec}(n,a;i,u.h)/u],\\
 p(0)&=0, & p(\Suc n)&=n,\\
 \operatorname{sg}(0)&=0, & \operatorname{sg}(\Suc n)&=1.
\end{align*}
No induction, numerical case distinction or logical rule is used. On normal terms each
equation is a literal identity of normal forms.
\end{lemma}
\begin{proof}
The equations for $P$ are the two computations with the body that returns the old counter
and ignores the state. In $C(0,a,b)$ the inner recursor is $\Suc a$ and in $C(\Suc n,a,b)$ it
is $\Suc b$; then apply $P(\Suc x)=x$. For recursion put
$T(n)=R(n,a;i,v.\Suc(h[P(v)/u]))$. Then $T(0)=\Suc a$ and
$T(\Suc n)=\Suc(h[n/i,P(T(n))/u])$, and applying $P$ gives both equations. The equations for
$p$ and $\operatorname{sg}$ are instances of those for $C$ and $P$. On normal terms each
step is a computation, so both sides have the same normal form by \cref{ctr:lem:nf}.
\end{proof}

The decoder $P$ is a left inverse of the successor obtained by one computation step, for an
arbitrary open argument. Its value at $0$ is $1$, not $0$; the predecessor of
\cref{base:thm:pred} is $p$. The state of $\operatorname{Rec}$ is stored as a successor and
the outer $P$ removes that wrapper after each observation, so the successor seed of $R$ does
not restrict the initial value.

\begin{remark}
\pkg{counter\_evaluation} and \pkg{uniform\_binding} use the names of
\eqref{ctr:eq:decoder}. \pkg{uniform\_recursion} writes $W$ for the decoder, $P$ for the
predecessor $p$ and $J$ for $\operatorname{Rec}$, and uses the term
$K(y)\defeq C(W(y),1,0)$, with $K(1)=1$ and $K(\Suc 1)=0$. \pkg{single\_constructor} defines
the predecessor and the selector from the recursion macro, as
$p(t)\defeq\operatorname{Rec}(t,0;i,v.i)$ and $C(t,a,b)\defeq\operatorname{Rec}(t,a;i,v.b)$;
they satisfy the same equations.
\end{remark}

\subsection{Primitive recursion without pairing}

\begin{theorem}[Direct compiler]\label{ctr:thm:compiler}
Every fixed primitive recursive description compiles to a term of the counter-aware
language whose defining equations are derivable from the computation equations and equality
reasoning alone. Zero, successor and projections are terms, composition is substitution, and
a recursion $f(0,\vec p)=g(\vec p)$, $f(\Suc n,\vec p)=h(n,f(n,\vec p),\vec p)$ compiles to
\[
 f^*(n,\vec p)=\operatorname{Rec}\bigl(n,g^*(\vec p);i,u.h^*(i,u,\vec p)\bigr).
\]
The compiler uses no pairing, no projection functions, no induction and no coverage. The
iterator of \cref{base:sec:terms} is the instance
$\It{n}{a}{u}{h}\defeq\operatorname{Rec}(n,a;i,u.h)$ with $i\notin\FV(h)$.
\end{theorem}
\begin{proof}
By external induction on the description. The recursion case is \cref{ctr:lem:decoder};
the other cases are immediate.
\end{proof}

For example $x+y=\operatorname{Rec}(y,x;i,v.\Suc v)$,
$x\cdot y=\operatorname{Rec}(y,0;i,v.v+x)$ and the truncated subtraction is
$d(x,y)=\operatorname{Rec}(y,x;i,v.p(v))$. The compiler is an external recursion on each
finite description. It is not a universal evaluator for description codes inside the theory.
The compiler of \cref{base:thm:primrec} for iterator terms stores the counter beside the
state by a pairing function and therefore comes after the construction of pairing; here the
counter is an argument of the body.

\subsection{The register iterator}\label{ctr:sec:register}

The register iterator carries the same information in a different form: the body returns an
instruction, and the axioms describe the transition graph.

\begin{definition}[The theory $\mathsf T_J$]\label{ctr:def:TJ}
The language has the terms of \cref{ctr:def:languages}(3) and the formulas
\[
 A ::= (s=t)\mid\bot\mid(A\land A)\mid(A\lor A)\mid(A\to A)\mid\forall x\,A\mid\exists x\,A .
\]
The logical rules are those of intuitionistic natural deduction with equality: hypothesis,
reflexivity, transport of an arbitrary formula along a proved equation, ex falso, and the
introduction and elimination rules of $\land,\lor,\to,\forall,\exists$ with the usual
eigenvariable conditions. There is no classical rule. For a fixed body $h(i,r,\vec p)$ write
$j_n=J(n;i,r.h)$ and $e_n=h(n,j_n,\vec p)$. The arithmetic axioms are the universal closures
of
\begin{align}
 & J(0;i,r.h)=1, \tag{J0}\label{ctr:eq:J0}\\
 & J(\Suc n;i,r.h)=b \;\leftrightarrow\;
   \bigl((e_n=0\land b=\Suc j_n)\lor e_n=\Suc b\bigr), \tag{STEP}\label{ctr:eq:STEP}\\
 & 0=1\to\bot, \tag{Sep}\label{ctr:eq:Sep}
\end{align}
where $b$ is a variable distinct from the passive parameters. The one arithmetic rule is
induction for literal equations with an arbitrary context:
\begin{equation}\label{ctr:eq:EI}
 \infrule{\Gamma\proves s(0)=t(0)\qquad \Gamma,\ s(z)=t(z)\proves s(\Suc z)=t(\Suc z)}
         {\Gamma\proves s(a)=t(a)}
 \qquad (z \text{ fresh for } \Gamma,a \text{ and the parameters}).
\end{equation}
An equation under a quantifier or a connective is not eligible as an invariant.
\end{definition}

The right-to-left direction of \eqref{ctr:eq:STEP} gives two conditional computation rules,
whose premises are proved equations and not tests performed by an evaluator:
\begin{align}
 e_n=0 &\;\Longrightarrow\; J(\Suc n;i,r.h)=\Suc j_n, \label{ctr:eq:inc}\\
 e_n=\Suc t &\;\Longrightarrow\; J(\Suc n;i,r.h)=t. \label{ctr:eq:load}
\end{align}

\begin{lemma}[Functional coverage]\label{ctr:lem:coverage}
Let $f(n,a)\defeq J(n;i,r.a)$ with $i,r$ fresh for $a$, and $W(a)\defeq f(1,a)$. Without
induction, $\mathsf T_J\proves a=0\lor a=\Suc(W(a))$. In particular
$\forall a\,(a=0\lor\exists x\,a=\Suc x)$ is derivable.
\end{lemma}
\begin{proof}
Take the constant body $a$, the counter $n=0$ and $b=f(1,a)$ in \eqref{ctr:eq:STEP}. Its
left side holds by reflexivity. Its right side is $(a=0\land b=\Suc 1)\lor a=\Suc b$; drop
the second conjunct of the first alternative.
\end{proof}

\begin{proposition}[The graph axiom]\label{ctr:prop:step}
Over intuitionistic logic with equality, and without \eqref{ctr:eq:Sep}, induction or
injectivity of the successor:
\begin{enumerate}
\item \eqref{ctr:eq:J0} and \eqref{ctr:eq:STEP} are equivalent to \eqref{ctr:eq:J0},
\eqref{ctr:eq:inc}, \eqref{ctr:eq:load} and numerical coverage
$\forall a\,(a=0\lor\exists x\,a=\Suc x)$;
\item both are equivalent to the single schema
\begin{equation}\label{ctr:eq:fullgraph}
 J(n;i,r.h)=b\;\leftrightarrow\;
 (n=0\land b=1)\lor\exists m\,\bigl(n=\Suc m\land
   ((e_m=0\land b=\Suc j_m)\lor e_m=\Suc b)\bigr).
\end{equation}
\end{enumerate}
\end{proposition}
\begin{proof}
(1) From \eqref{ctr:eq:STEP} the two conditional rules follow by its right-to-left
direction and coverage is \cref{ctr:lem:coverage}. Conversely the right-to-left direction of
\eqref{ctr:eq:STEP} is the two conditional rules. For the left-to-right direction assume
$J(\Suc n;i,r.h)=b$ and split coverage for $e_n$. If $e_n=0$, then \eqref{ctr:eq:inc} gives
$b=\Suc j_n$. If $e_n=\Suc t$, then \eqref{ctr:eq:load} gives $b=t$ and hence $e_n=\Suc b$.

(2) % ed.: the package states that the two schemas can be packaged into one graph schema without displaying it; the schema and the argument are supplied here, following the audit note of the package.
The right-to-left direction of \eqref{ctr:eq:fullgraph} gives \eqref{ctr:eq:J0} at
$n=0$, $b=1$ and both conditional rules at $n=\Suc m$. Its left-to-right direction applied
to the reflexive equation $J(n;i,r.0)=J(n;i,r.0)$ gives $n=0\lor\exists m\,n=\Suc m$. By
(1) this proves \eqref{ctr:eq:STEP}. Conversely, assume \eqref{ctr:eq:J0} and
\eqref{ctr:eq:STEP}. The right-to-left direction of \eqref{ctr:eq:fullgraph} is
\eqref{ctr:eq:J0} and the right-to-left direction of \eqref{ctr:eq:STEP} after transport
along $n=\Suc m$. For the left-to-right direction split coverage for $n$: at $n=0$ use
\eqref{ctr:eq:J0}, and at $n=\Suc m$ use \eqref{ctr:eq:STEP}.
\end{proof}

The biconditional is essential: the left-to-right direction of \eqref{ctr:eq:STEP} alone
does not assert that the computation is deterministic. A packaging such as
\eqref{ctr:eq:fullgraph} is also available for the iterator of \cref{base:sec:terms}, so a
smaller number of named schemas is not by itself a simplification.

\begin{lemma}[Bootstrap without induction]\label{ctr:lem:regboot}
Using only equality, \eqref{ctr:eq:J0} and the conditional rules, $\mathsf T_J$ proves
\begin{equation}\label{ctr:eq:f}
 f(0,a)=1,\qquad f(n,0)=\Suc n,\qquad f(\Suc n,\Suc a)=a,
\end{equation}
and for
\[
 Z(x)\defeq f(x,1),\quad W(x)\defeq f(1,x),\quad D(x)\defeq f(x,\Suc x),\quad
 P(x)\defeq W(D(x)),
\]
\begin{alignat*}{3}
 Z(0)&=1, &\qquad Z(\Suc x)&=0,\\
 W(0)&=\Suc 1, & W(\Suc x)&=x,\\
 D(0)&=1, & D(\Suc x)&=\Suc x,\\
 P(0)&=0, & P(\Suc x)&=x.
\end{alignat*}
Further, for
\[
 G(c,a)\defeq P^2 f(\Suc a,Z(c)),\qquad
 C(c,a,b)\defeq P^2 f\bigl(\Suc a,\,G(c,\Suc^3 b)\bigr),
\]
where $P^2$ and $\Suc^3$ are finite compositions,
\begin{equation}\label{ctr:eq:regsel}
 G(0,a)=0,\quad G(\Suc n,a)=a,\qquad C(0,a,b)=a,\quad C(\Suc n,a,b)=b .
\end{equation}
Consequently the successor is injective, and with \eqref{ctr:eq:Sep} no successor is $0$.
\end{lemma}
\begin{proof}
The first equation of \eqref{ctr:eq:f} is \eqref{ctr:eq:J0}, the second is the definition
of $\Suc$, and the third is a load step whose premise is the reflexive equation
$\Suc a=\Suc a$. The eight endpoint equations are instances: for example
$Z(\Suc x)=f(\Suc x,\Suc 0)=0$, $W(0)=f(1,0)=\Suc 1$, $P(0)=W(1)=W(\Suc 0)=0$ and
$P(\Suc x)=W(\Suc x)=x$. For $G$: at $c=0$ the inner term is $f(\Suc a,\Suc 0)=0$ and
$P^2(0)=0$; at $c=\Suc n$ it is $f(\Suc a,0)=\Suc^2 a$ and $P^2$ returns $a$. For $C$: at
$c=0$ the instruction $G(c,\Suc^3 b)$ is $0$, so the inner term is $\Suc^2 a$; at
$c=\Suc n$ the instruction is $\Suc^3 b$, and the counter $\Suc a$, which is visibly a
successor, loads $\Suc^2 b$. Applying $P$ to $\Suc x=\Suc y$ gives $x=y$. An equation
$\Suc x=0$ gives $Z(\Suc x)=Z(0)$, that is $0=1$.
\end{proof}

The selector $C$ returns arbitrary numerical outputs. It is not limited to Boolean results or
to values bounded by a preselected initializer; the direct use of the old register in the
increment branch is what makes this possible.

\begin{theorem}[Register compiler]\label{ctr:thm:regcompiler}
For terms $a$ and $h(u,v)$ of the register language, with $i,r$ fresh, put
\[
 H(i,r)\defeq \Suc\bigl(C(i,a,h(P(i),r))\bigr),\qquad
 \operatorname{Rec}_J(n,a;u,v.h)\defeq J(\Suc n;i,r.H(i,r)).
\]
Without induction, coverage, pairing or reflection of equality,
\[
 \operatorname{Rec}_J(0,a;u,v.h)=a,\qquad
 \operatorname{Rec}_J(\Suc n,a;u,v.h)=h\bigl(n,\operatorname{Rec}_J(n,a;u,v.h)\bigr).
\]
Hence every fixed primitive recursive description compiles to a term of the register
language with derivable defining equations.
\end{theorem}
\begin{proof}
The instruction is always a successor, so every update is a load. The first update has
counter $i=0$ and old register $1$; by \eqref{ctr:eq:regsel} its instruction is $\Suc a$ and
\eqref{ctr:eq:load} gives $J(1;i,r.H)=a$. For the second equation put
$r_0=J(\Suc n;i,r.H)$. The old counter is $\Suc n$, so \eqref{ctr:eq:regsel} and
$P(\Suc n)=n$ give $H(\Suc n,r_0)=\Suc(h(n,r_0))$, and one load gives
$J(\Suc(\Suc n);i,r.H)=h(n,r_0)$. Initial functions and composition translate by
substitution.
\end{proof}

The compact source grammar does not guarantee small expanded terms. The compiler shipped
with \pkg{register\_graph} uses the instruction $\Suc(C(Z(i),h(P(i),r),a))$, which has the
same endpoint values and a smaller expansion.

\subsection{A register with one bound variable}

\begin{proposition}[One-variable register]\label{ctr:prop:scalarreg}
Let $t ::= x\mid 0\mid J_s(n;r.h)$ with $\Suc a\defeq J_s(a;r.0)$, and let the axioms be
\eqref{ctr:eq:J0}, \eqref{ctr:eq:STEP} with $e_n=h(j_n,\vec p)$, \eqref{ctr:eq:Sep} and
\eqref{ctr:eq:EI}. Then \cref{ctr:lem:coverage,ctr:lem:regboot} hold unchanged, and for
\[
 H(v)\defeq\Suc^3\bigl(C(P(v),a,h(P^2 v))\bigr),\qquad
 K(n)\defeq J_s(\Suc n;v.H(v)),\qquad
 \widehat I(n,a;r.h)\defeq P^2K(n)
\]
the equations $\widehat I(0,a;r.h)=a$ and
$\widehat I(\Suc n,a;r.h)=h(\widehat I(n,a;r.h))$ are derivable without induction or
coverage.
\end{proposition}
\begin{proof}
The bootstrap uses constant bodies only. One load step gives, at every counter,
\[
 J_s(\Suc n;v.H)=\Suc^2\bigl(C(P(J_s(n;v.H)),a,h(P^2J_s(n;v.H)))\bigr),
\]
so every state after the initial one is visibly a double successor. At the first step the
previous register is $1$ and $P(1)=0$, so $K(0)=\Suc^2 a$. At a later step $P(K(n))$ is a
successor and the selector takes the update branch, so $K(\Suc n)=\Suc^2h(P^2K(n))$.
Applying $P$ twice gives both equations. The shape of the state is a consequence of one
computation, not an inferred invariant.
\end{proof}

The one-variable register keeps the binding arity of the iterator and still has neither a
successor constructor nor an initializer. Its simulation of general primitive recursion
passes through the iterator and therefore through the pairing of \cref{base:thm:primrec}.
The counter variable of the two-variable register is the cost of the direct compiler of
\cref{ctr:thm:regcompiler}.

\begin{proposition}[Reflexivity from conditional computation]\label{ctr:prop:regrefl}
Consider \eqref{ctr:eq:J0}, \eqref{ctr:eq:inc} and \eqref{ctr:eq:load} as rules whose
premises are derivations, together with transport of an equation along an equation. Then
$t=t$ is derivable for every term $t$, without a primitive reflexivity rule.
\end{proposition}
\begin{proof}
(i) \eqref{ctr:eq:J0} for the zero body is literally $1=1$. (ii) A load step with the
constant instruction $1=\Suc 0$, whose premise is (i), gives $J(1;i,r.1)=0$; transporting
this equation along itself gives $0=0$. (iii) An increment step with the zero body at the
counter $t$, whose premise is (ii), gives the literal equation
$\Suc(\Suc t)=\Suc(\Suc t)$, since $J(t;i,r.0)$ is $\Suc t$ by definition. (iv) A load of the
constant instruction $\Suc(\Suc t)$, whose premise is (iii), gives an equation with right
side $\Suc t$, and transport along itself gives $\Suc t=\Suc t$. (v) A load of the constant
instruction $\Suc t$ gives an equation with right side $t$, and transport along itself gives
$t=t$.
\end{proof}

By \cref{ctr:prop:regrefl} the register terms are compatible with a native kernel in the
style of \cref{base:sec:kernel}, in which a conditional computation step takes a proved
equation as premise. No checker for such a kernel is supplied (\cref{ctr:sec:questions}).

\subsection{Without a counter}

The seeded recursor with one bound variable has the partial selector
$B(n,a,b)\defeq R(n,a;u.b)$, $u\notin\FV(b)$, with $B(0,a,b)=\Suc a$ and
$B(\Suc n,a,b)=b$, but its body has no access to the counter, and the decoder of
\cref{ctr:lem:decoder} is not available. A predecessor exists nevertheless: by
\cref{seed:thm:succ-pred} there is a fixed term $P$ of the seeded language with
\begin{equation}\label{ctr:eq:seedpred}
 P(0)=0,\qquad P(\Suc a)=a
\end{equation}
derivable in the native calculus and in the ordinary theory of the seeded recursor, by
native numerical cases, equality reasoning and induction on literal equations. By
\cref{seed:thm:succ-iter} the term
\[
 \widehat I(n,a;u.h)\defeq P\bigl(R(n,a;v.\Suc(h[P(v)/u]))\bigr)
\]
then satisfies the two equations of the iterator, and $C(n,a,b)\defeq\widehat I(n,a;u.b)$ is
a selector with $C(0,a,b)=a$ and $C(\Suc n,a,b)=b$.

The difference between the two arities is therefore the length of the bootstrap and the
need for pairing, not expressive strength: with the counter the predecessor is one
computation step and primitive recursion is \cref{ctr:thm:compiler}; without it the
predecessor is the construction of \cref{seed:sec:succ} and primitive recursion is compiled
as in \cref{base:thm:primrec}. If the independent initializer is removed as well and no
equality test is added, no term denotes the predecessor (\cref{ctr:thm:unary}). The
witness-free introductions that use the computing predecessor of the counter-aware
recursor are treated in \cref{wit:sec:recursor,wit:sec:copy}, on the normal terms of
\cref{ctr:def:normal}.

\section{Translations to iterator terms and the transfer of the bridge}\label{ctr:sec:bridge}

This section draws on \pkg{register\_graph}, \pkg{canonical\_recursors},
\pkg{clausal\_canonical}, \pkg{counter\_evaluation}, \pkg{uniform\_binding},
\pkg{uniform\_recursion} and \pkg{single\_constructor}: each of them proves, for its own
syntax, a bi-interpretation with $\iPRA$. Here the statement is proved once, relative to
$\Teq$, and \cref{base:thm:bridge} is then applied.

\subsection{The ordinary theories}

\begin{definition}[Ordinary theories]\label{ctr:def:ordinary}
The theories $\mathsf T_R^1$ and $\mathsf T_R^2$ have the raw terms of the seeded recursor
with one and with two bound variables, the formulas and logical rules of
\cref{ctr:def:TJ}, and the following arithmetic part: the two computation equations
(\eqref{ctr:eq:r1zero}, \eqref{ctr:eq:r1succ}, respectively \eqref{ctr:eq:r2zero},
\eqref{ctr:eq:r2succ}); coverage $\forall x\,(x=0\lor\exists y\,x=\Suc y)$; separation
$0=\Suc 0\to\bot$; and the rule \eqref{ctr:eq:EI} of induction for literal equations with an
arbitrary context. The theory $\mathsf T_J$ is that of \cref{ctr:def:TJ}; coverage is a
theorem of it (\cref{ctr:lem:coverage}). The theory $\mathsf T_R^1$ is the theory
$\mathsf T_R$ of \cref{seed:def:succ}.
\end{definition}

These are the theory $\Teq$ of \cref{base:sec:teq} with the terms exchanged. The separation
axiom is part of each theory: the positive formulas read in isolation admit a one-element
structure. Since a fixed body is a function symbol of its parameters, ordinary equality
logic gives congruence in the parameters of a body. Replacing a body by a different body
that is provably equal to it pointwise is not an instance of congruence; it is the next
lemma.

\begin{lemma}[Recurrence uniqueness]\label{ctr:lem:unique}
Let $\mathsf T$ be $\Teq$, $\mathsf T_R^1$ or $\mathsf T_R^2$, and let $F(z)$, $G(z)$,
$k(z,w)$, $k'(z,w)$ be terms with $z,w$ fresh for $\Gamma$. If
\[
 \Gamma\proves F(0)=G(0),\quad \Gamma\proves F(\Suc z)=k(z,F(z)),\quad
 \Gamma\proves G(\Suc z)=k'(z,G(z)),\quad \Gamma\proves k(z,w)=k'(z,w),
\]
then $\Gamma\proves F(t)=G(t)$ for every term $t$.
\end{lemma}
\begin{proof}
Apply \eqref{ctr:eq:EI} to the literal equation $F(z)=G(z)$, with all parameters fixed. The
base is given. Under $F(z)=G(z)$,
\[
 F(\Suc z)=k(z,F(z))=k(z,G(z))=k'(z,G(z))=G(\Suc z),
\]
by congruence in the argument $w$ of $k$ and the instance $w\defeq G(z)$ of the fourth
hypothesis.
\end{proof}

\begin{lemma}[Body extensionality]\label{ctr:lem:bodyext}
In $\mathsf T_R^1$ and $\mathsf T_R^2$, if $\Gamma\proves a=b$ and $\Gamma\proves h=k$ with
the bound variables $\vec u$ free in this premise and fresh for $\Gamma$, then
$\Gamma\proves R(n,a;\vec u.h)=R(n,b;\vec u.k)$ for every $n$.
\end{lemma}
\begin{proof}
For one bound variable this is \cref{seed:lem:succ-unique}. In general take
$F(z)=R(z,a;\vec u.h)$ and $G(z)=R(z,b;\vec u.k)$ in \cref{ctr:lem:unique}. At $0$ the
two sides are $\Suc a$ and $\Suc b$, equal by congruence. The successor equations are the
computation equations, and the updates agree by the second premise.
\end{proof}

\begin{lemma}[Normalization is sound]\label{ctr:lem:normsound}
$\mathsf T_R^1$ and $\mathsf T_R^2$ prove $t=\nu(t)$ for every raw term $t$, and
$\mathsf T_R^2$ proves $t=\nf(t)$. Hence every formula is provably equivalent to its
normalization, and a raw substitution instance is provably equivalent to the normalized one.
\end{lemma}
\begin{proof}
For $\nu$ induct on $t$. The counter and the initializer change by congruence, the body by
\cref{ctr:lem:bodyext} and the induction hypothesis with the bound variables free. If the
counter is literally $0$, the zero computation, applied to the term and to its
normalization, identifies both with $\Suc a$.

For $\nf$ prove $\mathsf T_R^2\proves t[\rho]=\operatorname{ev}(t,\rho)$ for every normal
environment $\rho$, by induction on $t$. At a recursor node the induction hypotheses give
$n[\rho]=c$, $a[\rho]=v$ and, at the fresh variables, $h[\rho]=H(i,u)$; hence
$R(n,a;i,u.h)[\rho]=R(c,v;i,u.H(i,u))$ by \cref{ctr:lem:bodyext}. It remains to prove
$R(c,v;i,u.H(i,u))=\operatorname{run}(c,v,H)$ by induction on the successor prefix of $c$.
If $c\equiv 0$ this is the zero computation. If $c$ is neutral the two sides are identical.
If $c\equiv\Suc d$, the successor computation gives
$H(i,u)[d/i,R(d,v;i,u.H(i,u))/u]$; replace the state by $\operatorname{run}(d,v,H)=:w$ using
the hypothesis for $d$, and observe that $H(i,u)[d/i,w/u]=h[\rho][d/i,w/u]=H(d,w)$ by the
induction hypothesis for $h$ in the two environments. The statements on formulas follow by
transport of formulas along equations.
\end{proof}

\subsection{The translations}

\begin{definition}[Translations]\label{ctr:def:translations}
Let $\mathsf T$ be one of $\mathsf T_R^1$, $\mathsf T_R^2$, $\mathsf T_J$. The translation
$t\mapsto t^\iota$ from $\mathsf T$ to $\Teq$ fixes variables and $0$ and is defined on the
binding constructor as follows.
\begin{enumerate}
\item For $\mathsf T_R^1$: $R(n,a;u.h)^\iota=\It{n^\iota}{\Suc a^\iota}{u}{h^\iota}$, as in
\cref{seed:def:succ-trans}.
\item For $\mathsf T_R^2$: $R(n,a;i,u.h)^\iota=F(n^\iota)$, where $F$ is the term given by
the compiler of \cref{base:thm:primrec} for the recursion
$F(0)=\Suc a^\iota$, $F(\Suc z)=h^\iota[z/i,F(z)/u]$.
\item For $\mathsf T_J$: $J(n;i,r.h)^\iota=F(n^\iota)$, where $F$ is given by the same
compiler for the recursion
\[
 F(0)=\Suc 0,\qquad F(\Suc z)=\sel\bigl(e,\ \Suc F(z),\ \mathrm{pd}(e)\bigr),\qquad
 e=h^\iota[z/i,F(z)/r],
\]
and $\mathrm{pd}$ is the predecessor term of \cref{base:thm:pred}, written $P$ there.
\end{enumerate}
The compiler of \cref{base:thm:primrec} provides, for terms $a$ and $g(y,v)$ of $\Teq$, a
defined term with the value $a$ at $0$ and the value $g(c,\cdot)$ of the preceding value at
$\Suc c$; it is built with pairing. That derived recursor of $\Teq$ is not the constructor
$R$ of this chapter, whose value at $0$ is $\Suc a$.
The translation $s\mapsto s^\kappa$ from $\Teq$ to $\mathsf T$ fixes variables and $0$, sends
$\Suc s$ to the defined successor of $s^\kappa$, and sends the iterator to the initialized
iteration of $\mathsf T$:
\[
 \It{n}{a}{u}{h}^\kappa=
 \begin{cases}
  \widehat I(n^\kappa,a^\kappa;u.h^\kappa) & \text{in } \mathsf T_R^1
   \text{ (\cref{seed:thm:succ-iter}; the translation $\rho$ of \cref{seed:def:succ-trans})},\\
  \operatorname{Rec}(n^\kappa,a^\kappa;i,u.h^\kappa),\ i \text{ fresh}, & \text{in } \mathsf T_R^2
   \text{ (\cref{ctr:lem:decoder})},\\
  \operatorname{Rec}_J(n^\kappa,a^\kappa;i,u.h^\kappa),\ i \text{ fresh}, & \text{in } \mathsf T_J
   \text{ (\cref{ctr:thm:regcompiler})}.
 \end{cases}
\]
Both translations commute with substitution and are extended to formulas so that they
commute with every connective and quantifier.
\end{definition}

\begin{theorem}[Transfer of the bridge]\label{ctr:thm:transfer}
Let $\mathsf T$ be $\mathsf T_R^1$, $\mathsf T_R^2$ or $\mathsf T_J$.
\begin{enumerate}
\item If $\mathsf T\proves A$ then $\Teq\proves A^\iota$, and if $\Teq\proves B$ then
$\mathsf T\proves B^\kappa$.
\item $\mathsf T\proves t^{\iota\kappa}=t$ for every term $t$ of $\mathsf T$, and
$\Teq\proves s^{\kappa\iota}=s$ for every term $s$ of $\Teq$.
\item Consequently $\mathsf T\proves A^{\iota\kappa}\leftrightarrow A$ and
$\Teq\proves B^{\kappa\iota}\leftrightarrow B$; both translations are faithful; and
$\mathsf T$ is bi-interpretable with $\Teq$, hence by \cref{base:thm:bridge} with $\Nat$ and
with $\iPRA$, by interpretations whose domains are the whole numerical domain, whose
equality is actual equality and whose numerical composites have the graph $x=y$. The four
theories are equiconsistent.
\end{enumerate}
The same holds for the one-variable register of \cref{ctr:prop:scalarreg}, with
$J_s(n;r.h)^\iota=\It{n^\iota}{\Suc 0}{v}{\sel(h^\iota(v),\Suc v,\mathrm{pd}(h^\iota(v)))}$
and the iteration $\widehat I$ of that proposition for $\kappa$.
\end{theorem}
\begin{proof}
For $\mathsf T_R^1$ the three statements are \cref{seed:thm:succ-terms,seed:thm:succ-bridge}.
The proof is written for $\mathsf T_R^2$ and $\mathsf T_J$.

(1) \emph{From $\mathsf T$ to $\Teq$.} The logical rules translate to themselves. The
translation of the defined successor satisfies
\begin{equation}\label{ctr:eq:succiota}
 \Teq\proves(\Suc t)^\iota=\Suc(t^\iota).
\end{equation}
For the seeded recursor this is the base equation of the translated recursion. For the
register, $(\Suc t)^\iota=F_0(t^\iota)$, where $F_0$ translates the zero body and satisfies
$F_0(0)=\Suc 0$ and $F_0(\Suc z)=\sel(0,\Suc F_0(z),\mathrm{pd}(0))=\Suc F_0(z)$; induction on
the literal equation $F_0(z)=\Suc z$ gives \eqref{ctr:eq:succiota}. Thus the defined
successor is interpreted as the successor of $\Teq$ at every element of the domain, not only
at numerals. With \eqref{ctr:eq:succiota} and transport, the computation equations of the
seeded recursor translate to the defining equations of the translated recursion, coverage
and separation translate to the corresponding axioms of $\Teq$, and an instance of
\eqref{ctr:eq:EI} translates to an instance of the induction rule of $\Teq$ for a literal
equation. For $\mathsf T_J$, \eqref{ctr:eq:J0} is the base equation. For
\eqref{ctr:eq:STEP} write $j=F(n)$ and $e$ for the translated instruction, so that
$F(\Suc n)=\sel(e,\Suc j,\mathrm{pd}(e))$, and split coverage for $e$. If $e=0$ then
$F(\Suc n)=\Suc j$, and both sides of the biconditional are equivalent to $b=\Suc j$, the
second alternative being excluded by separation. If $e=\Suc y$ then $F(\Suc n)=y$, and both
sides are equivalent to $b=y$, by injectivity of the successor. \eqref{ctr:eq:Sep} is
separation.

\emph{From $\Teq$ to $\mathsf T$.} The two computation equations of the iterator translate
to the equations of $\operatorname{Rec}$ and $\operatorname{Rec}_J$, which are
derivable by \cref{ctr:lem:decoder,ctr:thm:regcompiler}. Coverage and separation are axioms
or theorems of $\mathsf T$, and the induction rule for literal equations is
\eqref{ctr:eq:EI}.

(2) Induct externally on the term. Let $t=R(n,a;i,u.h)$ in $\mathsf T_R^2$ and let
$F(z)$ be the $\kappa$-translation of the compiled recursion of
\cref{ctr:def:translations}(2). By part~(1), $\mathsf T_R^2$ proves
$F(0)=\Suc a^{\iota\kappa}$ and $F(\Suc z)=h^{\iota\kappa}[z/i,F(z)/u]$. By the induction
hypotheses, $a^{\iota\kappa}=a$ and $h^{\iota\kappa}=h$ with $i,u$ free. The term
$G(z)=R(z,a;i,u.h)$ has the same base and the same update, so \cref{ctr:lem:unique} gives
$F(z)=G(z)$, and the instance at $n^{\iota\kappa}=n$ is the claim.
In the other direction, for $s=\It{n}{a}{u}{h}$ the term $s^{\kappa\iota}$ is $F(n^{\kappa\iota})$
for a term $F$ of $\Teq$ that satisfies the two equations of the iterator with initial value
$a^{\kappa\iota}$ and update $h^{\kappa\iota}$, by part~(1); \cref{ctr:lem:unique} in $\Teq$
and the induction hypotheses give $s^{\kappa\iota}=s$. For $s=\Suc s_0$ use
\eqref{ctr:eq:succiota}.

For $t=J(n;i,r.h)$ in $\mathsf T_J$ the double translation is $F(n^{\iota\kappa})$ with
$F(0)=1$ and $F(\Suc z)=C'(e',\Suc F(z),P'(e'))$, where $e'=h^{\iota\kappa}[z/i,F(z)/r]$ and
$C',P'$ are the translated selector and predecessor, which have their endpoint equations in
$\mathsf T_J$. Apply \eqref{ctr:eq:EI} to the equation $F(z)=J(z;i,r.h)$. The base is
\eqref{ctr:eq:J0}. Under the induction equation and the induction hypothesis for $h$, the
two instructions $e'$ and $e_z$ are equal. Split coverage for this common value. In the zero
branch both sides increment equal registers, by \eqref{ctr:eq:inc}; in the successor branch
both load the same number, by \eqref{ctr:eq:load}. The invariant is a single equation, and
the case split is supplied by coverage inside the step.

(3) Since the translations commute with the logical operations, part~(2) and transport give
the two formula equivalences by induction on formulas. Faithfulness follows: if
$\Teq\proves A^\iota$ then $\mathsf T\proves A^{\iota\kappa}$ and hence $\mathsf T\proves A$.
Both translations fix the domain, equality and $\bot$, so the identity graph is a provably
total bijection that preserves every operation in both composites. The one-variable register
is treated in the same way: coverage validates its graph axiom for the displayed iterator,
induction on a literal equation shows that its successor abbreviation agrees with the
successor, and recurrence uniqueness gives both composites.
\end{proof}

\begin{remark}[The direct interpretations of the packages]\label{ctr:rem:direct}
Each package states the theorem against $\iPRA$ directly, with the same two maps: a recursor
or register description is interpreted by the primitive recursive function symbol with the
same base and step, and a function symbol of $\iPRA$ is interpreted by the compiler of
\cref{ctr:thm:compiler} or \cref{ctr:thm:regcompiler}, without pairing. By
\cref{ctr:lem:unique} these maps agree provably with the composites of $\iota$ and $\kappa$
with the interpretations of \cref{base:thm:bridge}. \pkg{canonical\_recursors} and
\pkg{clausal\_canonical} use the first clause of \cref{ctr:def:translations} and the
iteration $\widehat I$, and compile primitive recursion with pairing as in
\cref{base:thm:primrec}. The mere agreement of two theories on standard numerals, or mutual
interpretability, would not give \cref{ctr:thm:transfer}; the composites are proved inside
the theories by \cref{ctr:lem:unique}.
\end{remark}

\begin{proposition}[Quantifier-free induction]\label{ctr:prop:qfi}
Each of $\mathsf T_R^1$, $\mathsf T_R^2$ and $\mathsf T_J$ proves the successor to be
injective, proves equality decidable, and admits induction for all quantifier-free formulas
with an arbitrary context.
\end{proposition}
\begin{proof}
Let $A(z)$ be quantifier-free in the language of $\mathsf T$. By \cref{base:thm:qfi},
$\Teq$ proves the induction principle for the quantifier-free formula $A^\iota$ as a
sentence, with its parameters universally quantified. By \cref{ctr:thm:transfer}(1) the
theory $\mathsf T$ proves the $\kappa$-translation of this sentence, in which the successor
is the defined successor, and by \cref{ctr:thm:transfer}(3) the formula $A^{\iota\kappa}$ may
be replaced by $A$. A sentence of this form yields the rule with an arbitrary context.
Injectivity and decidability of equality transfer in the same way from
\cref{base:thm:pred,base:thm:reflect}.
\end{proof}

\begin{remark}[The reflection argument in the packages]\label{ctr:rem:reflection}
\pkg{register\_graph}, \pkg{counter\_evaluation}, \pkg{uniform\_binding},
\pkg{uniform\_recursion} and \pkg{single\_constructor} do not transfer
\cref{base:thm:reflect,base:thm:qfi}; they repeat the proofs with the terms of
\cref{ctr:sec:counter},
where the predecessor and the selector are available before any induction. With addition and
truncated subtraction compiled by \cref{ctr:thm:compiler}, four inductions on literal
equations give $0+y=y$, $\Suc x+y=\Suc(x+y)$, $d(\Suc x,\Suc y)=d(x,y)$ and
$d(\Suc a,a)=1$; a fifth, on the single equation
\[
 C\bigl(\operatorname{sg}(d(x,y)),\,x,\,y+d(x,y)\bigr)=x
\]
with $x$ fixed, in whose step every case split is supplied by coverage, gives
$d(\Suc a,b)=1\leftrightarrow a=b$. The reflection term is
$\epsilon(a,b)=\operatorname{sg}(K(d(\Suc a,b)))$ with $K(n)=C(n,1,p(n))$, and the
characteristic term of a quantifier-free formula is defined with $0$ for truth by
$b_{s=t}=\epsilon(s,t)$, $b_\bot=1$, $b_{A\land B}=C(b_A,b_B,1)$, $b_{A\lor B}=C(b_A,0,b_B)$,
$b_{A\to B}=C(b_A,b_B,0)$. \pkg{register\_graph} uses the control $d(x,y)$ in place of its
sign, the term $e(a,b)=Z(W(d(\Suc a,b)))$ and characteristic terms with $1$ for truth. An
independent audit shipped with \pkg{uniform\_binding} confirms that the only induction
instances in this argument are literal equations. The argument is that of
\cref{base:thm:reflect,base:thm:qfi} and is not repeated.
\end{remark}

\section{A register language without constants}\label{ctr:sec:constfree}

This section draws on \pkg{constantfree}, which builds on \pkg{register\_graph}.

The register language of \cref{ctr:def:languages}(3) has one constant, $0$. It can be
removed. Three occurrences of $J$ over a variable form a term that is provably constant,
and an equation $z=Z(z)$ then defines zero without a closed term. This section states the
theory without constants, proves that it defines zero and the successor, translates it to
and from $\mathsf T_J$, and obtains the bi-interpretation with $\iPRA$ from
\cref{ctr:thm:transfer}. Only what the removal of the constant requires is proved here; the
bootstrap, the compiler and the reflection argument are those of \cref{ctr:sec:register}
and \cref{ctr:rem:reflection}.

\subsection{The language and the theory}

\begin{definition}[Terms without constants]\label{ctr:def:cfterms}
The terms are
\[
 t ::= x\mid J(n;i,r.h),
\]
with $i,r$ bound in $h$ only. There is no numerical constant, no successor, no initializer
and no pairing. With fresh bound variables in every abbreviation put
\begin{align*}
 \mathsf K(t)&\defeq J(t;i,r.i), & \mathsf U(t)&\defeq J(t;i,r.\mathsf K(r)),\\
 \mathsf Z(t)&\defeq J(\mathsf U(t);i,r.r), & \mathsf S(t)&\defeq J(t;i,r.\mathsf Z(r)).
\end{align*}
Fully expanded, $\mathsf Z(t)=J\bigl(J(t;i,r.J(r;j,s.j));k,v.v\bigr)$; thus $\mathsf U$ has
two occurrences of $J$ and $\mathsf Z$ has three. The successor $\mathsf S$ is a fixed term
with the single free variable $t$. For an arbitrary term $a$, the \emph{anchor}, put
\[
 0_a\defeq\mathsf Z(a),\qquad 1_a\defeq\mathsf S(0_a),\qquad 2_a\defeq\mathsf S(1_a).
\]
Here $a$ is not a distinguished element; it is a passive parameter.
\end{definition}

The sans-serif letters distinguish these terms from the successor $\Suc$ of $\mathsf T_J$
and from the term $Z(x)=f(x,1)$ of \cref{ctr:lem:regboot}. \pkg{constantfree} writes
$K,U,Z,S$ for $\mathsf K,\mathsf U,\mathsf Z,\mathsf S$ and $q$ for the term $Z$ of
\cref{ctr:lem:regboot}.

\begin{lemma}[No closed terms]\label{ctr:lem:cffree}
Every term of \cref{ctr:def:cfterms} has a free variable.
\end{lemma}
\begin{proof}
By induction on the term. A variable is free in itself. The counter $n$ of $J(n;i,r.h)$ is
outside the binder, so a free variable of $n$ is free in $J(n;i,r.h)$.
\end{proof}

The schemas of $\mathsf T_J$ are written once more, with the zero and the successor as
parameters, so that $\mathsf T_J$ and the theory without constants are two instances.

\begin{definition}[Anchored register schemas]\label{ctr:def:anchored}
Let $o$ be a term and $\sigma(x)$ a term with a distinguished variable $x$. For a fixed body
$h(i,r,\vec p)$ write $j_n=J(n;i,r.h)$ and $e_n=h(n,j_n,\vec p)$. The \emph{anchored
schemas at $[o,\sigma]$} are the universal closures of
\begin{align}
 & J(o;i,r.h)=\sigma(o), \tag{J0$[o,\sigma]$}\label{ctr:eq:aJ0}\\
 & J(\sigma(n);i,r.h)=b\;\leftrightarrow\;
   \bigl((e_n=o\land b=\sigma(j_n))\lor e_n=\sigma(b)\bigr), \tag{STEP$[o,\sigma]$}\label{ctr:eq:aSTEP}\\
 & o=\sigma(o)\to\bot, \tag{Sep$[o,\sigma]$}\label{ctr:eq:aSep}
\end{align}
with $b$ a variable distinct from the parameters, together with the rule
\begin{equation}\label{ctr:eq:aEI}
 \infrule{\Gamma\proves E(o)\qquad \Gamma,\ E(z)\proves E(\sigma(z))}{\Gamma\proves E(t)}
 \tag{EI$[o,\sigma]$}
\end{equation}
for literal equations $E$, with $z$ fresh for $\Gamma$, $t$, $o$ and the parameters.
\end{definition}

The theory $\mathsf T_J$ of \cref{ctr:def:TJ} is ordinary logic with the anchored schemas at
$[0,\Suc]$, where $\Suc x=J(x;i,r.0)$: the instances are literally \eqref{ctr:eq:J0},
\eqref{ctr:eq:STEP}, \eqref{ctr:eq:Sep} and \eqref{ctr:eq:EI}.

\begin{definition}[The theory $\mathsf T_J^{\varnothing}$]\label{ctr:def:TJ0}
The formulas of $\mathsf T_J^{\varnothing}$ are those of \cref{ctr:def:TJ} over the terms of
\cref{ctr:def:cfterms}, and its logical rules are those of \cref{ctr:def:TJ}. Its arithmetic
part is the anchored schemas at $[\mathsf Z(a),\mathsf S]$ for every term $a$, with bodies
ranging over the terms without constants:
\[
 J(0_a;i,r.h)=1_a,\qquad
 J(\mathsf S n;i,r.h)=b\leftrightarrow\bigl((e_n=0_a\land b=\mathsf S j_n)\lor e_n=\mathsf S b\bigr),
 \qquad 0_a=1_a\to\bot,
\]
and the rule \eqref{ctr:eq:aEI} with base $E(0_a)$, step $E(z)\proves E(\mathsf S z)$ and
$z$ fresh also for the anchor. The anchor is part of each instance and is universally
quantified in the axioms.
\end{definition}

As in every first-order logic with a nonempty domain, a derivation may contain a free
variable that occurs neither in its context nor in its conclusion. When the variables in
scope are recorded in the judgment, as in the checker of \pkg{constantfree}, this is the
explicit rule
\begin{equation}\label{ctr:eq:param}
 \infrule{\Gamma\proves_{V\cup\{a\}}A}{\Gamma\proves_V A}
 \qquad a\notin V,\ a\notin\FV(\Gamma,A).
 \tag{Param}
\end{equation}
It discharges an arbitrary element, not an assumption. By \cref{ctr:lem:cffree} it cannot be
replaced by the instantiation of a closed term.

\begin{proposition}[Param is needed in an empty scope]\label{ctr:prop:cfparam}
In the calculus of \cref{ctr:def:TJ0} with variable scopes, in which every term of a
derivation is in scope and every axiom instance is a universal closure or an instance of one
at terms in scope, the sentence $\exists x\,(x=x)$ is not derivable in the empty scope
without \eqref{ctr:eq:param}. With \eqref{ctr:eq:param} it is derivable.
\end{proposition}
% ed.: the package asserts that a scoped checker without Param is incomplete but gives no proof; the empty-structure argument below is supplied here.
\begin{proof}
Call a judgment $\Gamma\proves_V A$ valid if $V\neq\varnothing$, or if $V=\varnothing$ and
$A$ is true in the structure with the empty domain whenever all formulas of $\Gamma$ are.
Every rule other than \eqref{ctr:eq:param} preserves validity. A rule with a conclusion in a
nonempty scope has a valid conclusion. In the empty scope there is no term, by
\cref{ctr:lem:cffree}, so reflexivity, transport, $\forall$-elimination,
$\exists$-introduction, the instances of axioms at terms, and \eqref{ctr:eq:aEI}, whose
conclusion and anchor are terms, do not occur there. Universal closures are true in the
empty structure, and so is every universal conclusion of $\forall$-introduction. The
propositional rules preserve truth. For $\exists$-elimination in the empty scope, the
major premise $\exists x\,A$ is false in the empty structure, so its validity makes the
context false and the conclusion valid. Since $\exists x\,(x=x)$ is false in the empty
structure, it has no derivation in the empty scope. With \eqref{ctr:eq:param}: in scope
$\{a\}$, reflexivity gives $a=a$ and $\exists$-introduction gives $\exists x\,(x=x)$;
\eqref{ctr:eq:param} discharges $a$.
\end{proof}

\begin{remark}[Bodies as function symbols]\label{ctr:rem:cfdescriptors}
As in \cref{ctr:sec:bridge}, a fixed body is a function symbol $J_h(n,\vec p)$ of its
parameters, and the schemas are supplied for each one. A capture-avoiding substitution
into a body yields another symbol; its agreement with the substituted old one is not an
axiom. It follows by external induction on the body: inner occurrences of $J$ are handled by
the induction hypothesis, and the outer occurrence by \cref{ctr:lem:cfext} below.
\pkg{constantfree} gives this argument; its auditors add that it needs neither injectivity
of the successor nor coverage.
\end{remark}

\subsection{Zero and the successor}

The following two lemmas hold for every instance of \cref{ctr:def:anchored}. They are
applied to $\mathsf T_J^{\varnothing}$ at $[\mathsf Z(a),\mathsf S]$ and to $\mathsf T_J$
at $[0,\Suc]$.

\begin{lemma}[Constant one and zero]\label{ctr:lem:cfconst}
Let a theory contain ordinary logic with equality, prove every instance of
\eqref{ctr:eq:aJ0} and of the right-to-left direction of \eqref{ctr:eq:aSTEP}, and be
closed under \eqref{ctr:eq:aEI}, at one pair $[o,\sigma]$. Then it proves, for every term
$t$,
\[
 \mathsf K(o)=\sigma(o),\quad \mathsf K(\sigma(o))=\sigma(\sigma(o)),\quad
 \mathsf U(t)=\sigma(o),\quad J(\sigma(o);i,r.r)=o,\quad \mathsf Z(t)=o .
\]
Neither \eqref{ctr:eq:aSep}, nor coverage, nor the left-to-right direction of
\eqref{ctr:eq:aSTEP} is used, and the only induction is one instance of
\eqref{ctr:eq:aEI} on the literal equation $\mathsf U(x)=\sigma(o)$.
\end{lemma}
\begin{proof}
\eqref{ctr:eq:aJ0} for the body $i$ gives $\mathsf K(o)=\sigma(o)$. In
\eqref{ctr:eq:aSTEP} for the body $i$ at the counter $n\defeq o$ the instruction is
$e_o=o$, literally, and the old register is $j_o=\mathsf K(o)$. Take
$b\defeq\sigma(\sigma(o))$. The increment disjunct $e_o=o\land b=\sigma(\mathsf K(o))$ holds
by reflexivity and by congruence from $\mathsf K(o)=\sigma(o)$, so the right-to-left
direction gives $\mathsf K(\sigma(o))=\sigma(\sigma(o))$.

\eqref{ctr:eq:aJ0} for the body $\mathsf K(r)$ gives $\mathsf U(o)=\sigma(o)$. Assume
$\mathsf U(z)=\sigma(o)$. In \eqref{ctr:eq:aSTEP} for the body $\mathsf K(r)$ at the counter
$z$ the instruction is $\mathsf K(\mathsf U(z))$, which equals
$\mathsf K(\sigma(o))=\sigma(\sigma(o))$ by congruence and the preceding paragraph. With
$b\defeq\sigma(o)$ the load disjunct $e_z=\sigma(b)$ holds, and the right-to-left direction
gives $\mathsf U(\sigma(z))=\sigma(o)$. Now \eqref{ctr:eq:aEI}, applied to the literal
equation $\mathsf U(x)=\sigma(o)$ in the empty context, gives $\mathsf U(t)=\sigma(o)$.

Let $B(n)\defeq J(n;i,r.r)$. \eqref{ctr:eq:aJ0} gives $B(o)=\sigma(o)$. In
\eqref{ctr:eq:aSTEP} for the body $r$ at the counter $o$ the instruction is the old register
$B(o)$; with $b\defeq o$ the load disjunct $B(o)=\sigma(o)$ holds, so $B(\sigma(o))=o$.
Finally $\mathsf Z(t)=B(\mathsf U(t))$ literally, and transport along
$\mathsf U(t)=\sigma(o)$ gives $\mathsf Z(t)=B(\sigma(o))=o$.
\end{proof}

\begin{lemma}[Register body extensionality]\label{ctr:lem:cfext}
Let a theory contain ordinary logic with equality, prove every instance of
\eqref{ctr:eq:aJ0} and of both directions of \eqref{ctr:eq:aSTEP}, and be closed under
\eqref{ctr:eq:aEI}, at one pair $[o,\sigma]$. Let $h(i,r)$ and $k(i,r)$ be bodies, with
arbitrary parameters, and let $i,r$ be fresh for $\Gamma$. If $\Gamma\proves h=k$, then
$\Gamma\proves J(n;i,r.h)=J(n;i,r.k)$ for every term $n$. Coverage, separation and
injectivity of the successor are not used.
\end{lemma}
\begin{proof}
Apply \eqref{ctr:eq:aEI} to the literal equation $E(x)\colon J(x;i,r.h)=J(x;i,r.k)$ in the
context $\Gamma$, with $z$ fresh. The base holds since both sides equal $\sigma(o)$ by
\eqref{ctr:eq:aJ0}. For the step write $j=J(z;i,r.h)$, $j'=J(z;i,r.k)$, $e=h[z/i,j/r]$ and
$e'=k[z/i,j'/r]$, and assume $j=j'$. Since $i,r$ are not free in $\Gamma$, the premise gives
$h[z/i,j/r]=k[z/i,j/r]$, and congruence along $j=j'$ gives $e=e'$. Put
$c\defeq J(\sigma(z);i,r.h)$. The left-to-right direction of \eqref{ctr:eq:aSTEP} for $h$
at the counter $z$, applied to the reflexive equation $c=c$, gives
\[
 (e=o\land c=\sigma(j))\lor e=\sigma(c).
\]
Transport along $e=e'$ and $j=j'$ gives $(e'=o\land c=\sigma(j'))\lor e'=\sigma(c)$, and the
right-to-left direction of \eqref{ctr:eq:aSTEP} for $k$ at the counter $z$ gives
$J(\sigma(z);i,r.k)=c$, which is $E(\sigma(z))$ after symmetry.
\end{proof}

In $\mathsf T_J^{\varnothing}$ the lemma is applied with an anchor $a$ that is any term in
scope, for instance the counter $n$ itself; by \cref{ctr:lem:cffree} the conclusion always
contains one. For $\mathsf T_J$, \cref{ctr:lem:cfext} is a proof of the identification of
two register computations without the coverage split used in part~(2) of the proof of
\cref{ctr:thm:transfer}.

\begin{corollary}[Zero without a zero term]\label{ctr:cor:cfzero}
Write $\operatorname{Zero}(z)\defeq(z=\mathsf Z(z))$. The theory $\mathsf T_J^{\varnothing}$
proves, for all terms $a,b,t$,
\[
 \mathsf U(t)=\mathsf S(\mathsf Z(a)),\qquad \mathsf Z(t)=\mathsf Z(a),\qquad
 \mathsf S(t)=J(t;i,r.\mathsf Z(a)),
\]
and the sentences
\[
 \forall a\,\forall b\ \mathsf Z(a)=\mathsf Z(b),\qquad
 \exists z\,\operatorname{Zero}(z),\qquad
 \forall x\,\forall y\,\bigl(\operatorname{Zero}(x)\land\operatorname{Zero}(y)\to x=y\bigr).
\]
\end{corollary}
\begin{proof}
The first two equations are \cref{ctr:lem:cfconst} at $[\mathsf Z(a),\mathsf S]$. The third
is \cref{ctr:lem:cfext} for the bodies $\mathsf Z(r)$ and $\mathsf Z(a)$, whose premise
$\mathsf Z(r)=\mathsf Z(a)$ is the second equation with $r$ free; the anchor $a$ may be
taken for the lemma. The first sentence follows by $\forall$-introduction. For existence,
the second equation at $t\defeq\mathsf Z(a)$ gives $\mathsf Z(\mathsf Z(a))=\mathsf Z(a)$,
so $\mathsf Z(a)$ is a witness; the free variable $a$ is discharged as in
\eqref{ctr:eq:param}. For uniqueness, $x=\mathsf Z(x)=\mathsf Z(y)=y$.
\end{proof}

\begin{corollary}[The terms in $\mathsf T_J$]\label{ctr:cor:cfTJ}
$\mathsf T_J$ proves $\mathsf U(t)=1$, $\mathsf Z(t)=0$, $\mathsf S(t)=\Suc t$ and
$\operatorname{Zero}(z)\leftrightarrow z=0$ for all terms $t$ and variables $z$.
\end{corollary}
\begin{proof}
\Cref{ctr:lem:cfconst} at $[0,\Suc]$ gives the first two. \Cref{ctr:lem:cfext} for the
bodies $\mathsf Z(r)$ and $0$, with the premise $\mathsf Z(r)=0$, gives
$\mathsf S(t)=J(t;i,r.0)=\Suc t$. Finally $\mathsf Z(z)=0$ gives the last equivalence.
\end{proof}

\begin{remark}[The bootstrap at an anchor]
\pkg{constantfree} also carries out, at a fixed anchor, the bootstrap of
\cref{ctr:lem:regboot} with $f(n,0)=\mathsf S n$ obtained from the third equation of
\cref{ctr:cor:cfzero}, coverage as in \cref{ctr:lem:coverage}, the compiler of
\cref{ctr:thm:regcompiler} in the variant with the instruction
$\mathsf S(C(i,b,h(p(i),r)))$, and the reflection argument of \cref{ctr:rem:reflection} with
characteristic terms that use $1$ for truth. These are the constructions of
\cref{ctr:sec:register} with $\mathsf Z(a)$ for $0$ and $\mathsf S$ for $\Suc$. By
\cref{ctr:thm:cfbridge} below they are not needed separately: they hold in
$\mathsf T_J^{\varnothing}$ as translations of theorems of $\mathsf T_J$.
\end{remark}

\subsection{Translations and the bi-interpretation}

\begin{definition}[Translations]\label{ctr:def:cftrans}
The translation from $\mathsf T_J^{\varnothing}$ to $\mathsf T_J$ is the inclusion: every
term and formula without constants is a term or formula of $\mathsf T_J$. For a term or
formula $X$ of $\mathsf T_J$ and a variable $z$ that does not occur in $X$, let $X_z$ be
$X$ with every occurrence of the constant $0$ replaced by $z$, including the occurrences
inside the abbreviations $\Suc$ and $1$; thus $(\Suc t)_z=J(t_z;i,r.z)$. The translation
from $\mathsf T_J$ to $\mathsf T_J^{\varnothing}$ is
\[
 A^\flat\defeq\exists z\,\bigl(\operatorname{Zero}(z)\land A_z\bigr),
\]
and it interprets a term $t$ of $\mathsf T_J$ by the graph
$\exists z\,(\operatorname{Zero}(z)\land y=t_z)$, with $z$ chosen fresh. Both translations
fix the variables, the domain and equality, and $\bot^\flat$ is provably equivalent to $\bot$
(\cref{ctr:lem:cfflat}(3) below).
\end{definition}

The replacement commutes with substitution, $(X[s/x])_z=X_z[s_z/x]$, and with every
connective and quantifier whose bound variable is not $z$.

\begin{lemma}[Inclusion]\label{ctr:lem:cfdown}
If $\Gamma\proves A$ in $\mathsf T_J^{\varnothing}$, then $\Gamma\proves A$ in
$\mathsf T_J$.
\end{lemma}
\begin{proof}
By induction on the derivation. The logical rules are shared. A use of \eqref{ctr:eq:param}
is admissible in $\mathsf T_J$: substitute $0$ for the discharged variable throughout the
derivation of its premise, renaming eigenvariables. By \cref{ctr:cor:cfTJ}, $\mathsf T_J$
proves $\mathsf Z(a)=0$ and $\mathsf S(x)=\Suc x$ for all terms $a,x$. Transport along these
equations, at each occurrence, turns every instance of an anchored schema at
$[\mathsf Z(a),\mathsf S]$ into the instance with the same body at $[0,\Suc]$, which is an
axiom of $\mathsf T_J$. For an instance of \eqref{ctr:eq:aEI} at $[\mathsf Z(a),\mathsf S]$
the induction hypotheses give $\Gamma\proves E(\mathsf Z(a))$ and
$\Gamma,E(z)\proves E(\mathsf S z)$ in $\mathsf T_J$; transport gives
$\Gamma\proves E(0)$ and $\Gamma,E(z)\proves E(\Suc z)$, and \eqref{ctr:eq:EI} gives
$\Gamma\proves E(t)$.
\end{proof}

\begin{lemma}[Zero as a parameter]\label{ctr:lem:cfup}
Let $\Gamma\proves A$ in $\mathsf T_J$ and let $z$ be a variable that occurs nowhere in the
derivation. Then $\Gamma_z,\operatorname{Zero}(z)\proves A_z$ in
$\mathsf T_J^{\varnothing}$.
\end{lemma}
\begin{proof}
Under the hypothesis $\operatorname{Zero}(z)$ the theory $\mathsf T_J^{\varnothing}$ proves
\begin{equation}\label{ctr:eq:cfpar}
 \mathsf Z(a)=z,\qquad J(t;i,r.z)=\mathsf S(t)
\end{equation}
for all terms $a,t$: the first is $\mathsf Z(a)=\mathsf Z(z)=z$ by
\cref{ctr:cor:cfzero}, and the second is \cref{ctr:lem:cfext} for the bodies $z$ and
$\mathsf Z(r)$, with the premise $z=\mathsf Z(r)$ given by the first. Now induct on the
derivation. A logical rule maps to the same rule with $\operatorname{Zero}(z)$ added to the
context; eigenvariables differ from $z$ and do not occur in $\Gamma$, hence not in
$\Gamma_z,\operatorname{Zero}(z)$; transport and $\forall$-elimination commute with the
replacement.

Axioms. $(J(0;i,r.h)=\Suc 0)_z$ is $J(z;i,r.h_z)=J(z;i,r.z)$. The instance of
\eqref{ctr:eq:aJ0} at the anchor $z$ is $J(\mathsf Z(z);i,r.h_z)=\mathsf S(\mathsf Z(z))$, and
transport along $\mathsf Z(z)=z$ and $\mathsf S(z)=J(z;i,r.z)$ from \eqref{ctr:eq:cfpar}
gives the translated axiom. The translation of \eqref{ctr:eq:STEP} is
\[
 J\bigl(J(n_z;i,r.z);i,r.h_z\bigr)=b\leftrightarrow
 \bigl((e'=z\land b=J(j';i,r.z))\lor e'=J(b;i,r.z)\bigr),
\]
with $j'=J(n_z;i,r.h_z)$ and $e'=h_z[n_z/i,j'/r]$, since the replacement commutes with
substitution. The instance of \eqref{ctr:eq:aSTEP} at the anchor $z$ for the body $h_z$ and
the counter $n_z$ differs from it only by $\mathsf Z(z)$ in place of $z$ and $\mathsf S(x)$
in place of $J(x;i,r.z)$ for $x=n_z,j',b$, so transport along \eqref{ctr:eq:cfpar} proves
it. The translation of \eqref{ctr:eq:Sep} is $z=J(z;i,r.z)\to\bot$, and transport turns the
instance $\mathsf Z(z)=\mathsf S(\mathsf Z(z))\to\bot$ of \eqref{ctr:eq:aSep} into it.

Induction. For an instance of \eqref{ctr:eq:EI} with premises $\Gamma\proves E(0)$ and
$\Gamma,E(y)\proves E(\Suc y)$, the induction hypotheses give
$\Gamma_z,\operatorname{Zero}(z)\proves E_z(z)$ and
$\Gamma_z,\operatorname{Zero}(z),E_z(y)\proves E_z(J(y;i,r.z))$. Transport along
\eqref{ctr:eq:cfpar} gives $E_z(\mathsf Z(z))$ and $E_z(\mathsf S y)$, and
\eqref{ctr:eq:aEI} at the anchor $z$ gives $E_z(t_z)$. The invariant $E_z$ is a literal
equation, and $y$ differs from $z$ and is fresh for the extended context and the target.
\end{proof}

\begin{lemma}[The translation $\flat$]\label{ctr:lem:cfflat}
In $\mathsf T_J^{\varnothing}$:
\begin{enumerate}
\item $\operatorname{Zero}(z)\proves A^\flat\leftrightarrow A_z$ for every formula $A$ of
$\mathsf T_J$ and every variable $z$ not in $A$;
\item $\flat$ commutes with every connective and quantifier up to provable equivalence:
$(A\circ B)^\flat\leftrightarrow A^\flat\circ B^\flat$ for $\circ\in\{\land,\lor,\to\}$,
$\bot^\flat\leftrightarrow\bot$, $(\forall x\,A)^\flat\leftrightarrow\forall x\,A^\flat$ and
$(\exists x\,A)^\flat\leftrightarrow\exists x\,A^\flat$;
\item $A^\flat\leftrightarrow A$ for every formula $A$ without constants.
\end{enumerate}
\end{lemma}
\begin{proof}
(1) If $\operatorname{Zero}(w)$ and $A_w$, uniqueness in \cref{ctr:cor:cfzero} gives $w=z$,
and transport gives $A_z$. Conversely $z$ is a witness. (2) Let $z$ be a variable free in
neither side. By existence in \cref{ctr:cor:cfzero} and $\exists$-elimination it
suffices to prove the equivalence under $\operatorname{Zero}(z)$. There,
by (1), applied under the quantifier for $\forall x$ and $\exists x$ since $x$ is not $z$,
both sides are equivalent to $(A\circ B)_z=A_z\circ B_z$, respectively to
$\forall x\,A_z$ and $\exists x\,A_z$. (3) Here $A_z$ is $A$, and
$\exists z\,(\operatorname{Zero}(z)\land A)\leftrightarrow A$ by existence.
\end{proof}

\begin{theorem}[Constant-free presentation]\label{ctr:thm:cfbridge}
\leavevmode
\begin{enumerate}
\item If $\mathsf T_J^{\varnothing}\proves A$ then $\mathsf T_J\proves A$. If
$\Gamma\proves A$ in $\mathsf T_J$ then $\Gamma^\flat\proves A^\flat$ in
$\mathsf T_J^{\varnothing}$.
\item $\mathsf T_J^{\varnothing}\proves A^\flat\leftrightarrow A$ for every formula $A$ of
$\mathsf T_J^{\varnothing}$, and $\mathsf T_J\proves A^\flat\leftrightarrow A$ for every
formula $A$ of $\mathsf T_J$. The graph of each numerical composite is $x=y$.
\item With $\iota_\varnothing(A)\defeq A^\iota$ and
$\kappa_\varnothing(B)\defeq(B^\kappa)^\flat$, where $\iota,\kappa$ are the translations of
\cref{ctr:def:translations} for $\mathsf T_J$, the theory $\mathsf T_J^{\varnothing}$ is
bi-interpretable with $\Teq$, and hence by \cref{base:thm:bridge} with $\Nat$ and with
$\iPRA$, by interpretations whose domains are the whole numerical domain, whose equality is
actual equality and whose numerical composites have the graph $x=y$. All these theories are
equiconsistent.
\end{enumerate}
\end{theorem}
\begin{proof}
(1) The first statement is \cref{ctr:lem:cfdown}. For the second take $z$ fresh. By
\cref{ctr:lem:cfup}, $\Gamma_z,\operatorname{Zero}(z)\proves A_z$. Under
$\operatorname{Zero}(z)$, \cref{ctr:lem:cfflat}(1) turns every formula of $\Gamma^\flat$ into
the corresponding one of $\Gamma_z$ and $A_z$ into $A^\flat$. Hence
$\Gamma^\flat,\operatorname{Zero}(z)\proves A^\flat$, and $\exists$-elimination on
$\exists z\,\operatorname{Zero}(z)$ removes the hypothesis, since $z$ is free in neither
$\Gamma^\flat$ nor $A^\flat$.

(2) The first composite is \cref{ctr:lem:cfflat}(3). For the second, $\mathsf T_J$ proves
$\operatorname{Zero}(z)\leftrightarrow z=0$ by \cref{ctr:cor:cfTJ}, so
$A^\flat\leftrightarrow\exists z\,(z=0\land A_z)\leftrightarrow A$ by transport, since
$A_z[0/z]$ is $A$. For a term $t$ of $\mathsf T_J$ the composite graph
$\exists z\,(\operatorname{Zero}(z)\land y=t_z)$ is equivalent in $\mathsf T_J$ to $y=t$ in
the same way, and a term without constants is mapped to itself by the inclusion and has the
graph $\exists z\,(\operatorname{Zero}(z)\land y=t)$, equivalent to $y=t$ by
\cref{ctr:lem:cfflat}(3).

(3) By (1) and \cref{ctr:thm:transfer}(1), $\mathsf T_J^{\varnothing}\proves A$ implies
$\Teq\proves A^\iota$, and $\Teq\proves B$ implies $\mathsf T_J\proves B^\kappa$ and
$\mathsf T_J^{\varnothing}\proves(B^\kappa)^\flat$. For a formula $A$ of
$\mathsf T_J^{\varnothing}$, \cref{ctr:thm:transfer}(3) gives
$\mathsf T_J\proves A^{\iota\kappa}\leftrightarrow A$; by (1) and \cref{ctr:lem:cfflat}(2)
the theory $\mathsf T_J^{\varnothing}$ proves $(A^{\iota\kappa})^\flat\leftrightarrow A^\flat$,
and $A^\flat\leftrightarrow A$ by (2). For a formula $B$ of $\Teq$, (2) gives
$\mathsf T_J\proves(B^\kappa)^\flat\leftrightarrow B^\kappa$; by
\cref{ctr:thm:transfer}(1), since $\iota$ commutes with the connectives,
$\Teq\proves((B^\kappa)^\flat)^\iota\leftrightarrow B^{\kappa\iota}$, and
$B^{\kappa\iota}\leftrightarrow B$ by \cref{ctr:thm:transfer}(3). The four translations fix
the domain, equality and the variables, so the numerical composites have the graph $x=y$,
by (2) and \cref{ctr:thm:transfer}(2). Both translations preserve derivability and map
$\bot$ to $\bot$ up to provable equivalence ($\bot^\iota=\bot$ and
$\bot^\flat\leftrightarrow\bot$), so a derivation of $\bot$ in one theory yields one in the
other.
\end{proof}

\begin{remark}[The direct proof of the package]
\pkg{constantfree} proves \cref{ctr:thm:cfbridge}(3) against $\iPRA$ directly. A body is
interpreted by the primitive recursive function $j_h$ with $j_h(0,\vec p)=1$ and
$j_h(n+1,\vec p)=C(e,\,j_h(n,\vec p)+1,\,p(e))$, $e=h(n,j_h(n,\vec p),\vec p)$; in $\iPRA$
the defining equations give $\mathsf K(1)=2$ and $J(1;i,r.r)=0$, one induction on an
equation gives $\mathsf U(t)=1$ and $\mathsf Z(t)=0$, and a second gives $\mathsf S(n)=n+1$,
before the anchored schemas are verified. In the other direction it adds a constant $c$
with $c=\mathsf Z(c)$, which is a conservative definitional extension by
\cref{ctr:cor:cfzero}, compiles primitive recursion at the anchor $c$, and removes $c$ by
the translation $\flat$. \Cref{ctr:lem:cfup} is the same argument with $c$ kept as a
variable $z$. The two routes give the same maps up to provable equality, by
\cref{ctr:lem:unique} in $\Teq$ and \cref{ctr:lem:cfext} in $\mathsf T_J^{\varnothing}$. The
package presents the bi-interpretation as a written proof, with an independent written
audit; it is not formalized.
\end{remark}

\subsection{Sizes of the constant terms}

Let the terms of \cref{ctr:def:cfterms} denote functions of their free variables on the
natural numbers: $J(n;i,r.h)$ is $r_n$, where $r_0=1$, and $r_{m+1}=r_m+1$ if
$h(m,r_m)=0$ and $r_{m+1}=h(m,r_m)-1$ otherwise. This is the standard model of
$\mathsf T_J$, so by \cref{ctr:lem:cfdown} it satisfies $\mathsf T_J^{\varnothing}$.

\begin{proposition}[Sharp sizes]\label{ctr:prop:cfsize}
Count occurrences of $J$ in fully expanded terms. A term that denotes the constant function
$1$ has at least two occurrences, and a term that denotes the constant function $0$ has at
least three. The terms $\mathsf U(x)$ and $\mathsf Z(x)$ attain these bounds.
\end{proposition}
\begin{proof}
A term without $J$ is a variable and is not constant. A term with exactly one $J$ is
$J(x;i,r.v)$ with $x$ a free variable and $v$ a variable. It takes the value $1$ at $x=0$.
It also takes the value $0$: if $v$ is $r$, at $x=1$ the instruction $r_0=1$ loads $0$; if
$v$ is $i$, at $x=2$ the instructions $0$ and $1$ increment to $2$ and then load $0$; if
$v$ is free, possibly $x$ itself, at $x=v=1$ the instruction $1$ loads $0$. Hence no term
with at most one $J$ is constant, and every such term takes the value $0$ at some valuation.

A term with exactly two occurrences is $J(n;i,r.h)$ with $n$ a variable or a term with one
occurrence. The counter $n$ is outside the binder, so its free variables are free in the
whole term, and by the preceding paragraph some valuation makes $n=0$. There the term takes
the value $1$, whatever $h$ is. A term with at most two occurrences therefore takes the
value $1$ somewhere and is not constantly $0$.

$\mathsf U(x)$ has two occurrences and $\mathsf Z(x)$ three, and by \cref{ctr:cor:cfTJ} they
denote $1$ and $0$.
\end{proof}

The bound concerns raw term size only. It is not a bound on the size of an axiomatization or
of a checker.

\subsection{The reference checker}

The checker of \pkg{constantfree} is a standalone program for $\mathsf T_J^{\varnothing}$ in
ordinary logic. Its numerical term tags are a variable and $J$ with a counter, two bound
names and a body; zero and the successor are expanded raw terms. The input is a JSON object
with a closed goal and a proof tree in which every node states its conclusion. Every
arithmetic instance names its anchor, which is checked to be in scope; the anchor of an
induction may not depend on its eigenvariable, and the invariant must be a literal
equation. The biconditional \eqref{ctr:eq:aSTEP} is expanded into a conjunction of
implications, so increments and loads are checked by logical rules and not by evaluation.
\eqref{ctr:eq:param} is a primitive instruction. Four certificates are supplied:
$\forall a\,\forall t\ \mathsf U(t)=\mathsf S(\mathsf Z(a))$,
$\forall a\,\forall t\ \mathsf Z(t)=\mathsf Z(a)$, $\exists z\,(z=\mathsf Z(z))$, which uses
\eqref{ctr:eq:param}, and
$\forall x\,\forall y\,((x=\mathsf Z(x)\land y=\mathsf Z(y))\to x=y)$. The validation runs
$19$ certificate cases in ordinary and optimized Python, $38$ executions, including the
purely logical theorem $\exists x\,(x=x)$ and fourteen rejected variants (an escaping
anchor, an anchor depending on the induction variable, a zero constructor, induction on a
conjunction, malformed input), and $6$ checks of capture-avoiding substitution. These are
finite evidence for the implementation. The compact refinements shipped in the same package
concern the scalar kernel and are treated with the certificate checkers.

\section{Guarded scalar registers}\label{ctr:sec:guarded}

This section draws on \pkg{guarded\_register}.

The register languages of \cref{ctr:sec:register,ctr:sec:constfree} state a computation
step by the graph biconditional \eqref{ctr:eq:STEP}, whose right-to-left direction is the
pair of conditional rules \eqref{ctr:eq:inc} and \eqref{ctr:eq:load}. A restriction on the
bodies of the register with one bound variable makes both computations unconditional
equations: a body either does not depend on the state, or is literally a successor. The
resulting native calculus $\mathcal K_{\mathrm G}$ has the formulas and the six logical rules
of $\Nat$ and two computation equations in place of \rulename{I0} and \rulename{IS}. It has
no successor constructor, no initializer, no reflexivity rule and no conditional increment
or load. \Cref{ctr:thm:gequiv} proves it equivalent to $\Nat$, with both composites.

\subsection{Terms and the guard}

\begin{definition}[Guarded terms]\label{ctr:def:gterms}
The raw terms are
\[
 t ::= x\mid 0\mid J(n;u.h),
\]
with $u$ bound in $h$ only; the counter $n$ is outside the binder. Put
\[
 \Suc a\defeq J(a;v.0),\qquad 1\defeq\Suc 0,\qquad 2\defeq\Suc 1,\qquad
 \Suc^{k+1}a\defeq\Suc(\Suc^k a),
\]
with $v$ unused. A raw term is \emph{guarded} if at every occurrence $J(n;u.h)$ the terms
$n$ and $h$ are guarded and at least one of the following holds:
\begin{enumerate}
\item $u\notin\FV(h)$: the body is \emph{state-free}; passive parameters are allowed;
\item $h\equiv\Suc(g)$, that is $h\equiv J(g;v.0)$, for some term $g$, which may contain $u$:
the body is \emph{successor-headed}.
\end{enumerate}
The two conditions may overlap. They are tests on literal syntax, not on provable equality
or on values. Below every term is guarded, and a \emph{template} $h(v)$ is a guarded term
with a distinguished variable $v$.
\end{definition}

As with $R$ in \cref{ctr:def:languages}, the number of bound variables distinguishes this
$J$ from the register $J(n;i,r.h)$ of \cref{ctr:sec:register}; the one-variable register of
\cref{ctr:prop:scalarreg} is written $J_s$. On the natural numbers $J(n;u.h)$ denotes $r_n$,
where
\[
 r_0=h(0)+1,\qquad r_{m+1}=\begin{cases} r_m+1, & h(r_m)=0,\\ h(r_m)-1, & h(r_m)>0.\end{cases}
\]
The transition is that of $J_s$: the instruction $0$ increments the register and an
instruction $b+1$ loads $b$. The start differs. The register $J_s$ starts at $1$, while
$J(0;u.h)$ is $\Suc(h(0))$, which is the base equation of the self-seeded constructor $F$
of \cref{ctr:sec:bounds}. That constructor has no term for the predecessor
(\cref{ctr:thm:unary}); the guarded register has one (\cref{ctr:lem:gboot}), because its step
at a successor-headed body removes the successor head. The guard excludes the bodies of
$J_s$ that depend on the state and are not literally successors, for instance the
identity body $u$; the definition of zero in \cref{ctr:def:cfterms} uses bodies of that
kind, and the guarded language keeps the constant $0$. At a successor-headed body only the
load branch of the transition can occur, and the equation \eqref{ctr:eq:gR} below states
it. At a state-free body the instruction $a$ does not change along the computation, and the
step at an arbitrary $a$ follows from the two cases $a=0$ and $a=\Suc z$
(\cref{ctr:lem:grec}). Neither case needs the case distinction on a state-dependent
instruction that \eqref{ctr:eq:STEP} provides.

\begin{lemma}[Substitution closure]\label{ctr:lem:gsubst}
If $t$ and $s$ are guarded, then the capture-avoiding substitution $t\subst{s}{x}$ is
guarded. In particular $h\subst{s}{u}$ is guarded for a guarded body $h$ and a guarded term
$s$, and renaming of bound variables preserves the guard.
\end{lemma}
\begin{proof}
Induct on $t$, with the bound variable of each occurrence of $J$ renamed away from $x$ and
from $\FV(s)$. Variables and $0$ are immediate. For $t=J(n;u.h)$,
$t\subst{s}{x}=J(n\subst{s}{x};u.h\subst{s}{x})$, and the counter and the body are guarded
by the induction hypothesis. If $u\notin\FV(h)$, then $u\notin\FV(h\subst{s}{x})$, since $u$
is not free in $s$. If $h\equiv J(g;v.0)$, then $h\subst{s}{x}\equiv J(g\subst{s}{x};v.0)$,
since the zero body has no free variable. Renaming changes neither the set of free variables
of a body other than its bound variable nor the shape $J(g;v.0)$.
\end{proof}

\begin{definition}[The calculus $\mathcal K_{\mathrm G}$]\label{ctr:def:gcalc}
The formulas are the native formulas of \cref{base:sec:kernel} over guarded terms, with the
reading of \cref{base:sec:reading}, and contexts are finite lists of formulas. The rules are
\rulename{Hyp}, \rulename{Tr}, \rulename{QE}, \rulename{DE}, \rulename{DI} and \rulename{QI}
of \cref{base:def:kernel}, with the same side conditions, where the template of
\rulename{Tr} is a guarded template and $\Suc z$ in \rulename{DI} and \rulename{QI} is the
defined successor $J(z;v.0)$, together with the two computation equations, in every
context,
\begin{align}
 & J(0;u.h)=\Suc\bigl(h\subst{0}{u}\bigr), \tag{GB}\label{ctr:eq:gB}\\
 & J(\Suc n;u.\Suc g)=g\subst{J(n;u.\Suc g)}{u}, \tag{GR}\label{ctr:eq:gR}
\end{align}
for every guarded $J(0;u.h)$ and all guarded $n$ and $g$. The counter of \eqref{ctr:eq:gR}
is literally $J(n;v.0)$, and its body $\Suc g$ is admitted whatever $g$ is. The designated
contradiction is $0=1$.
\end{definition}

Compared with $\Nat$, the terms $0,\Suc,I$ are replaced by $0,J$, and the computations
\rulename{I0}, \rulename{IS} by \eqref{ctr:eq:gB}, \eqref{ctr:eq:gR}; the other six rules are
literally those of \cref{base:def:kernel}. As in $\Nat$ (\cref{base:rem:inventory}) there is
no reflexivity rule. Compared with $\mathsf T_J$, the conditional rules \eqref{ctr:eq:inc}
and \eqref{ctr:eq:load} are not primitive (\cref{ctr:prop:grules}). The guard, the test
that the state does not occur and the test of the successor head belong to the trusted
specification; by \cref{ctr:lem:gsubst} they are stable under substitution.

\begin{lemma}[Structural rules and equality]\label{ctr:lem:gstruct}
In $\mathcal K_{\mathrm G}$:
\begin{enumerate}
\item weakening, substitution and cut are admissible, as in \cref{base:lem:struct,base:lem:cut};
\item $\Gamma\proves J(1;u.\Suc t)=t$ for $u\notin\FV(t)$, and the equality rules of
\cref{base:lem:eq} hold: reflexivity, symmetry, transitivity, congruence in every guarded
template, transport in equation templates, and simultaneous congruence;
\item equation induction (\cref{base:lem:eqind}), numerical cases (\cref{base:cor:cases}),
congruence of $Q$ and $D$ (\cref{base:lem:bindercong}) and the Leibniz rule
(\cref{base:thm:leibniz}) hold, with $\Suc z$ the defined successor.
\end{enumerate}
\end{lemma}
\begin{proof}
The proofs in \cref{base:sec:structural} use the computation rules only in the instance
cases of \cref{base:lem:struct,base:lem:cut} and in \cref{base:lem:eq}(i). (1) For a
substitution $\theta$ with $u$ fresh for $\theta$,
\[
 J(0;u.h)[\theta]=J(0;u.h[\theta]),\qquad \Suc\bigl(h\subst{0}{u}\bigr)[\theta]=
 \Suc\bigl(h[\theta]\subst{0}{u}\bigr),
\]
and $J(\Suc n;u.\Suc g)[\theta]=J(\Suc(n[\theta]);u.\Suc(g[\theta]))$,
$\bigl(g\subst{J(n;u.\Suc g)}{u}\bigr)[\theta]
=g[\theta]\subst{J(n[\theta];u.\Suc(g[\theta]))}{u}$. The substituted terms are guarded by
\cref{ctr:lem:gsubst}, so an instance of \eqref{ctr:eq:gB} or \eqref{ctr:eq:gR} becomes an
instance of the same equation, and the remaining cases are those of the two base lemmas.
(2) \eqref{ctr:eq:gR} at the counter $0$, with the constant $g\defeq t$ and $u\notin\FV(t)$,
is $J(1;u.\Suc t)=t$, since $1\equiv\Suc 0$. \rulename{Tr} with the template $v$, applied to
this equation as both premises, gives $t=t$. This replaces the instances of \rulename{I0}
in \cref{base:lem:eq}(i); parts (ii)--(v) and \cref{base:lem:cong} are proved by
\rulename{Tr} alone, with templates that are guarded by \cref{ctr:lem:gsubst}. (3) The proofs
of these statements use \rulename{Hyp}, \rulename{Tr}, \rulename{QE}, \rulename{DE},
\rulename{DI}, \rulename{QI}, parts (1) and (2), and the term $\Suc z$ only as it occurs in
the conclusions of \rulename{DI} and \rulename{QI}. They apply verbatim.
\end{proof}

\subsection{Arithmetic from the two equations}

With fresh bound variables, and $u\notin\FV(a)$ in the first definition, put
\begin{align*}
 f(n,a)&\defeq J(n;u.a), & W(x)&\defeq f(1,x), & Y(n)&\defeq f(n,1),\\
 Z(c)&\defeq W^2\bigl(Y(Y(c))\bigr), & G(c,a)&\defeq W^3\bigl(f(\Suc^2a,Z(c))\bigr), &
 C(c,a,b)&\defeq W^3\bigl(f(\Suc^2a,G(c,\Suc^4b))\bigr),
\end{align*}
and $P(c)\defeq C(c,0,W(c))$, where the powers are finite compositions. In this section the
letters $f,W,Z,G,C,P$ denote these terms; the terms of the same names in
\cref{ctr:lem:regboot} are their analogues for the register with two bound variables.
\pkg{guarded\_register} writes $Z_2$ for $Y$. Every occurrence of $J$ in these terms has a
state-free body, so they are guarded for all guarded arguments, and a template such as
$f(\Suc^2a,v)=J(\Suc^2 a;u.v)$, with the hole in the body, is guarded.

\begin{lemma}[Bootstrap from two equations]\label{ctr:lem:gboot}
Using only \eqref{ctr:eq:gB}, \eqref{ctr:eq:gR} and the equality rules of
\cref{ctr:lem:gstruct}(2), $\mathcal K_{\mathrm G}$ derives in every context, for all
terms,
\begin{equation}\label{ctr:eq:gf}
 f(0,a)=\Suc a,\qquad f(n,0)\equiv\Suc n,\qquad f(\Suc n,\Suc a)=a,
\end{equation}
and
\begin{alignat*}{3}
 W(0)&=2, &\qquad W(\Suc x)&=x, &\qquad Y(0)&=2,\quad Y(\Suc n)=0,\\
 W^3(0)&=0, & W^3(\Suc^3 a)&=a, & Z(0)&=1,\quad Z(\Suc n)=0,\\
 G(0,a)&=0, & G(\Suc n,a)&=a, & C(0,a,b)&=a,\quad C(\Suc n,a,b)=b,\\
 P(0)&=0, & P(\Suc n)&=n.
\end{alignat*}
Consequently $\Suc x=\Suc y\proves x=y$ and $0=1\proves a=b$ for all terms $a,b$.
\end{lemma}
\begin{proof}
In \eqref{ctr:eq:gf} the first equation is \eqref{ctr:eq:gB} for the state-free body $a$,
the second is the definition of $\Suc$ up to the name of the bound variable, and the third
is \eqref{ctr:eq:gR} with the constant $g\defeq a$. Hence $W(0)=f(1,0)\equiv\Suc 1=2$,
$W(\Suc x)=f(\Suc 0,\Suc x)=x$, $Y(0)=f(0,1)=\Suc 1$ and $Y(\Suc n)=f(\Suc n,\Suc 0)=0$. The
map $W$ removes one successor and cycles through $0,2,1$: $W(0)=2$, $W(2)=W(\Suc 1)=1$ and
$W(1)=W(\Suc 0)=0$, with congruence at each composition. Thus $W^3(0)=0$ and
$W^3(\Suc^3a)=a$.

At $c=0$ the two inner applications of $Y$ give $Y(Y(0))=Y(\Suc 1)=0$, and $W^2(0)=W(2)=1$;
at $c=\Suc n$ they give $Y(0)=2$, and $W^2(2)=W(1)=0$. For $G(0,a)$ the instruction is
$Z(0)=1=\Suc 0$, so congruence in the template $W^3(f(\Suc^2a,v))$ and \eqref{ctr:eq:gR} give
$W^3(f(\Suc(\Suc a),\Suc 0))=W^3(0)=0$. For $G(\Suc n,a)$ the instruction is $0$, and
$f(\Suc^2a,0)\equiv\Suc^3a$, so the value is $W^3(\Suc^3a)=a$. For $C(0,a,b)$ the instruction
$G(0,\Suc^4b)$ is $0$, giving $W^3(\Suc^3 a)=a$. For $C(\Suc n,a,b)$ the instruction is
$\Suc^4b$, and $f(\Suc(\Suc a),\Suc(\Suc^3b))=\Suc^3b$ by the third equation of
\eqref{ctr:eq:gf}, so the value is $W^3(\Suc^3b)=b$. The rows of $P$ follow:
$P(0)=C(0,0,W(0))=0$ and $P(\Suc n)=C(\Suc n,0,W(\Suc n))=W(\Suc n)=n$.

From $\Suc x=\Suc y$ congruence in $W$ gives $W(\Suc x)=W(\Suc y)$, hence $x=y$. From $0=1$,
\rulename{Tr} with the template $C(v,a,b)$ and the minor premise $C(0,a,b)=a$ gives
$a=C(1,a,b)$, and $C(1,a,b)=C(\Suc 0,a,b)=b$.
\end{proof}

No induction, numerical case distinction, coverage or previously available selector is used.
The finite decoder $W^3$ fixes $0$ and removes three successors; this is the property used in
$G$ and $C$.

\begin{proposition}[The conditional rules are derivable]\label{ctr:prop:grules}
For a guarded term $J(n;u.h)$ write $j_n\defeq J(n;u.h)$ and $j_{\Suc n}\defeq J(\Suc n;u.h)$.
The rules
\[
 \infrule{\Gamma\proves h\subst{j_n}{u}=\Suc b}{\Gamma\proves j_{\Suc n}=b}\ \rulename{Load},
 \qquad
 \infrule{\Gamma\proves h\subst{j_n}{u}=0}{\Gamma\proves j_{\Suc n}=\Suc j_n}\ \rulename{Inc}
\]
are derivable in $\mathcal K_{\mathrm G}$ from \eqref{ctr:eq:gB}, \eqref{ctr:eq:gR} and
\rulename{Tr}, without \rulename{QE}, \rulename{DE}, \rulename{DI} and \rulename{QI}. They
are the native forms of \eqref{ctr:eq:load} and \eqref{ctr:eq:inc}.
\end{proposition}
\begin{proof}
\rulename{Load}. If $h\equiv\Suc g$, \eqref{ctr:eq:gR} gives $j_{\Suc n}=g\subst{j_n}{u}$, and
the premise $\Suc(g\subst{j_n}{u})=\Suc b$ gives $g\subst{j_n}{u}=b$ by injectivity
(\cref{ctr:lem:gboot}). If $u\notin\FV(h)$, write $h\equiv a$; the premise is $a=\Suc b$,
and congruence in the template $f(\Suc n,v)$ gives
$j_{\Suc n}=f(\Suc n,a)=f(\Suc n,\Suc b)=b$ by \eqref{ctr:eq:gf}.

\rulename{Inc}. If $u\notin\FV(h)$, write $h\equiv a$; the premise is $a=0$. Congruence in
$f(n,v)$ and in $f(\Suc n,v)$ gives $j_n=f(n,0)\equiv\Suc n$ and
$j_{\Suc n}=f(\Suc n,0)\equiv\Suc(\Suc n)$, and congruence in $\Suc v$ gives
$\Suc j_n=\Suc(\Suc n)$; transitivity concludes. If $h\equiv\Suc g$, the premise is
$\Suc(g\subst{j_n}{u})=0$; congruence in $Z$ and \cref{ctr:lem:gboot} give
$0=Z(\Suc(g\subst{j_n}{u}))=Z(0)=1$, and the equational explosion of \cref{ctr:lem:gboot}
gives the conclusion.
\end{proof}

\Cref{ctr:prop:regrefl} derives reflexivity from the conditional rules; in
$\mathcal K_{\mathrm G}$ reflexivity comes from \eqref{ctr:eq:gR}
(\cref{ctr:lem:gstruct}), and the conditional rules come from the two equations and
\rulename{Tr}. The body restriction is what permits their removal from the rules.

\subsection{Initialized iteration}

\begin{theorem}[Initialized iteration]\label{ctr:thm:giter}
Let $a$ be a guarded term and $h(v)$ a guarded template, with all passive parameters fixed.
With $w$ fresh put
\[
 H(w)\defeq\Suc\Bigl(C\bigl(w,\,a,\,\Suc^2\bigl(h(P^2(w))\bigr)\bigr)\Bigr),\qquad
 A(k)\defeq J(k;w.H(w)),\qquad
 I_{\mathrm G}(n,a;v.h)\defeq P^2\bigl(A(n)\bigr).
\]
These terms are guarded, since $H$ is literally successor-headed, even when $h$ depends on
its variable. $\mathcal K_{\mathrm G}$ derives in every context, for every term $n$,
\begin{equation}\label{ctr:eq:gIC}
 I_{\mathrm G}(0,a;v.h)=a,\qquad I_{\mathrm G}(\Suc n,a;v.h)=h\bigl(I_{\mathrm G}(n,a;v.h)\bigr).
\end{equation}
The proof uses \eqref{ctr:eq:gB}, \eqref{ctr:eq:gR}, equality and one equation induction,
and no coverage, pairing or logical connective.
\end{theorem}
\begin{proof}
By the zero row of $C$, $H(0)=\Suc a$, and \eqref{ctr:eq:gB} gives $A(0)=\Suc(H(0))=\Suc^2a$;
hence $P^2(A(0))=a$ by the rows of $P$. This is the first equation. Let $k$ be fresh for
the context, $n$, $a$ and $h$, and consider the literal equation
\begin{equation}\label{ctr:eq:ginv}
 A(k)=\Suc^2\bigl(P^2(A(k))\bigr).
\end{equation}
At $k=0$ both sides equal $\Suc^2a$. For the step assume \eqref{ctr:eq:ginv} and write
$q\defeq P^2(A(k))$. \eqref{ctr:eq:gR} for the successor-headed body $H$ gives
\begin{equation}\label{ctr:eq:gstep}
 A(\Suc k)=C\bigl(A(k),\,a,\,\Suc^2(h(q))\bigr).
\end{equation}
Transport of \eqref{ctr:eq:ginv} in the first argument of $C$ makes the control
$\Suc(\Suc q)$, and the successor row of $C$ gives $A(\Suc k)=\Suc^2(h(q))$. Applying $P$
twice gives $P^2(A(\Suc k))=h(q)$, and congruence in $\Suc^2v$ gives \eqref{ctr:eq:ginv} at
$\Suc k$. Equation induction (\cref{ctr:lem:gstruct}(3)) proves \eqref{ctr:eq:ginv} at every
term, in particular at $n$. Repeating the step calculation with this instance at $n$ in
place of the induction hypothesis gives $A(\Suc n)=\Suc^2(h(P^2(A(n))))$, and applying $P$
twice gives the second equation of \eqref{ctr:eq:gIC}.
\end{proof}

If the counter $n$ also occurs in $a$ or $h$, those occurrences stay passive: the induction
is on the fresh $k$, and $n$ is substituted only in the proved instance. In
\cref{ctr:prop:scalarreg} the shape of the state is a consequence of one computation; here it
is the invariant \eqref{ctr:eq:ginv}, proved by one induction on a literal equation.

\begin{lemma}[Unconditional recursion]\label{ctr:lem:grec}
Let $\mathrm{Sel}(c,a,b)$ and $\mathrm{Pd}(c)$ be guarded templates such that
$\mathcal K_{\mathrm G}$ proves $\mathrm{Sel}(0,a,b)=a$, $\mathrm{Sel}(\Suc c,a,b)=b$ and
$\mathrm{Pd}(\Suc c)=c$ for all terms $a,b,c$. For every guarded term $J(n;u.h)$, with
$j\defeq J(n;u.h)$, $\mathcal K_{\mathrm G}$ derives in every context
\[
 J(0;u.h)=\Suc\bigl(h\subst{0}{u}\bigr),\qquad
 J(\Suc n;u.h)=\mathrm{Sel}\bigl(h\subst{j}{u},\,\Suc j,\,\mathrm{Pd}(h\subst{j}{u})\bigr).
\]
In particular this holds for $\mathrm{Sel}=C$ and $\mathrm{Pd}=P$. The proof uses numerical
cases and no coverage.
\end{lemma}
\begin{proof}
The first equation is \eqref{ctr:eq:gB}. If $h\equiv\Suc g$, put $g'\defeq g\subst{j}{u}$;
\eqref{ctr:eq:gR} gives $J(\Suc n;u.h)=g'$, and
$\mathrm{Sel}(\Suc g',\Suc j,\mathrm{Pd}(\Suc g'))=\mathrm{Pd}(\Suc g')=g'$. Otherwise
$u\notin\FV(h)$; write $h\equiv a$, so that $h\subst{j}{u}\equiv a$. Let $x$ be a fresh
variable and consider the literal equation
\[
 E(x)\colon\quad \mathrm{Sel}\bigl(x,\,\Suc(J(n;u.x)),\,\mathrm{Pd}(x)\bigr)=J(\Suc n;u.x),
\]
whose bodies are state-free. At $x=0$ the zero row of $\mathrm{Sel}$ makes the left side
equal to $\Suc(J(n;u.0))$, which is literally $\Suc(\Suc n)$; the right side $J(\Suc n;u.0)$
is the same term up to the name of the bound variable. At $x=\Suc z$, with $z$
fresh, the left side is $\mathrm{Pd}(\Suc z)=z$ and the right side is $z$ by
\eqref{ctr:eq:gR} with the constant $g\defeq z$. Numerical cases
(\cref{ctr:lem:gstruct}(3)) give $E(a)$, which is the claim, since
$J(n;u.x)\subst{a}{x}\equiv J(n;u.a)=j$ and $u\notin\FV(a)$.
\end{proof}

\subsection{Equivalence with the native core}

\begin{definition}[Translations]\label{ctr:def:gtrans}
The translation $s\mapsto s^\lambda$ from terms of $\Nat$ to guarded terms fixes variables
and $0$ and sets
\[
 (\Suc s)^\lambda\defeq J(s^\lambda;v.0),\qquad
 \It{n}{a}{u}{h}^\lambda\defeq I_{\mathrm G}(n^\lambda,a^\lambda;u.h^\lambda).
\]
The translation $t\mapsto t^\mu$ from guarded terms to terms of $\Nat$ fixes variables and
$0$ and sets, with $\mathrm{pd}$ the predecessor term of \cref{base:thm:pred} (written $P$
there) and $\sel$ the selector of \cref{base:sec:terms},
\[
 J(n;u.h)^\mu\defeq
 \It{n^\mu}{\Suc\bigl(h^\mu\subst{0}{u}\bigr)}{u}{\sel\bigl(h^\mu,\Suc u,\mathrm{pd}(h^\mu)\bigr)}.
\]
Both are extended to native formulas and contexts so that they fix $=$, $\Q{y}$ and
$\D{x}{y}$. Put $\sigma(x)\defeq(\Suc x)^\mu=\It{x}{\Suc 0}{u}{\sel(0,\Suc u,\mathrm{pd}(0))}$.
\end{definition}

Both translations are compositional, so they commute with capture-avoiding substitution and
preserve the free variables of a term; they map literal equations to literal equations.
Every $\lambda$-image is guarded, by \cref{ctr:thm:giter,ctr:lem:gsubst}. The term
$(\Suc s)^\lambda$ is literally the defined successor of $s^\lambda$, while $(\Suc t)^\mu$
is $\sigma(t^\mu)$, which is only provably the successor.

\begin{lemma}[The reverse translation of the axioms]\label{ctr:lem:gsigma}
$\Nat$ proves $\sigma(t)=\Suc t$ for every term $t$, and every $\mu$-translation of an
instance of \eqref{ctr:eq:gB} or \eqref{ctr:eq:gR}, in every context.
\end{lemma}
\begin{proof}
\emph{Successor.} Apply \cref{base:lem:eqind} to $\sigma(z)=\Suc z$. The base is
\rulename{I0}. In the step \rulename{IS} and \eqref{base:eq:sel} give
$\sigma(\Suc z)=\sel(0,\Suc\sigma(z),\mathrm{pd}(0))=\Suc\sigma(z)$, and congruence along the
hypothesis gives $\Suc\Suc z$.

\emph{\eqref{ctr:eq:gB}.} The left side translates to a term whose counter is $0$, so
\rulename{I0} gives $\Suc(h^\mu\subst{0}{u})$. The right side translates to
$\sigma(h^\mu\subst{0}{u})$, since $(h\subst{0}{u})^\mu=h^\mu\subst{0}{u}$, and this equals
$\Suc(h^\mu\subst{0}{u})$ by the first part.

\emph{\eqref{ctr:eq:gR}.} Put $s\defeq(\Suc g)^\mu=\sigma(g^\mu)$ and
$K(c)\defeq\It{c}{\Suc(s\subst{0}{u})}{u}{\sel(s,\Suc u,\mathrm{pd}(s))}$, so that
$J(m;u.\Suc g)^\mu=K(m^\mu)$ for every $m$. The left side translates to $K(\sigma(n^\mu))$.
By the first part and congruence it equals $K(\Suc n^\mu)$, and \rulename{IS} gives, with
$k\defeq K(n^\mu)$ and $g'\defeq g^\mu\subst{k}{u}$,
\[
 K(\Suc n^\mu)=\sel\bigl(\sigma(g'),\Suc k,\mathrm{pd}(\sigma(g'))\bigr)
 =\sel\bigl(\Suc g',\Suc k,\mathrm{pd}(\Suc g')\bigr)=\mathrm{pd}(\Suc g')=g',
\]
by the first part, \eqref{base:eq:sel} and \cref{base:thm:pred}. The right side translates to
$g^\mu\subst{K(n^\mu)}{u}=g'$.
\end{proof}

\begin{theorem}[Guarded registers and the native core]\label{ctr:thm:gequiv}
\leavevmode
\begin{enumerate}
\item If $\Gamma\proves_{\Nat}A$ then $\Gamma^\lambda\proves_{\mathcal K_{\mathrm G}}A^\lambda$.
If $\Gamma\proves_{\mathcal K_{\mathrm G}}A$ then $\Gamma^\mu\proves_{\Nat}A^\mu$.
\item $\mathcal K_{\mathrm G}\proves t^{\mu\lambda}=t$ for every guarded term $t$, and
$\Nat\proves s^{\lambda\mu}=s$ for every term $s$ of $\Nat$. Every formula is
interderivable with its double translation, in $\mathcal K_{\mathrm G}$ and in $\Nat$.
\item $\Gamma\proves_{\mathcal K_{\mathrm G}}A$ if and only if $\Gamma^\mu\proves_{\Nat}A^\mu$,
and $\Gamma\proves_{\Nat}A$ if and only if $\Gamma^\lambda\proves_{\mathcal K_{\mathrm G}}A^\lambda$.
Moreover $\proves_{\mathcal K_{\mathrm G}}0=1$ if and only if $\proves_{\Nat}0=\Suc 0$.
\item Consequently $\mathcal K_{\mathrm G}$ is bi-interpretable with $\Teq$ and with $\iPRA$,
by the compositions of $\lambda$ and $\mu$ with the translations of \cref{base:thm:bridge},
with whole domains, actual equality and numerical composites of graph $x=y$, and
$\mathcal K_{\mathrm G}$, $\Nat$, $\Teq$ and $\iPRA$ are equiconsistent.
\end{enumerate}
\end{theorem}
\begin{proof}
(1) \emph{From $\Nat$ to $\mathcal K_{\mathrm G}$.} Induct on the derivation. \rulename{Hyp}
is unchanged. An instance of \rulename{I0} translates to the first equation of
\eqref{ctr:eq:gIC}. An instance of \rulename{IS} translates to
$I_{\mathrm G}(J(n^\lambda;v.0),a^\lambda;u.h^\lambda)=
h^\lambda\subst{I_{\mathrm G}(n^\lambda,a^\lambda;u.h^\lambda)}{u}$, since $\lambda$
commutes with substitution, and this is the second equation of \eqref{ctr:eq:gIC}.
\rulename{Tr} with the template $h$ translates to \rulename{Tr} with the guarded template
$h^\lambda$. \rulename{QE} and \rulename{DE} translate to themselves, with the side
condition of \rulename{DE} preserved since $\lambda$ preserves free variables. In
\rulename{DI} and \rulename{QI} the translated conclusion of the premise is
$B^\lambda(t^\lambda,J(z;v.0))$, respectively $E^\lambda(J(z;v.0))$, which is literally the
premise of the rule of $\mathcal K_{\mathrm G}$; $E^\lambda$ is a literal equation.

\emph{From $\mathcal K_{\mathrm G}$ to $\Nat$.} Induct on the derivation. \rulename{Hyp},
\rulename{Tr}, \rulename{QE} and \rulename{DE} translate to themselves, as above. The
computation equations are \cref{ctr:lem:gsigma}. In \rulename{DI} the induction hypothesis
gives $\Gamma^\mu,B^\mu(t^\mu,0)\proves_{\Nat}B^\mu(t^\mu,\sigma(z))$. The Leibniz rule
(\cref{base:thm:leibniz}) along $\sigma(z)=\Suc z$ (\cref{ctr:lem:gsigma}) gives
$B^\mu(t^\mu,\Suc z)$, and \rulename{DI} of $\Nat$ applies; $z$ is fresh since $\mu$
preserves free variables. \rulename{QI} is the same with the literal equation $E^\mu$.

(2) \emph{Terms of $\Nat$.} Induct on $s$. For $s=\Suc s_0$, $s^{\lambda\mu}=\sigma(s_0^{\lambda\mu})$,
which equals $\Suc s_0$ by \cref{ctr:lem:gsigma} and the induction hypothesis. For
$s=\It{n}{a}{u}{h}$ let $\Phi(k)\defeq I_{\mathrm G}(k,a^\lambda;u.h^\lambda)^\mu$ with $k$
fresh, so that $s^{\lambda\mu}=\Phi(n^{\lambda\mu})$. The $\mu$-translations of the two
equations \eqref{ctr:eq:gIC} at $k$ are derivable in $\Nat$ by part~(1); they are
$\Phi(0)=a^{\lambda\mu}$ and $\Phi(\sigma(k))=h^{\lambda\mu}\subst{\Phi(k)}{u}$, and
\cref{ctr:lem:gsigma} turns the second into $\Phi(\Suc k)=h^{\lambda\mu}\subst{\Phi(k)}{u}$.
By the induction hypotheses $a^{\lambda\mu}=a$ and $h^{\lambda\mu}=h$, the latter with $u$
free. The term $\It{k}{a}{u}{h}$ has the base $a$ and the update $h$. The proof of
\cref{ctr:lem:unique} uses equality, an instance of a derived equation and equation
induction, and is available in $\Nat$ (\cref{base:lem:eq,base:lem:eqind}); it gives
$\Phi(k)=\It{k}{a}{u}{h}$, and the instance at $n^{\lambda\mu}=n$ is the claim.

\emph{Guarded terms.} Induct on $t$; variables and $0$ are fixed. Let $t=J(n;u.h)$. With
$\mathrm{Sel}(c,a,b)\defeq I_{\mathrm G}(c,a;w.b)$, $w$ not free in $b$, and
$\mathrm{Pd}(c)\defeq\bigl(\mathrm{pd}(x)\bigr)^\lambda\subst{c}{x}$,
\[
 t^{\mu\lambda}=\Psi(n^{\mu\lambda}),\qquad
 \Psi(k)\defeq I_{\mathrm G}\Bigl(k,\ \Suc\bigl(h^{\mu\lambda}\subst{0}{u}\bigr);\
 u.\,\mathrm{Sel}\bigl(h^{\mu\lambda},\Suc u,\mathrm{Pd}(h^{\mu\lambda})\bigr)\Bigr),
\]
because $\lambda$ commutes with substitution and $\sel(c,a,b)$ is $\It{c}{a}{w}{b}$. By
\eqref{ctr:eq:gIC}, $\mathrm{Sel}(0,a,b)=a$ and $\mathrm{Sel}(\Suc c,a,b)=b$. By
\cref{base:thm:pred}, $\Nat\proves\mathrm{pd}(\Suc x)=x$, and part~(1) gives
$\mathcal K_{\mathrm G}\proves\mathrm{Pd}(\Suc x)=x$, since $(\Suc x)^\lambda$ is literally
$\Suc x$; substitution gives it for every $c$. So \cref{ctr:lem:grec} applies to
$\mathrm{Sel}$ and $\mathrm{Pd}$. Apply equation induction to $\Psi(k)=J(k;u.h)$ with $k$
fresh. At $0$, \eqref{ctr:eq:gIC}, the induction hypothesis $h^{\mu\lambda}=h$ with $u$ free,
substitution of $0$ for $u$ and \eqref{ctr:eq:gB} give
$\Psi(0)=\Suc(h\subst{0}{u})=J(0;u.h)$. In the step, with $j\defeq J(k;u.h)$ and the
hypothesis $\Psi(k)=j$, \eqref{ctr:eq:gIC}, congruence and the induction hypothesis
substituted at $j$ give
$\Psi(\Suc k)=\mathrm{Sel}(h\subst{j}{u},\Suc j,\mathrm{Pd}(h\subst{j}{u}))$, which is
$J(\Suc k;u.h)$ by \cref{ctr:lem:grec}. The instance at $n^{\mu\lambda}=n$ is the claim.
This compares the recurrence of the double translation directly with the original register;
no intermediate term $J(n;u.h^{\mu\lambda})$, which need not be guarded, is formed.

\emph{Formulas.} Induct on the formula. For an equation use the term identities with
symmetry and transitivity. For $\Q{y}B$ and $\D{x}{y}B$ the induction hypothesis gives
interderivability of $B$ and its double translation, with $x,y$ free, and the congruence of
$Q$ and $D$ (\cref{base:lem:bindercong}, and \cref{ctr:lem:gstruct}(3) for
$\mathcal K_{\mathrm G}$) applies.

(3) If $\Gamma^\mu\proves_{\Nat}A^\mu$, then
$\Gamma^{\mu\lambda}\proves_{\mathcal K_{\mathrm G}}A^{\mu\lambda}$ by (1), and (2) with cut
(\cref{ctr:lem:gstruct}(1)) gives $\Gamma\proves_{\mathcal K_{\mathrm G}}A$. The other
equivalence is symmetric, with \cref{base:lem:cut}. The formula $(0=\Suc 0)^\lambda$ is
literally $0=1$, and $(0=1)^\mu$ is $0=\sigma(0)$, which is interderivable in $\Nat$ with
$0=\Suc 0$ by \cref{ctr:lem:gsigma}.

(4) By (3) and \cref{base:thm:bridge}(i) and (iii), derivability in
$\mathcal K_{\mathrm G}$ corresponds to derivability in $\Teq$ and in $\iPRA$ under the
composite translations. For the round trips, a formula $A$ of $\mathcal K_{\mathrm G}$ is
interderivable with $A^{\mu\lambda}$ by (2), and the round trip of \cref{base:thm:bridge}
makes $A^\mu$ interderivable in $\Nat$ with its double translation through $\Teq$ or
$\iPRA$; by (1) the $\lambda$-images of interderivable formulas are interderivable. In the
other direction a formula $\psi$ of $\iPRA$ is equivalent to its round trip through $\Nat$
by \cref{base:thm:bridge}(iii), and the formula of $\Nat$ obtained on the way is
interderivable with its $\lambda\mu$-translation by (2); since derivability in $\Nat$ is
preserved by the translation to $\iPRA$, $\iPRA$ proves the equivalence of $\psi$ with its
round trip through $\mathcal K_{\mathrm G}$. The same argument applies to $\Teq$ with
\cref{base:thm:bridge}(i). The translations $\lambda$ and $\mu$ fix the variables, the
binders and equality, and on terms the composites are provably the identity by (2); with
\cref{base:thm:bridge}(iv) the numerical composites have the graph $x=y$. Equiconsistency
is (3) with \cref{base:thm:bridge}(v).
\end{proof}

\begin{remark}[The route of the package]\label{ctr:rem:groute}
\pkg{guarded\_register} proves \cref{ctr:thm:gequiv} with the same two translations. It
interprets $\mathcal K_{\mathrm G}$ in the ordinary scalar theory by $\mu$, reading
\rulename{DI} by coverage and \rulename{QI} by equation induction; it reconstructs the
logical connectives inside $\mathcal K_{\mathrm G}$ by gates and a selector of formulas built
from the term $C$, which is the construction of \cref{base:def:macros} with $C$ in place of
$\sel$; and it proves \cref{ctr:lem:grec} for a state-free body by numerical coverage of the
instruction and elimination of the derived disjunction. The proof above replaces coverage by
numerical cases at a fresh variable, so the logical interface is not needed for the
equivalence; the logic of $\mathcal K_{\mathrm G}$ is that of $\Nat$, transferred by
\cref{ctr:thm:gequiv}(3). The two composites are proved inside the theories, by equation
induction, and not by agreement on the standard model.
\end{remark}

\begin{remark}[The ordinary realization]\label{ctr:rem:gordinary}
\pkg{guarded\_register} also states an ordinary theory $\mathsf T_{\mathrm G}$: intuitionistic
first-order logic with equality and an independent $\bot$ over guarded terms, the universal
closures of \eqref{ctr:eq:gB} and \eqref{ctr:eq:gR}, coverage
$\forall x\,(x=0\lor\exists y\,x=\Suc y)$, the separation $0=1\to\bot$, and the rule
\eqref{ctr:eq:EI} of induction for literal equations. With $\lambda$ and $\mu$ extended to
commute with the connectives, it is bi-interpretable with $\Teq$ by the argument of
\cref{ctr:thm:gequiv}. In $\Teq$ the $\mu$-translations of the two equations hold by
\cref{ctr:lem:gsigma} and \cref{base:thm:bridge}(i), and coverage, separation and
induction translate to the corresponding principles of $\Teq$ after transport along
$\sigma(x)=\Suc x$. In $\mathsf T_{\mathrm G}$ the proofs of
\cref{ctr:lem:gboot,ctr:thm:giter} use only equality, the two equations and induction on a
literal equation; the axioms of $\Teq$ translate to \eqref{ctr:eq:gIC}, to coverage and to
separation, the last two literally. The term composites are proved as in
\cref{ctr:thm:gequiv}(2), with coverage in place of numerical cases, and the formula
composites follow by transport. The separation axiom is needed: a one-element structure
satisfies the two equations, coverage and every instance of induction, and refutes
$0=1\to\bot$.
\end{remark}

\begin{remark}[A closure meaning that the eliminations do not fit]\label{ctr:rem:gclosure}
\pkg{guarded\_register} records an obstruction for a different shortcut. Keep the
elimination \rulename{QE} and the introduction \rulename{QI} of $\Nat$, and read the closure
$\Q{y}E$ of a literal equation progressively, as $\forall y\,(E(y)\to E(\Suc y))$. The rules
are sound for this reading, but not complete for it: under the reading of
\cref{base:sec:reading}, $\Q{y}E$ is $E(0)\to\forall y\,E(y)$, every rule of $\Nat$ is
sound, and for $E(y)\defeq(y=1)$ the formulas $\Q{y}E$ and $E(1)$ are true on the natural
numbers while $E(2)$ is false. Hence $\Q{y}E,E(1)\not\proves_{\Nat}E(2)$, although the
sequent is valid under the progressive reading. A faithful progressive reading requires a
corresponding change of the elimination rule.
\end{remark}

\subsection{The strict checker}

The checker of \pkg{guarded\_register} is a standalone module that imports only the JSON
library, reads a certificate $[\mathit{goal},\mathit{proof}]$ with a closed goal and prints
\texttt{ACCEPT} only after checking the goal and the proof; every check is a runtime
condition, so optimized Python does not remove it. Variables are one-based de Bruijn indices
and the integer $0$ is the constant. A guarded term is $[\texttt{j},n,h]$, with $h$ binding
the index $1$; the state-free test is \texttt{subst(h,(0,),1)==h}, which replaces the bound
state by $0$ and keeps all outer indices, and the alternative test is that $h$ is a $J$ node
with literal zero body. Both computation nodes validate the whole constructed $J$ term, so an
unguarded body cannot enter through an axiom, and the node for \eqref{ctr:eq:gR} takes $n$
and $g$, has no equation premise and no output witness. The six logical nodes are those of
the matched control for iterator terms, which has the term tags $\Suc$ and $I$ and no body
restriction. The validation runs $108$ subprocess checks in ordinary and optimized Python
($48$ acceptances, $60$ rejections, no mismatch), among them rejected certificates that use
the removed increment and conditional load, unguarded bodies in auxiliary nodes and
quantified induction, and $12$ parameterized certificates of the equations of
\cref{ctr:lem:gboot} for $f$, $W$, $Z$, $G$, $C$ and $P$, checked with an empty context. These
are finite evidence for the implementation.

\section{The seeded recursor: unrestricted calculi and six rules}\label{ctr:sec:six}

This section draws on \pkg{canonical\_recursors}.

The native formulas are those of \cref{base:sec:kernel}, $A ::= (s=t)\mid\Q{y}A\mid\D{x}{y}A$,
over the terms of the seeded recursor with one bound variable, with the reading of
\cref{base:sec:reading}. Contexts are finite reusable lists of formulas, and the designated
contradiction is $0=\Suc 0$. The six-rule calculus has two tests on literal syntax, which
are not stable under substitution. The comparison with the raw calculus is therefore
established first for two unrestricted calculi, and the six rules are recovered afterwards.

\subsection{The raw and the canonical unrestricted calculus}

\begin{definition}\label{ctr:def:free}
The \emph{raw calculus} $\mathcal R_{\mathrm{raw}}$ uses raw terms and ordinary substitution.
Its rules are:
\begin{enumerate}
\item \rulename{Hyp}: $\Gamma\proves A$ if $A\in\Gamma$;
\item the two computation equations \eqref{ctr:eq:r1zero} and \eqref{ctr:eq:r1succ}, in
every context;
\item transport in every native formula template: from $s=t$ and $F\subst{s}{v}$ infer
$F\subst{t}{v}$;
\item \rulename{QE}: from $\Q{y}B(y)$ and $B(0)$ infer $B(t)$;
\item \rulename{DE}: from $\D{x}{y}B(x,y)$ and a derivation of $A$ from
$\Gamma,\Q{y}B(b,y)$ infer $A$, where $b$ is fresh for $\Gamma$, $A$ and the parameters of
$B$;
\item introduction by cases: from $\Gamma,B(t,0)\proves B(t,\Suc z)$ infer
$\D{x}{y}B(x,y)$;
\item introduction by progress: from $\Gamma,E(t,z)\proves E(t,\Suc z)$ infer
$\D{x}{y}E(x,y)$, for a literal equation $E$.
\end{enumerate}
In the two introductions $z$ is fresh for $\Gamma$, the witness $t$ and the parameters of the
matrix. The \emph{canonical calculus} $\mathcal C_\nu$ has the same rules on zero-body
canonical formulas with normalized substitution, except that its only computation schema is
\begin{equation}\label{ctr:eq:Rplus}
 R(\Suc n,a;u.h)=h\langle R(n,a;u.h)/u\rangle . \tag{$R_+$}
\end{equation}
\end{definition}

The left counter of \eqref{ctr:eq:Rplus} is never literally $0$, so all displayed terms are
canonical. The raw zero computation has become an identity of spelling. The raw calculus is
the kernel $\Nat$ on the seeded terms in an equivalent rule set: transport of formulas is
the Leibniz rule of \cref{base:sec:structural}, and progress concludes a paired formula
instead of a closure.

\begin{lemma}[Structural rules]\label{ctr:lem:struct}
Weakening, renaming of eigenvariables, cut of an ordinary assumption and capture-free
numerical substitution are admissible in $\mathcal R_{\mathrm{raw}}$ and $\mathcal C_\nu$.
\end{lemma}
\begin{proof}
Induct on derivations after freshening the local eigenparameters. For $\mathcal C_\nu$,
\cref{ctr:lem:coherence} shows that normalized substitution maps every rule instance to a
rule instance, including the instances of \eqref{ctr:eq:Rplus}. The distinction between the
two introductions depends only on the root of the matrix and is stable.
\end{proof}

\begin{lemma}[Elementary resources]\label{ctr:lem:resources}
Both unrestricted calculi derive reflexivity, symmetry, transitivity and congruence of
equality, introduction of $\Q{y}B$ by cases, numerical cases, and induction on literal
equations with an arbitrary context.
\end{lemma}
\begin{proof}
Let $\operatorname{wr}(a)\defeq R(\Suc 0,0;u.a)$ with $u\notin\FV(a)$. The successor
computation gives $\operatorname{wr}(a)=a$, also in $\mathcal C_\nu$. Transport of this
equation along itself in the template $v=a$ gives $a=a$. The templates $v=s$, $s=v$ and
$f(s)=f(v)$ give symmetry, transitivity and congruence.

From $\Gamma,B(0)\proves B(\Suc z)$, introduction by cases with an unused witness coordinate
gives $\D{w}{y}B(y)$, and \rulename{DE} followed by \rulename{Hyp} exports $\Q{y}B(y)$. With
proofs of $B(0)$ and $B(\Suc z)$, \rulename{QE} gives $B(t)$; this is numerical cases. The
same construction with introduction by progress gives $\Q{y}E(y)$ from
$\Gamma,E(z)\proves E(\Suc z)$, and with a proof of $E(0)$ it gives $E(t)$.
\end{proof}

\begin{lemma}[Body extensionality, native form]\label{ctr:lem:bodynative}
In both unrestricted calculi, if $\Gamma\proves h(v)=k(v)$ for a variable $v$ fresh for
$\Gamma$, then $\Gamma\proves R(n,a;u.h(u))=R(n,a;u.k(u))$.
\end{lemma}
\begin{proof}
Apply induction on literal equations to the equation $R(z,a;u.h)=R(z,a;u.k)$ at a fresh
counter $z$, with the initializer and the parameters fixed. At $0$ both terms are $\Suc a$,
by the zero computation in $\mathcal R_{\mathrm{raw}}$ and by identity of canonical spelling
and reflexivity in $\mathcal C_\nu$. At a successor the computation exposes the two updates.
Replace the old state in $h$ by the induction equation, and use the premise at the second
old state; this instance of the premise is justified by \cref{ctr:lem:struct}.
\end{proof}

The lemma is necessary. Changing a normalized body changes a finite description; it is not
congruence for one fixed function symbol.

\begin{lemma}[Normalization is provable in the raw calculus]\label{ctr:lem:normaleq}
For every raw term $t$, $\mathcal R_{\mathrm{raw}}\proves t=\nu(t)$. For every raw native
formula $A$, the raw calculus derives $A\proves\nu(A)$ and $\nu(A)\proves A$, uniformly in
the free parameters.
\end{lemma}
\begin{proof}
For terms induct on $t$. The counter and the initializer change by congruence. The body
changes by \cref{ctr:lem:bodynative} and the induction hypothesis with the bound variable
free. If the counter is literally $0$, the raw zero computation replaces the body by $0$.

For formulas induct on the formula. An equation follows from the two endpoint identities.
For $\Q{y}A(y)$, under this assumption introduce $\Q{y}\nu(A)(y)$ by cases: convert the
assumed normalized base to the old base, apply \rulename{QE} at the fresh successor and
convert back. The reverse direction is identical. For $\D{x}{y}A$, eliminate the old
formula at a fresh witness, introduce the new one at that witness by cases, convert the new
base to the old base, use the exported closure, and convert the successor conclusion. These
are derived congruence arguments for the two binders; no ordinary quantifier logic is
assumed.
\end{proof}

\begin{theorem}[Canonical and raw consequence]\label{ctr:thm:canonical}
For zero-body canonical $\Gamma,A$,
\[
 \Gamma\proves_{\mathcal C_\nu}A\iff\Gamma\proves_{\mathcal R_{\mathrm{raw}}}A .
\]
More generally, $\nu$ sends every raw derivation to a canonical derivation of the
normalized sequent.
\end{theorem}
\begin{proof}
From raw to canonical, normalize every sequent of the derivation. By
\cref{ctr:lem:coherence} an instance of transport, \rulename{QE}, \rulename{DE} or an
introduction becomes an instance of the same rule with normalized substitution. A raw zero
computation becomes $\Suc a=\Suc a$, which is derivable by \cref{ctr:lem:resources}, and a
successor computation becomes an instance of \eqref{ctr:eq:Rplus}.

Conversely, regard canonical expressions as raw ones and recurse on the canonical
derivation. A canonical rule instance differs from the raw instance only by normalization of
substituted expressions. \Cref{ctr:lem:normaleq} converts premises, local assumptions and
conclusions in both directions, and cut handles the local assumptions. For example,
canonical transport converts the minor premise $\nu(E[s/v])$ to $E[s/v]$, applies raw
transport and converts $E[t/v]$ to its normalization. Freshness is preserved because $\nu$
introduces no free variable. No substitution theorem for a restricted rule enters either
direction.
\end{proof}

\subsection{The six rules}

For a native formula define the root reversal by
$(s=t)^\dagger=(t=s)$, $(\Q{y}B)^\dagger=\Q{y}B$ and $(\D{x}{y}B)^\dagger=\D{x}{y}B$. It
commutes with normalized substitution.

\begin{definition}[The calculus $\mathcal K_6$]\label{ctr:def:six}
Contexts are finite reusable lists of zero-body canonical native formulas. The rules are:
\begin{enumerate}
\item \rulename{Hyp}: $\Gamma\proves A$ if $A\in\Gamma$.
\item Successor computation: every instance of \eqref{ctr:eq:Rplus}, in every context.
\item Crossed equational transport, for a template $E(v)$ that is a literal equation
(constant templates are allowed):
\[
 \infrule{\Gamma\proves s=t\qquad \Gamma\proves E\langle s/v\rangle}
         {\Gamma\proves (E\langle t/v\rangle)^\dagger}
 \qquad (s\not\equiv 0). \tag{X}
\]
\item Crossed closure elimination:
\[
 \infrule{\Gamma\proves \Q{y}B(y)\qquad \Gamma\proves B\langle 0/y\rangle}
         {\Gamma\proves (B\langle t/y\rangle)^\dagger}. \tag{QX}
\]
\item Paired elimination, with $b$ fresh for $\Gamma$, $A$ and the parameters of $B$:
\[
 \infrule{\Gamma\proves \D{x}{y}B(x,y)\qquad \Gamma,\ \Q{y}B\langle b/x\rangle\proves A}
         {\Gamma\proves A}. \tag{DE}
\]
\item One introduction, with $z$ fresh for $\Gamma$, $t$ and the parameters of $B$:
\[
 \infrule{\Gamma,\ B\langle t/x,\eta/y\rangle\proves B\langle t/x,\Suc z/y\rangle}
         {\Gamma\proves \D{x}{y}B}
 \qquad
 \eta=\begin{cases} z, & B \text{ is a literal equation and } t\equiv 0,\\
                    0, & \text{otherwise.}\end{cases} \tag{WI}
\]
\end{enumerate}
\end{definition}

There is no primitive reflexivity, symmetry, numerical case rule, induction rule,
zero-computation rule, logical connective or rule of numerical substitution. On iterator
terms, one-term transport at sources other than the literal $0$ together with the single
witness-directed introduction is the seven-rule calculus \Lean{MPrf}, which is formalized
and proved equivalent to the kernel (\Lean{mprf\_iff\_nprf}). $\mathcal K_6$ crosses its two
eliminations, takes equation templates, and has no zero computation; what is specific to it
is the canonical syntax. Weakening, renaming of eigenvariables and cut of
an ordinary assumption are admissible in $\mathcal K_6$ by induction on derivations: cut
replaces uses of an assumption by proofs and changes no witness, root or source term, and
renaming preserves literal zero and roots. Numerical substitution is not claimed at this
stage.

\begin{lemma}[Equality in $\mathcal K_6$]\label{ctr:lem:sixeq}
$\mathcal K_6$ derives reflexivity, symmetry, transitivity, congruence, uncrossed transport
in every literal-equation template without the restriction on the source, and the ordinary
rule \rulename{QE}.
\end{lemma}
\begin{proof}
Keep $\operatorname{wr}(a)=R(\Suc 0,0;u.a)$ and put $c\defeq\operatorname{wr}(0)$. Each
$\operatorname{wr}(a)$ is not literally $0$, and $\operatorname{wr}(a)=a$ is an instance of
\eqref{ctr:eq:Rplus}.

\emph{Symmetry.} (X) with major premise $c=0$ and a constant template reverses any proved
equation. Since the reversal of a compound formula is the identity, (QX) followed by a
reversal is ordinary \rulename{QE}, and (X) followed by a reversal is uncrossed transport at
sources that are not literally $0$.

\emph{Reflexivity.} Use $\operatorname{wr}(a)=a$ as both premises of (X) with the template
$v=a$.

\emph{Zero join.} Let $p,q\not\equiv 0$ with $p=0$ and $q=0$. (X) with major $p=0$, minor
$\operatorname{wr}(p)=p$ and template $\operatorname{wr}(v)=p$ gives $p=c$; similarly
$q=c$. Reverse the second and transport $p=c$ along $c=q$, whose source $c$ is allowed, to
get $p=q$.

\emph{Transitivity.} Given $s=t$ and $t=r$: if $t\not\equiv 0$ transport $s=t$ along
$t=r$. If $t\equiv 0$ and both endpoints are not literally $0$, reverse the second equation
and use the zero join. If an endpoint is literally $0$, the conclusion is one of the
premises.

\emph{Congruence.} Along $s=t$ with $s\not\equiv 0$, transport the reflexive equation
$f(s)=f(s)$ in the template $f(s)=f(v)$. For $0=t$ with $t\not\equiv 0$, reverse to $t=0$
and transport $f(t)=f(t)$ in the template $f(v)=f(t)$, which gives $f(0)=f(t)$. If both
endpoints are literally $0$ use reflexivity.

\emph{Transport.} From $s=t$ and $a(s)=b(s)$, congruence and transitivity give
$a(t)=a(s)=b(s)=b(t)$. Every step instantiates a schema at the terms at hand; none
substitutes into a derivation.
\end{proof}

\begin{lemma}[Introduction by cases]\label{ctr:lem:sixcase}
Every instance of introduction by cases is derivable in $\mathcal K_6$.
\end{lemma}
\begin{proof}
(WI) is introduction by cases when the matrix is compound or the witness is not literally
$0$. Let $E$ be an equation matrix with the premise $\Gamma,E(0,0)\proves E(0,\Suc z)$. Use
the witness $c$, which is not literally $0$ and is provably $0$. Under the assumption
$E(c,0)$, \cref{ctr:lem:sixeq} gives $E(0,0)$; cut with the premise gives $E(0,\Suc z)$, and
transport back gives $E(c,\Suc z)$. (WI) at the witness $c$ is in case mode and concludes
the unchanged formula $\D{x}{y}E(x,y)$. The witness $c$ does not contain $z$.
\end{proof}

\begin{lemma}[Introduction by progress]\label{ctr:lem:sixprog}
Every instance of introduction by progress for a literal equation is derivable in
$\mathcal K_6$.
\end{lemma}
\begin{proof}
Let the premise be $\Gamma,E(t,z)\proves E(t,\Suc z)$ and let $w$ be unused in $E(t,y)$.
(WI) with witness $0$ is in progress mode and proves $\Gamma\proves\D{w}{y}E(t,y)$. Under
the additional assumption $E(t,0)$ eliminate this formula. The exported closure is
$\Q{y}E(t,y)$, independent of the fresh witness, and ordinary \rulename{QE} gives
$E(t,\Suc z)$. Closing (DE) gives the premise of \cref{ctr:lem:sixcase} at the witness $t$.
\end{proof}

\begin{lemma}[Transport in compound templates]\label{ctr:lem:sixtransport}
$\mathcal K_6$ derives transport in every native formula template.
\end{lemma}
\begin{proof}
Induct on the template. Equations are \cref{ctr:lem:sixeq}. For $\Q{y}B$, introduce the new
closure by cases with an unused witness as in \cref{ctr:lem:resources}: convert its base
back along the equation by the induction hypothesis, apply the old closure at the fresh
successor, and convert forward. For $\D{x}{y}B$, (DE) opens the old formula at a fresh
witness; introduce the new formula by cases at that witness, convert its base back, use the
exported closure, and convert the successor forward. The other direction of the equation is
obtained by symmetry. Only symmetry, cut and renaming are used.
\end{proof}

\begin{theorem}[Six-rule equivalence]\label{ctr:thm:six}
For zero-body canonical $\Gamma,A$,
\[
 \Gamma\proves_{\mathcal K_6}A\iff\Gamma\proves_{\mathcal C_\nu}A .
\]
\end{theorem}
\begin{proof}
\Cref{ctr:lem:sixeq,ctr:lem:sixcase,ctr:lem:sixprog,ctr:lem:sixtransport} provide every rule
of $\mathcal C_\nu$ inside $\mathcal K_6$. Conversely an instance of (WI) is one of the two
unrestricted introductions, and (X) and (QX) are transport and \rulename{QE} followed by the
derived symmetry when the conclusion is an equation; the other rules are common. No
admissibility of numerical substitution in $\mathcal K_6$ is presupposed; it now transfers
through the equivalence.
\end{proof}

\begin{remark}
Normalization preserves the literal-zero status of a whole term and the root of a formula,
so it does not alter a test of (X) or (WI). An arbitrary numerical substitution can alter
such a test: a variable witness can become $0$. Keeping the unrestricted equivalence of
\cref{ctr:thm:canonical} and the recovery of the restricted rules separate avoids an invalid
rule-by-rule substitution argument.
\end{remark}

\subsection{The bridge for the seeded calculi}

\begin{theorem}\label{ctr:thm:sixbridge}
Extend $\iota$ and $\kappa$ of \cref{ctr:def:translations}(1) to native formulas, fixing
$=$, $\Q{y}$ and $\D{x}{y}$.
\begin{enumerate}
\item If $\Gamma\proves_{\mathcal R_{\mathrm{raw}}}A$ then $\Gamma^\iota\proves_{\Nat}A^\iota$,
and if $\Gamma\proves_{\Nat}A$ then $\Gamma^\kappa\proves_{\mathcal R_{\mathrm{raw}}}A^\kappa$.
\item $\mathcal R_{\mathrm{raw}}\proves t^{\iota\kappa}=t$ and $\Nat\proves s^{\kappa\iota}=s$;
every native formula is interderivable with its double translation in both calculi.
\item For zero-body canonical $\Gamma,A$:
$\Gamma\proves_{\mathcal K_6}A$ if and only if $\Gamma^\iota\proves_{\Nat}A^\iota$.
\end{enumerate}
Consequently $\mathcal K_6$, $\mathcal C_\nu$ and $\mathcal R_{\mathrm{raw}}$ are sound and
faithful for $\mathsf T_R^1$ under the reading of \cref{base:sec:reading}, and each is
equiconsistent and bi-interpretable with $\iPRA$, with whole domains, actual equality and
numerical composites of graph $x=y$. On raw syntax the composite through canonical syntax is
provably, not literally, the identity; on canonical syntax $\nu$ is literally idempotent.
\end{theorem}
\begin{proof}
(1) From $\mathcal R_{\mathrm{raw}}$ to $\Nat$: the computation equations translate to the
computation equations of the iterator, using $\Nat\proves(\Suc t)^\iota=\Suc(t^\iota)$ and
the Leibniz rule of \cref{base:sec:structural}; transport of formulas is that Leibniz rule;
\rulename{QE}, \rulename{DE} and introduction by cases are rules of $\Nat$; and introduction
by progress is derived by forming $\Q{y}E(t,y)$ by progress, applying \rulename{QE} under
the assumption $E(t,0)$ and introducing by cases at $t$. From $\Nat$ to
$\mathcal R_{\mathrm{raw}}$: the computation equations of the iterator are
\cref{seed:thm:succ-iter} with \eqref{ctr:eq:seedpred}; the one-term transport of $\Nat$ is
transport followed by symmetry; and the progress rule of $\Nat$ is introduction by progress
with an unused witness followed by \rulename{DE} and \rulename{Hyp}, as in
\cref{ctr:lem:resources}.

(2) By \cref{ctr:lem:struct,ctr:lem:resources} the proof of \cref{ctr:lem:unique} is
available in $\mathcal R_{\mathrm{raw}}$ and in $\Nat$: it uses equality reasoning, an
instance of a derived equation and induction on a literal equation. The argument of
\cref{ctr:thm:transfer}(2) then gives the two term identities. A formula differs from its
double translation by these provable identities at its terms, with the bound variables of
the formula free; the binder-congruence argument of \cref{ctr:lem:normaleq} gives
interderivability.

(3) By \cref{ctr:thm:six,ctr:thm:canonical}, derivability in $\mathcal K_6$ is derivability
in $\mathcal R_{\mathrm{raw}}$. One direction is (1). For the other, a derivation of
$\Gamma^\iota\proves_{\Nat}A^\iota$ gives $\Gamma^{\iota\kappa}\proves A^{\iota\kappa}$ in
$\mathcal R_{\mathrm{raw}}$, and (2) with cut gives $\Gamma\proves A$. The consequences
follow from \cref{base:thm:bridge} and \cref{ctr:thm:transfer}.
\end{proof}

The mathematical change in this section is the canonical syntax and the proof that it
preserves substitution, derivability and both composites. The inventory decreases from
seven rules to six, and the successor constructor is gone, but the omitted
zero computation is the law $R(0,a;u.h)\equiv R(0,a;u.0)$ of the syntax, which every
substitution and comparison of a complete checker enforces.

\section{Five rules with clause contexts}\label{ctr:sec:five}

This section draws on \pkg{clausal\_canonical}.

The calculus of this section is the clause presentation of \cref{base:sec:clauses} with the
one-term transport of \cref{base:sec:kernel}, stated on the zero-body canonical terms of the
seeded recursor with one bound variable. The closure binder is no longer a formula
constructor, and the hypothesis rule and the closure elimination become one rule of clause
use. Only what the canonical syntax and the missing reflexivity rule require is proved.

\subsection{The calculus}

\begin{definition}[The calculus $\mathcal K_5$]\label{ctr:def:five}
Terms are zero-body canonical terms of the seeded recursor, and all substitutions are
normalized. Formulas are $F ::= (s=t)\mid\D{x}{y}F$. A context is a finite reusable list of
clauses $[A;\,y.\,B(y)]$ with $y\notin\FV(A)$; clauses are context objects, not formulas or
goals. Let
\[
 \operatorname{wr}(a)\defeq R(\Suc 0,0;u.a),\qquad c\defeq\operatorname{wr}(0),\qquad
 H\defeq(c=0).
\]
An ordinary assumption $F$ is the constant clause $[H;\,v.\,F]$ with $v$ unused; this applies
also to the local assumption of rule (I). The rules are:
\begin{enumerate}
\item Computation, in every context:
\[
 R(\Suc n,a;u.h)=h\langle R(n,a;u.h)/u\rangle . \tag{R}
\]
\item One-term crossed transport, for a term template $h(v)$ and an arbitrary source:
\[
 \infrule{\Gamma\proves s=t\qquad \Gamma\proves h(s)=r}{\Gamma\proves r=h(t)}. \tag{T}
\]
\item Clause use, for $[A;\,y.\,B(y)]\in\Gamma$:
\[
 \infrule{\Gamma\proves A}{\Gamma\proves B(t)}. \tag{U}
\]
\item Paired elimination, with $b$ fresh for $\Gamma$, $F$ and the parameters of $B$:
\[
 \infrule{\Gamma\proves\D{x}{y}B(x,y)\qquad \Gamma,\ [B(b,0);\,y.\,B(b,y)]\proves F}
         {\Gamma\proves F}. \tag{E}
\]
\item Witness-directed introduction, with $z$ fresh for $\Gamma$, $t$ and the parameters of
$B$:
\[
 \infrule{\Gamma,\ B(t,\eta)\proves B(t,\Suc z)}{\Gamma\proves\D{x}{y}B(x,y)}
 \qquad
 \eta=\begin{cases} z, & B \text{ is a literal equation and } t\equiv 0,\\
                    0, & \text{otherwise.}\end{cases} \tag{I}
\]
\end{enumerate}
\end{definition}

There is no separate hypothesis rule, zero-computation rule, closure constructor, closure
elimination, reflexivity axiom, case rule or induction rule. The neutral antecedent is
$H$ and not $0=0$, because $H$ is an instance of (R) and reflexivity is not yet available.
The clause pair, its bound variable and its substitution are part of the specification and
of the implementation; they are not a saving that goes uncounted.

Computation proves $\operatorname{wr}(a)=a$ in every context, in particular $H$. Use of an
ordinary assumption is therefore computation of $H$ followed by (U). Weakening, renaming of
fresh parameters and cut of an ordinary assumption follow by induction on derivations; in a
cut each use of the constant clause is replaced by the supplied proof of its conclusion. No
substitution theorem for rule (I) is presupposed.

\begin{lemma}[Equality from one term template]\label{ctr:lem:fiveeq}
$\mathcal K_5$ derives reflexivity, symmetry, transitivity, congruence and transport in
every literal-equation template.
\end{lemma}
\begin{proof}
Use $\operatorname{wr}(a)=a$ as both premises of (T) with the identity template to get
$a=a$. To reverse $s=t$, take the major premise $H$ and the constant template $h(v)=s$; (T)
returns $t=s$. Given $s=t$ and $t=r$, reverse the former and apply (T) with the identity
template to the major $t=r$ and the minor $t=s$, which gives $s=r$. Congruence is (T) with
the major $s=t$ and the reflexive minor $h(s)=h(s)$. For a template $a(v)=b(v)$ the two
applications
\[
 s=t,\ a(s)=b(s)\ \Longrightarrow\ b(s)=a(t),\qquad
 s=t,\ b(s)=a(t)\ \Longrightarrow\ a(t)=b(t)
\]
give transport. Each template is fixed finite syntax; no function variable occurs.
\end{proof}

\begin{lemma}[All case introductions]\label{ctr:lem:fivecase}
For every formula $B$ and term $t$ the rule ``from $\Gamma,B(t,0)\proves B(t,\Suc z)$ infer
$\D{x}{y}B(x,y)$'' is admissible in $\mathcal K_5$.
\end{lemma}
\begin{proof}
This is the repair of \cref{claus:lem:cases}, at the second name $c=\operatorname{wr}(0)$
of zero. The only instance that is not primitive has an equation matrix $E$ and the witness
$0$. Under $E(c,0)$, the theorem $c=0$ and \cref{ctr:lem:fiveeq} give $E(0,0)$; cut with the
supplied branch gives $E(0,\Suc z)$; transport back gives $E(c,\Suc z)$. Rule (I) at the
witness $c$ is in case mode and proves the original goal.
\end{proof}

\begin{lemma}[Transport in every formula]\label{ctr:lem:fivetransport}
For every template $F(v)$, from $s=t$ and $F(s)$ one derives $F(t)$ in $\mathcal K_5$, under
arbitrary retained clauses.
\end{lemma}
\begin{proof}
The proof of \cref{claus:lem:ftransport}, with rules (E) and (U) for \rulename{Elim} and
\rulename{Use}, and \cref{ctr:lem:fiveeq,ctr:lem:fivecase} for
\cref{claus:lem:eq,claus:lem:cases}.
\end{proof}

\subsection{Closure as an abbreviation}

For a fresh unused variable $w$ write $q_yB(y)\defeq\D{w}{y}B(y)$.

\begin{lemma}\label{ctr:lem:fiveclosure}
In $\mathcal K_5$:
\begin{enumerate}
\item from $\Gamma,B(0)\proves B(\Suc z)$ one derives $\Gamma\proves q_yB$;
\item from $q_yB$ and $B(0)$ one derives $B(t)$;
\item in every context that contains $[B(0);\,y.\,B(y)]$ one derives $q_yB$;
\item for a literal equation $E$, from $\Gamma,E(z)\proves E(\Suc z)$ one derives
$\Gamma\proves q_yE$; hence induction on literal equations with an arbitrary context,
numerical cases for all formulas, and introduction by progress at every witness are
derivable.
\end{enumerate}
\end{lemma}
\begin{proof}
Items (1)--(3) and the first part of (4) are \cref{claus:lem:closure}, and the derived rules
of (4) are \cref{claus:cor:induction}; their proofs use only (E), (U), (I) and
\cref{ctr:lem:fivecase}, which here play the parts of \rulename{Elim}, \rulename{Use},
\rulename{Intro} and \cref{claus:lem:cases}.
\end{proof}

Items (1)--(3) are introduction, elimination and eta expansion of a closure resource. The
third is what translates a branch of (DE) in $\mathcal K_6$, which assumes a closure formula.

\subsection{Comparison with the six rules}

Define $T$ from the formulas of $\mathcal K_6$ to those of $\mathcal K_5$ by
\[
 T(s=t)=(s=t),\qquad T(\D{x}{y}B)=\D{x}{y}T(B),\qquad T(\Q{y}B)=q_yT(B),
\]
with every inserted coordinate fresh and unused. $T$ preserves free variables, the
distinction between equation and compound roots, and normalized substitution. Let $U$ be the
inclusion of the formulas of $\mathcal K_5$ into the native formulas, and let
$\operatorname{ord}(\Delta)$ replace each formula of a list $\Delta$ by its constant clause.

\begin{theorem}[Preservation in both directions]\label{ctr:thm:fivesix}
For every native context $\Gamma$ and formula $A$,
\[
 \Gamma\proves_{\mathcal K_6}A\ \Longrightarrow\
 \operatorname{ord}(T\Gamma)\proves_{\mathcal K_5}T(A),
\]
and for formulas $\Delta,F$ of $\mathcal K_5$,
\[
 \operatorname{ord}(\Delta)\proves_{\mathcal K_5}F\ \Longrightarrow\
 \Delta\proves_{\mathcal K_6}F .
\]
\end{theorem}
\begin{proof}
\emph{Forward.} Induct on the derivation in $\mathcal K_6$. \rulename{Hyp} uses $H$ and the
constant clause. Computation is unchanged. (X) is derived transport followed by symmetry.
(QX) is \cref{ctr:lem:fiveclosure}(2), followed by symmetry when the conclusion is an
equation. (WI) is (I): the translation preserves roots and the witness is unchanged. For
(DE), the translated major premise proves $\D{x}{y}T(B)$; rule (E) opens
$[T(B)(b,0);\,y.\,T(B)(b,y)]$; \cref{ctr:lem:fiveclosure}(3) derives the translated
assumption $q_yT(B)(b,y)$ of the old branch, and cut inserts this proof into the translated
branch. All new variables are freshened first.

\emph{Backward.} Track each context occurrence by the rule that introduced it. A constant
clause is represented by its formula as an assumption of $\mathcal K_6$, and a clause opened
by (E) by the assumption $\Q{y}B(b,y)$. (U) on a constant clause is \rulename{Hyp}, and (U)
on an opened clause is ordinary \rulename{QE}, which $\mathcal K_6$ derives
(\cref{ctr:lem:sixeq}). (R) is the same computation, (T) is derived one-term transport, (E)
is (DE), and (I) is (WI). Induction on the derivation gives the claim.
\end{proof}

The two directions are those of \cref{claus:thm:five}(ii) and (i), with $\mathcal K_6$ in
the place of $\Nat$. A clause supplied from outside, which is neither constant nor opened by
(E), is treated as in \cref{claus:rem:external}; a checker for closed certificates needs no
such clause.

\begin{theorem}[Both composites]\label{ctr:thm:fivecomp}
$T(U(F))=F$ literally for every formula $F$ of $\mathcal K_5$, and
$A\dashv\vdash_{\mathcal K_6}U(T(A))$ for every native formula $A$. Both numerical
translations are the identity. Consequently $\mathcal K_5$ is equiconsistent and
bi-interpretable with $\iPRA$, with whole domains, actual equality and numerical composites
of graph $x=y$; the designated contradiction $0=\Suc 0$ is unchanged.
\end{theorem}
\begin{proof}
The first identity is immediate from the clauses of $T$. For the second use
$\Q{y}B(y)\dashv\vdash\D{w}{y}B(y)$ in $\mathcal K_6$ with $w$ unused: forward, under $B(0)$
specialize the closure at $\Suc z$ and introduce the paired formula by cases; backward, (DE)
opens exactly $\Q{y}B$ at an unused witness and \rulename{Hyp} closes the branch. The
binder congruence of \cref{ctr:lem:normaleq} lifts this through $\Q{y}$ and $\D{x}{y}$, by
induction on the formula. No quotient, retract or reinterpretation of equality is
introduced. The consequence follows from \cref{ctr:thm:fivesix,ctr:thm:sixbridge}.
\end{proof}

Under the reading $[A;\,y.\,B]^\circ=A^\circ\to\forall y\,B^\circ(y)$ the five rules are
directly sound in $\mathsf T_R^1$: (R) is a defining equation, (T) is transport followed by
symmetry, (U) is implication and universal elimination, (E) is existential elimination, and
(I) is coverage in case mode and induction on an equation under the assumed zero instance in
progress mode, followed by existential introduction. Certificates of $\mathcal K_5$ can be
longer than those of $\mathcal K_6$ while the checker is shorter: an ordinary assumption
needs the neutral computation premise, and the removal of the closure constructor adds
elimination and eta-expansion steps.

\section{Counter evaluation}\label{ctr:sec:ce}

This section draws on \pkg{counter\_evaluation} and \pkg{uniform\_binding}.

On the counter-aware recursor, the computation rule and the transport rule merge into one
rule that places the supplied equation in the counter. The rule has no analogue on iterator
terms, since it needs a body that sees the counter.

\subsection{The clause syntax and a liberal calculus}

Throughout this section and the next, terms are canonical terms of the counter-aware
language, either zero-body canonical (with $\nu$) or normal (with $\nf$), and every
displayed substitution is normalized accordingly. Formulas are
$A ::= (s=t)\mid\D{x}{y}B$ with the reading of \cref{base:sec:reading}. A context is a
finite reusable list of clauses $[A;\,y.\,B(y)]$, $y\notin\FV(A)$, read as
$A^\circ\to\forall y\,B^\circ(y)$ as in \cref{base:sec:clauses}; a clause may also have an
empty antecedent, $[\,;\,y.\,B(y)]$, read as $\forall y\,B^\circ(y)$. An ordinary assumption
$A$ is the constant clause $[0=0;\,v.\,A]$, or $[\,;\,v.\,A]$ where empty antecedents are
used, with $v$ unused. The designated contradiction is $0=1$.

The calculi below have rules with tests on literal syntax. The following calculus has none
and serves as the common reference.

\begin{definition}[The liberal clause calculus $\mathcal L$]\label{ctr:def:liberal}
The rules of $\mathcal L$ are:
\begin{enumerate}
\item \rulename{Refl}: $\Gamma\proves t=t$.
\item \rulename{Comp}: $\Gamma\proves R(\Suc n,a;i,u.h)=h\langle n/i,R(n,a;i,u.h)/u\rangle$.
\item \rulename{Tr}: from $s=t$ and $h\langle s/v\rangle=r$ infer $r=h\langle t/v\rangle$, for a
term template $h$.
\item \rulename{Use}: for $[A;\,y.\,B(y)]\in\Gamma$, from $\Gamma\proves A$ infer
$\Gamma\proves B\langle t/y\rangle$; for a clause with empty antecedent there is no premise.
\item \rulename{DE}: from $\D{x}{y}B(x,y)$ and a derivation of $A$ from
$\Gamma,[B(b,0);\,y.\,B(b,y)]$, with $b$ fresh for $\Gamma$, $A$ and the parameters of $B$,
infer $A$.
\item \rulename{Case}: from $\Gamma,B(t,0)\proves B(t,\Suc z)$ infer $\D{x}{y}B(x,y)$, for
every matrix and witness.
\item \rulename{Prog}: from $\Gamma,E(t,z)\proves E(t,\Suc z)$ infer $\D{x}{y}E(x,y)$, for a
literal equation $E$.
\end{enumerate}
In the last two rules $z$ is fresh for $\Gamma$, $t$ and the parameters of the matrix. On
normal terms both sides of \rulename{Comp} have the same spelling, and the rule is an
instance of \rulename{Refl}.
\end{definition}

\begin{lemma}[Structure of $\mathcal L$]\label{ctr:lem:libstruct}
The following are admissible in $\mathcal L$: weakening; injective renaming of parameters;
numerical substitution; cut of an ordinary assumption; and cut of a clause, that is, from
$\Gamma,[A;\,y.\,B(y)]\proves F$ and $\Gamma,A\proves B(y)$ with $y$ fresh for $\Gamma,A$
infer $\Gamma\proves F$.
\end{lemma}
\begin{proof}
Induct on derivations after freshening eigenparameters. Every rule is a schema without a
test on the shape of a term, so \cref{ctr:lem:coherence}, respectively \cref{ctr:lem:nf},
shows that a substitution instance of a rule instance is a rule instance. For the cut of an
assumption replace each use of its constant clause by the supplied proof. For the cut of a
clause replace each use of the clause at a term $t$, whose premise proves $A$, by the second
derivation with $t$ substituted for $y$, followed by a cut of the assumption $A$.
\end{proof}

\begin{lemma}[Derived rules of $\mathcal L$]\label{ctr:lem:libderived}
$\mathcal L$ derives:
\begin{enumerate}
\item symmetry, transitivity, congruence, and transport in every equation template;
\item numerical cases: from $\Gamma\proves F(0)$ and $\Gamma\proves F(\Suc z)$, $z$ fresh,
infer $\Gamma\proves F(t)$;
\item equation induction: from $\Gamma\proves E(0)$ and $\Gamma,E(z)\proves E(\Suc z)$, for a
literal equation $E$ and $z$ fresh, infer $\Gamma\proves E(t)$;
\item transport in every formula template: from $s=t$ and $F(s)$ infer $F(t)$;
\item binder congruence: if $B(x,y)$ and $B'(x,y)$ are interderivable in every context
uniformly in $x,y$, then so are $\D{x}{y}B$ and $\D{x}{y}B'$.
\end{enumerate}
\end{lemma}
\begin{proof}
(1) As in \cref{ctr:lem:fiveeq}, with \rulename{Refl} in place of the two uses of
$\operatorname{wr}$: to reverse $s=t$ apply \rulename{Tr} to the major $0=0$ and the
constant template.

(2), (3) are \cref{claus:cor:induction}, with \rulename{Case} for \cref{claus:lem:cases}
and \rulename{Prog} at the witness $0$ for the progress closure: for an unused coordinate $w$
and $q_yA\defeq\D{w}{y}A(y)$, \rulename{Case} introduces $q_yF$ from
$\Gamma\proves F(\Suc z)$, ignoring the assumed base, \rulename{Prog} introduces $q_yE$
from the step, and \rulename{DE} and \rulename{Use} at $t$ with the given base finish.

(4) is \cref{claus:lem:ftransport}, with (1) for the equations and \rulename{Case} at the
fresh witness.

(5) The construction of (4) at the outer binder, with the two conversions given by the
hypothesis.
\end{proof}

\subsection{The five patterns of counter evaluation}

\begin{definition}[The calculus $\mathcal K_{\mathrm{CE}}$]\label{ctr:def:ce}
Terms are zero-body canonical. Write $\varrho(t,a,h)$ for the canonical term
$R(t,a;i,u.h)$, which is $\Suc a$ when $t\equiv 0$. Every clause has an antecedent, and an
ordinary assumption $A$ is $[0=0;\,v.\,A]$. The rules are \rulename{Refl}, \rulename{Use}
and \rulename{DE} of \cref{ctr:def:liberal}, the rule of counter evaluation
\begin{equation}\label{ctr:eq:CE}
 \infrule{\Gamma\proves s=t\qquad \Gamma\proves R(\Suc s,a;i,u.h)=r}
         {\Gamma\proves r=h\langle t/i,\varrho(t,a,h)/u\rangle}\tag{CE}
\end{equation}
and the witness-directed introduction, with $z$ fresh for $\Gamma$, $w$ and the parameters
of $B$,
\begin{equation}\label{ctr:eq:CEDI}
 \infrule{\Gamma,\ B(w,\eta)\proves B(w,\Suc z)}{\Gamma\proves\D{x}{y}B(x,y)}
 \qquad
 \eta=\begin{cases} z, & B \text{ is a literal equation and } w\equiv 0,\\
                    0, & \text{otherwise.}\end{cases}\tag{DI}
\end{equation}
\end{definition}

In \eqref{ctr:eq:CE} the two substitutions in $h$ are simultaneous. The seed, the body and
the passive parameters are read from the second premise and are kept fixed: occurrences of
$s$ inside $a$ or $h$ are not replaced. The first endpoint $s$ is an arbitrary term,
including the literal $0$. There is no further computation axiom, equality rule, induction
rule, successor axiom or logical connective. The count of five patterns counts
\eqref{ctr:eq:CEDI} as one schema with two modes; it is not a claim under every convention
of counting rules. \pkg{uniform\_binding} calls \rulename{DE} and \eqref{ctr:eq:CEDI}
\rulename{Open} and \rulename{Close}.

\begin{theorem}[Arithmetic reconstruction]\label{ctr:thm:arith}
\rulename{Refl} and \eqref{ctr:eq:CE} alone derive the successor computation
\rulename{Comp}, symmetry, transitivity, congruence in every term template, and the crossed
transport \rulename{Tr}. The derivations use no logical rule, no numerical cases and no
induction.
\end{theorem}
\begin{proof}
\emph{Computation.} Apply \eqref{ctr:eq:CE} to the reflexive premises $n=n$ and $B=B$ with
$B=R(\Suc n,a;i,u.h)$. The conclusion is $B=h\langle n/i,\varrho(n,a,h)/u\rangle$.

\emph{Root composition.} Let $p$ prove $s=t$ and $q$ prove $s=r$, and put
$K_s\defeq R(\Suc s,0;i,u.i)$. \eqref{ctr:eq:CE} applied to $q$ and $K_s=K_s$ gives
$K_s=r$, and applied to $p$ and this result it gives $r=t$. Write
$\mathsf J(p,q)$ for this derivation of $r=t$ from $s=t$ and $s=r$.

\emph{Symmetry and transitivity.} $\mathsf J(\rulename{Refl}(s),p)$ proves $t=s$ from
$p\colon s=t$. Given $p\colon s=t$ and $q\colon t=v$, the derivation
$\mathsf J(q,\mathrm{Sym}(p))$ proves $s=v$.

\emph{Congruence.} For a template $F(v)$ choose the binders fresh and put
$B_F\defeq R(\Suc s,0;i,u.F(i))$. \eqref{ctr:eq:CE} on two reflexive premises gives
$B_F=F(s)$, and with first premise $p\colon s=t$ and this second premise it gives
$F(s)=F(t)$. The template may have repeated occurrences and occurrences under the binders of
nested recursors.

\emph{Crossed transport.} Apply $\mathsf J$ to $F(s)=F(t)$ and the supplied $q\colon F(s)=r$;
the result is $r=F(t)$. Every substitution in the argument is an instance of
\eqref{ctr:eq:CE}; no congruence is presumed in order to establish congruence.
\end{proof}

\begin{proposition}[Proof costs]\label{ctr:prop:cost}
Count every primitive inference, including a \rulename{Refl} leaf, as one node and count
repeated subtrees repeatedly. Root composition adds three nodes to its premises, congruence
adds four, and crossed transport adds seven. Each supplied premise occurs once.
\end{proposition}
\begin{proof}
$\mathsf J(p,q)$ is two applications of \eqref{ctr:eq:CE} and one leaf. Congruence is two
applications and two leaves. Crossed transport is $\mathsf J$ applied to a congruence,
$4+3=7$. The term labels contain copies of $F$ and of the endpoints, so these are bounds on
inference nodes, not on bytes or running time.
\end{proof}

\begin{theorem}[Equality-interface equivalence]\label{ctr:thm:interface}
Let $\mathcal K_{\mathrm{CT}}$ be $\mathcal K_{\mathrm{CE}}$ with \rulename{Refl} and
\eqref{ctr:eq:CE} replaced by \rulename{Comp} and \rulename{Tr}. A sequent is derivable in
$\mathcal K_{\mathrm{CE}}$ if and only if it is derivable in $\mathcal K_{\mathrm{CT}}$, and
the two proof translations are effective and preserve end sequents literally.
\end{theorem}
\begin{proof}
\Cref{ctr:thm:arith} replaces each \rulename{Comp} or \rulename{Tr} node. Conversely, in
$\mathcal K_{\mathrm{CT}}$ the computation with the counter projection gives $K_s=s$, and
\rulename{Tr} applied to two copies of this equation with the identity template gives
$s=s$. \rulename{Tr} with a constant template gives symmetry, with the identity template
and a reversed premise it gives transitivity, and with a reflexive minor premise it gives
congruence. For \eqref{ctr:eq:CE}: from $s=t$, congruence in the counter gives
$R(\Suc s,a;i,u.h)=R(\Suc t,a;i,u.h)$, and \rulename{Comp} gives
$R(\Suc t,a;i,u.h)=h\langle t/i,\varrho(t,a,h)/u\rangle$; reverse the second premise
$R(\Suc s,a;i,u.h)=r$ and compose. The three logical rules are common. The use of ordinary
equality in the converse is legitimate, since it is derived in $\mathcal K_{\mathrm{CT}}$
before \eqref{ctr:eq:CE} is simulated; the forward direction never assumes it.
\end{proof}

\begin{lemma}[Both introductions]\label{ctr:lem:ceintro}
\rulename{Case} and \rulename{Prog} are derivable in $\mathcal K_{\mathrm{CE}}$ for every
matrix and witness. Before this, weakening, renaming of fresh parameters and cut of an
ordinary assumption are admissible in $\mathcal K_{\mathrm{CE}}$.
\end{lemma}
\begin{proof}
The three structural operations do not change the literal-zero test of
\eqref{ctr:eq:CEDI}; for cut, replace each \rulename{Use} of the constant clause
$[0=0;\,v.\,A]$ by the supplied proof of $A$ and discard the proof of its antecedent.

The only instance of \rulename{Case} that is not an instance of \eqref{ctr:eq:CEDI} has an
equation matrix $E$ and the witness $0$. It is repaired as in \cref{claus:lem:cases}, at
the second name $Z\defeq R(\Suc 0,0;i,u.0)$ of zero, which is not literally $0$ and for
which \cref{ctr:thm:arith} proves $Z=0$: under $E(Z,0)$, transport in the equation gives
$E(0,0)$; cut with the premise gives $E(0,\Suc z)$; transport back gives $E(Z,\Suc z)$; and
\eqref{ctr:eq:CEDI} at the witness $Z$ is in case mode. Only transport in an equation is
used, not transport through a binder.

\rulename{Prog} at a witness $w$ is the third statement of \cref{claus:cor:induction}, with
the same proof: \eqref{ctr:eq:CEDI} in progress mode at the witness $0$ of an unused
coordinate gives $\D{v}{y}E(w,y)$, \rulename{DE} opens $[E(w,0);\,y.\,E(w,y)]$, and under
$E(w,0)$ \rulename{Use} and \rulename{Case} at $w$ conclude.
\end{proof}

\begin{theorem}\label{ctr:thm:ceequiv}
On zero-body canonical terms, $\mathcal K_{\mathrm{CE}}$, $\mathcal K_{\mathrm{CT}}$ and
$\mathcal L$ derive the same sequents whose clauses have antecedents. In particular
numerical substitution is admissible in $\mathcal K_{\mathrm{CE}}$.
\end{theorem}
\begin{proof}
Every rule of $\mathcal L$ is derivable in $\mathcal K_{\mathrm{CE}}$ by
\cref{ctr:thm:arith,ctr:lem:ceintro}. Conversely \eqref{ctr:eq:CE} is derivable in
$\mathcal L$ as in \cref{ctr:thm:interface}, and an instance of \eqref{ctr:eq:CEDI} is an
instance of \rulename{Case} or of \rulename{Prog}. Substitution is
\cref{ctr:lem:libstruct}.
\end{proof}

The test of \eqref{ctr:eq:CEDI} is not invariant under substitution: a variable witness can
become $0$. A progress introduction keeps its equation root and its literal witness $0$,
but a case introduction whose witness becomes $0$ has to be replaced by the derived
\rulename{Case}. Admissibility of substitution is obtained through $\mathcal L$, not by a
rule-by-rule argument.

\subsection{Ordinary logic with the decoder}\label{ctr:sec:macros}

The derivation of intuitionistic logic from the paired binder is that of the gates,
selectors and connectives of \cref{base:sec:structural}; it needs a selector with the two
endpoint equations. With the
decoder $P$ of \cref{ctr:lem:decoder} the conjunction and the universal quantifier take a
shorter form: no parity iterator and no separate zero instance are needed. The construction
is given in full for $\mathcal L$, on either kind of canonical terms.

\begin{definition}[Gates, selection and the logical macros]\label{ctr:def:macros}
With binders renamed away from the control $v$, put
\[
 G_v(s=t)\defeq\bigl(C(v,0,s)=C(v,0,t)\bigr),\qquad G_v(\D{x}{y}B)\defeq\D{x}{y}G_vB,
\]
and $L(F)\defeq\D{w}{y}G_yF$ with $w,y$ fresh. Every formula is a finite chain of paired
binders ending in an equation. For formulas $A,B$ apply $L$ to the shorter one until the
depths agree, align the bound names, and at the equations put
\[
 M_v\bigl((s=t),(a=b)\bigr)\defeq\bigl(C(v,s,a)=C(v,t,b)\bigr),
\]
carrying the matched binders outside. With all auxiliary binders fresh and $w$ unused
define
\begin{equation}\label{ctr:eq:macros}
\begin{aligned}
 \bot_*&\defeq(0=1), &
 A\to_*B&\defeq\D{w}{y}M_y(A,B),\\
 A\land_*B&\defeq\D{w}{y}G_yM_{P(y)}(A,B), &
 \exists_*x\,A(x)&\defeq\D{x}{y}G_yA(x),\\
 A\lor_*B&\defeq\exists_*x\,M_x(A,B), &
 \forall_*x\,A(x)&\defeq\D{w}{y}G_yA(P(y)).
\end{aligned}
\end{equation}
These are finite formulas of the same syntax, not further constructors.
\end{definition}

\begin{lemma}[Gates and selection]\label{ctr:lem:gates}
In $\mathcal L$, for all formulas and in every context:
$\proves G_0F$; $G_{\Suc z}F\dashv\vdash F$; $L(F)\dashv\vdash F$;
$M_0(A,B)\dashv\vdash A$; and $M_{\Suc z}(A,B)\dashv\vdash B$.
\end{lemma}
\begin{proof}
For an equation the first two statements are the endpoint equations of $C$, reflexivity and
transport. For $\D{x}{y}B$ the first is \rulename{Case} with the gate theorem of the matrix,
the assumed base being unused, and the second is binder congruence
(\cref{ctr:lem:libderived}). For $L(F)$: forward, \rulename{Case} uses $F$ to prove the
successor gate; backward, \rulename{DE}, the theorem $G_0F$ and \rulename{Use} at $1$ give
$G_1F$ and hence $F$. For $M$: at the equations the endpoint equations of $C$ give the two
statements for the padded formulas, binder congruence carries them through the common
binders, and $L(F)\dashv\vdash F$ removes the padding.
\end{proof}

\begin{lemma}[Natural deduction]\label{ctr:lem:macros}
In $\mathcal L$ the macros \eqref{ctr:eq:macros} admit every introduction and elimination
rule of intuitionistic first-order natural deduction, with the usual discharge and
eigenvariable conditions, and ex falso for $\bot_*$. Moreover $\mathcal L$ proves coverage
$n=0\lor_*\exists_*m\,(n=\Suc m)$.
\end{lemma}
\begin{proof}
\emph{Implication.} Given a derivation of $B$ under the assumption $A$, \rulename{Case}
introduces $A\to_*B$: the assumed zero instance gives $A$ and the successor instance is
equivalent to $B$. For elimination, \rulename{DE} retains the selector clause, the supplied
$A$ proves its antecedent, and \rulename{Use} at $1$ gives $B$.

\emph{Conjunction.} From $A$ and $B$, numerical cases prove $M_v(A,B)$ for every term $v$,
in particular for $v=P(\Suc z)$. The successor gate and \rulename{Case} give $A\land_*B$;
the antecedent of its clause is a gate theorem. For elimination, \rulename{DE} and
\rulename{Use} at $1$ and at $\Suc 1$ give the selector at $P(1)=0$ and at $P(\Suc 1)=1$,
hence $A$ and $B$ by transport. No parity or range theorem is used.

\emph{Existential quantifier.} From $A(t)$, \rulename{Case} at the witness $t$ introduces
the gated formula. Conversely \rulename{DE} at a fresh $b$, the zero gate and
\rulename{Use} at $1$ give $A(b)$; cut this into a branch whose conclusion and context do
not contain $b$, and close \rulename{DE}.

\emph{Disjunction.} The witnesses $0$ and $1$ give the two injections. For elimination,
discharge the two branches as $A\to_*F$ and $B\to_*F$. Numerical cases, in the original
context, prove $M_v(A,B)\to_*F$ for every $v$. Existential elimination supplies
$M_b(A,B)$ at a fresh tag $b$, and implication elimination gives $F$. The tag is not
substituted into a retained assumption, since the numerical cases are performed before that
assumption is added.

\emph{Universal quantifier.} A derivation of $A(z)$ with $z$ fresh gives
$G_{\Suc z}A(P(\Suc z))$ by $P(\Suc z)=z$ and the gate equivalence, and \rulename{Case} with
an unused witness gives $\forall_*x\,A(x)$. For elimination, \rulename{DE}, the zero gate
and \rulename{Use} at $\Suc t$ give $G_{\Suc t}A(P(\Suc t))$ and hence $A(t)$. Every
numerical target is covered directly; the value $P(0)=1$ plays no role, because the outer
gate makes the zero instance provable.

\emph{Falsity.} Under $0=1$, transport in the control of $C(v,s,t)$ and the two endpoint
equations give $s=t$ for arbitrary terms. Induction on the number of binders, with
\rulename{Case}, gives every formula.

\emph{Coverage.} Numerical cases on $n$: the left injection at $0$, and the right injection
with the witness $z$ at $\Suc z$.
\end{proof}

\begin{remark}[Variants of the macros]\label{ctr:rem:macrovariants}
The form \eqref{ctr:eq:macros} is that of \pkg{uniform\_binding} and
\pkg{single\_constructor}. \pkg{counter\_evaluation} uses a term $\pi$ with $\pi(1)=1$ and
$\pi(\Suc 1)=0$ in the conjunction and defines
$\forall_*x\,A(x)\defeq A(0)\land_*\D{w}{x}A(x)$. \pkg{uniform\_recursion} uses
$K(y)=C(P(y),1,0)$ in the conjunction and the form above for the universal quantifier.
\pkg{clausal\_canonical}, which has no decoder, uses
$p(z)=\widehat I(z,0;u.C(u,1,0))$ with $p(1)=1$, $p(\Suc 1)=0$ and the universal quantifier
with a separate zero instance. \pkg{single\_constructor} defines $M_v$ by recursion on the
pair of formulas, padding an equation $E$ against a paired formula by $G_yE$ under the
aligned binder; the result is interderivable with the definition above.
\end{remark}

\subsection{Adequacy and the main theorem}

\begin{theorem}[Adequacy of $\mathcal L$]\label{ctr:thm:adequacy}
Let $A\mapsto A^*$ translate the formulas of $\mathsf T_R^2$ to the clause syntax: an
equation is normalized, $\bot$ becomes $\bot_*$, and the connectives and quantifiers become
the macros \eqref{ctr:eq:macros}.
\begin{enumerate}
\item If $\Gamma\proves_{\mathcal L}A$ then $\mathsf T_R^2$ proves $A^\circ$ from the
readings of the clauses of $\Gamma$.
\item If $\mathsf T_R^2$ proves $B$ from assumptions $\Delta$, then
$\Delta^*\proves_{\mathcal L}B^*$, the formulas of $\Delta^*$ being ordinary assumptions.
\item $\mathsf T_R^2\proves(B^*)^\circ\leftrightarrow B$ for every formula $B$ of
$\mathsf T_R^2$, and $(A^\circ)^*\dashv\vdash_{\mathcal L}A$ for every formula $A$ of the
clause syntax, under arbitrary retained clauses.
\end{enumerate}
Hence $\mathcal L$ is faithful for $\mathsf T_R^2$ in both directions, and both numerical
maps of this comparison are the identity.
\end{theorem}
\begin{proof}
(1) Induct on the derivation. By \cref{ctr:lem:normsound} a normalized instance is provably
equivalent to the raw instance. \rulename{Refl}, \rulename{Comp} and \rulename{Tr} are
equality reasoning and a computation equation. \rulename{Use} is implication and universal
elimination, and \rulename{DE} is existential elimination. For \rulename{Case}, assume
$B(t,0)$; coverage and the premise give $B(t,y)$ for every $y$; then introduce the witness
$t$. For \rulename{Prog}, assume $E(t,0)$ and apply \eqref{ctr:eq:EI} with the premise as
step.

(2) The logical rules are \cref{ctr:lem:macros} and the equality rules are
\cref{ctr:lem:libderived}. The translation of the zero computation has the same canonical
term on both sides, and that of the successor computation is \rulename{Comp}. Coverage is
\cref{ctr:lem:macros}. Separation translates to $(0=1)\to_*(0=1)$. The rule
\eqref{ctr:eq:EI} is \cref{ctr:lem:libderived}(3), with the unchanged context.

(3) For the first statement induct on $B$: in $\mathsf T_R^2$ the gate and selector
endpoints, coverage for an arbitrary control and $P(\Suc t)=t$ show that the reading of each
macro is equivalent to the connective or quantifier it represents; for example
$\forall y\,(M_y(A,B))^\circ$ is equivalent to $A^\circ\land B^\circ$ by coverage, so the
reading of $A\to_*B$ is equivalent to $A^\circ\to B^\circ$. For the second induct on $A$.
Equations are unchanged. For $\D{x}{y}B$ apply the induction hypothesis to the matrix,
uniformly, and prove
\[
 \D{x}{y}B(x,y)\dashv\vdash
 \exists_*x\bigl(B(x,0)\to_*\forall_*y\,B(x,y)\bigr).
\]
Forward, open the left side at $b$; under $B(b,0)$ its clause gives $B(b,z)$ for a fresh
$z$; universal, implication and existential introduction build the right side; close
\rulename{DE}. Backward, existential elimination gives a fresh $b$ and the implication;
introduce $\D{x}{y}B$ by \rulename{Case} at $b$: the local base feeds the implication, and
universal elimination at $\Suc z$ gives the successor instance. Faithfulness follows: if
$\mathsf T_R^2$ proves $A^\circ$ from the readings of $\Gamma$, then by (2) and (3)
$\mathcal L$ derives $A$ from $\Gamma$, each clause $[A_0;\,y.\,B_0]$ providing the
translation of its reading by implication and universal introduction.
\end{proof}

\begin{theorem}[Counter evaluation]\label{ctr:thm:cemain}
The five-pattern calculus $\mathcal K_{\mathrm{CE}}$ is sound and faithful for
$\mathsf T_R^2$, and it is equiconsistent and bi-interpretable with $\iPRA$: both numerical
domains are the whole domain, equality is actual equality, the numerical composites have the
graph $x=y$, and every formula is provably equivalent to its double translation. The
translations preserve the full consequence relation, including open sequents and formulas
with unbounded quantifiers.
\end{theorem}
\begin{proof}
Combine \cref{ctr:thm:ceequiv,ctr:thm:adequacy,ctr:thm:transfer}. The contradiction
$0=1$ reads as an equation that $\mathsf T_R^2$ refutes by separation and that implies every
formula, so inconsistency transfers in both directions.
\end{proof}

\begin{remark}
\pkg{counter\_evaluation} and \pkg{uniform\_binding} present the same calculus.
\pkg{counter\_evaluation} contributes \cref{ctr:thm:arith,ctr:prop:cost,ctr:thm:interface}
and a checker that synthesizes the conclusion of each proof node; its ordinary theory has
the zero computation as an identity of canonical spelling. \pkg{uniform\_binding} takes the
ordinary theory over raw terms with both computation equations, reads canonical terms as
raw ones, justifies the normalization by \cref{ctr:lem:normsound}, and contributes the
macros \eqref{ctr:eq:macros} and the goal-directed checker with one traversal
(\cref{ctr:sec:binding}). The equivalence of the two certificate formats is
\cref{ctr:prop:checkers}.
\end{remark}

\section{Three families on normal forms}\label{ctr:sec:three}

This section draws on \pkg{uniform\_recursion} and \pkg{single\_constructor}.

On normal terms (\cref{ctr:def:normal}) both computation equations are identities of
spelling, and the equations of \cref{ctr:lem:decoder} are literal identities of normal
forms. A calculus then needs neither a computation rule nor, with a suitable elimination, a
transport rule. Two such calculi are stated on the counter-aware recursor. Both take the
clause syntax of \cref{ctr:sec:ce} with empty antecedents, and both contain in every context
the fixed clause
\begin{equation}\label{ctr:eq:reflclause}
 \mathcal A\defeq[\,;\,y.\,y=y],
\end{equation}
which is the initial context of a closed proof; it is the reflexivity clause $\mathcal R$ of
\cref{claus:def:v}, and the convention is that of \cref{disch:lem:conventions}. Reflexivity
is thus an axiom clause; it has not been deleted without replacement. An ordinary assumption
$F$ is $[\,;\,v.\,F]$. On iterator terms in full normal form, the two calculi are the
calculi $\mathcal W^{\mathcal R}$ and $\mathcal O^{\mathcal R}$ of \pkg{witness\_discharge}
and \pkg{inferred\_witness}, treated with their Lean formalizations in
\cref{disch:sec:name,disch:sec:inferred} (\cref{disch:thm:name,disch:thm:inferred}); the
counter-aware forms below are not formalized. The arguments of that chapter are stated for
calculi with the listed rules; apart from the packets, they use from the term syntax only
the literal test for $0$, the endpoint identities of the selector and the coherence of
normalization with substitution. On the normal terms of \cref{ctr:def:normal} these are
\cref{ctr:lem:decoder} and \eqref{ctr:eq:SC}, so such arguments are cited and not repeated. What the new syntax requires
is proved here: the packets that carry transport use the decoder, and the repair of the
introduction at the witness $0$ cannot use a closed term that is provably but not literally
$0$, because normal terms have none; the compound envelope of \cref{comp:lem:envelope} is
used instead.

\subsection{Use, Close and Name}

\begin{definition}[The calculus $\mathcal K_{\mathrm{Name}}$]\label{ctr:def:name}
Terms are normal and all displayed instances are normalized by $\nf$. The rules are:
\begin{enumerate}
\item \rulename{Use}: for $[A;\,y.\,B(y)]\in\Gamma$ and a term $t$, from $\Gamma\proves A$
infer $\Gamma\proves B(t)$; for an empty antecedent there is no premise. In particular
$\mathcal A$ gives $t=t$.
\item \rulename{Close}: the rule \eqref{ctr:eq:CEDI}, its local assumption being an
unconditional constant clause.
\item \rulename{Name}: for a fresh parameter $b$ and the clause
$\mathcal B_b\defeq[B(b,0);\,y.\,B(b,y)]$,
\begin{equation}\label{ctr:eq:name}
 \infrule{\Gamma\proves\D{x}{y}B\qquad \Gamma,\mathcal B_b\proves b=t\qquad
          \Gamma,\mathcal B_b\proves F}
         {\Gamma\proves F\langle t/b\rangle}
 \qquad b\notin\FV(\Gamma,F\langle t/b\rangle).
\end{equation}
The term $t$ and the formula $F$ may contain $b$. The substitution replaces the free
occurrences of $b$ once and does not revisit an inserted copy of $t$. The case $t\equiv b$
is permitted.
\end{enumerate}
\end{definition}

There is no transport rule, no hypothesis rule apart from \rulename{Use}, no computation
inference, no successor axiom and no separate induction rule. The rule \eqref{ctr:eq:name}
is the rule \rulename{Name} of \cref{disch:sec:name} on the present normal terms, with the
same soundness argument (\cref{disch:lem:name-sound}). Weakening, exchange, contraction,
injective renaming of parameters and cut of an ordinary assumption are admissible by
induction on derivations, as in \cref{disch:lem:structural}; these operations preserve the
test of \rulename{Close}.

\begin{lemma}[Elimination, closure and cases]\label{ctr:lem:nameelim}
In $\mathcal K_{\mathrm{Name}}$:
\begin{enumerate}
\item from $\Gamma\proves\D{x}{y}B$ and $\Gamma,\mathcal B_b\proves A$ with
$b\notin\FV(\Gamma,A)$ one infers $\Gamma\proves A$;
\item writing $q_yA(y)\defeq\D{x}{y}A(y)$ for an unused $x$: from
$\Gamma,A(0)\proves A(\Suc z)$ one infers $\Gamma\proves q_yA$, and from $q_yA$ and $A(0)$
one infers $A(t)$;
\item numerical cases: from $\Gamma\proves A(0)$ and $\Gamma\proves A(\Suc z)$, $z$ fresh,
one infers $\Gamma\proves A(t)$.
\end{enumerate}
\end{lemma}
\begin{proof}
(1) is the proof of \cref{disch:lem:name-elim}, with $\mathcal A$ for \rulename{Refl}: take
$t\equiv b$ in \eqref{ctr:eq:name}. (2) and (3) are \cref{disch:lem:closure}, whose proof
uses \rulename{Use}, the introduction at the unused witness $1$ (in case mode, since $1$ is
not literally $0$), weakening, and (1) for a matrix with an unused witness coordinate.
\end{proof}

\begin{lemma}[Transport from one discharge]\label{ctr:lem:nametransport}
For every term template $h(w)$, from $\Gamma\proves s=t$ and $\Gamma\proves h(s)=r$ one
infers $\Gamma\proves h(t)=r$ in $\mathcal K_{\mathrm{Name}}$. Hence symmetry,
transitivity, congruence and the crossed transport \rulename{Tr} are derivable.
\end{lemma}
\begin{proof}
Let $K(y)\defeq C(P(y),1,0)$, so that $K(1)\equiv 1$ and $K(\Suc 1)\equiv 0$ as normal
forms. Choose $x,y$ fresh and put
\[
 B(x,y)\defeq\Bigl(C\bigl(y,0,C(K(y),x,h(x))\bigr)=C\bigl(y,0,C(K(y),t,r)\bigr)\Bigr).
\]
Its normalized rows are $B(x,0)\equiv(0=0)$, $B(x,1)\equiv(h(x)=r)$ and
$B(x,\Suc 1)\equiv(x=t)$, and its row at $\Suc z$ is
$C(K(\Suc z),x,h(x))=C(K(\Suc z),t,r)$. To introduce $\D{x}{y}B$ at the witness $s$, apply
numerical cases to the equation $C(v,s,h(s))=C(v,t,r)$ and instantiate it at
$v=K(\Suc z)$: its zero case is $s=t$ and its successor case is $h(s)=r$. This derivation
ignores the local assumption of \rulename{Close} and is valid in either mode, including
$s\equiv 0$. Now apply \eqref{ctr:eq:name} with a fresh $b$. In the opened context the
antecedent $0=0$ holds by $\mathcal A$; \rulename{Use} at $\Suc 1$ gives the equality
premise $b=t$, and \rulename{Use} at $1$ gives the formula premise $h(b)=r$. The export is
$h(t)=r$ by \eqref{ctr:eq:SC}; the passive data do not contain $b$. This is the derivation
of \cref{disch:thm:name-transport} with the packet above in place of that of
\cref{disch:def:packet}. The remaining equality rules follow from this uncrossed transport
and $\mathcal A$ by \cref{disch:lem:equality}.
\end{proof}

The argument is not circular: ordinary elimination, closure and numerical cases
(\cref{ctr:lem:nameelim}) are available before any equality rule, and the packet needs only
the literal endpoint identities and no range theorem for $K$.

\begin{lemma}[Formula transport]\label{ctr:lem:nameformula}
From $s=t$ and $F(s)$ one derives $F(t)$ in $\mathcal K_{\mathrm{Name}}$, for every formula
template $F$.
\end{lemma}
\begin{proof}
\Cref{disch:lem:formula}, whose hypotheses hold by
\cref{ctr:lem:nameelim,ctr:lem:nametransport}: the atomic case is transport in an equation
template, and at a paired binder the introduction is at a fresh parameter, which is not
literally $0$, so that \rulename{Close} is in case mode.
\end{proof}

\begin{lemma}[The introduction at the witness $0$]\label{ctr:lem:envelope}
For an equation $E=(s=t)$ put $G_vE\defeq(C(v,0,s)=C(v,0,t))$ and
$L(E)\defeq q_vG_vE$. Then $G_0E\equiv(0=0)$ and $G_{\Suc w}E\equiv E$ literally, and
$E\dashv\vdash L(E)$ in every context. Consequently \rulename{Case} is derivable in
$\mathcal K_{\mathrm{Name}}$ for every matrix and witness, and so are \rulename{Prog} and
induction on literal equations.
\end{lemma}
\begin{proof}
The two literal identities are the endpoint equations of $C$ on normal forms
(\cref{ctr:lem:decoder}). With them, $G_vE$ is a gate in the sense of
\cref{comp:lem:envelope}, and the proof of that lemma gives $E\dashv\vdash L(E)$, with the
closure rules of \cref{ctr:lem:nameelim}(2) in place of the closure by cases and
\rulename{QE}, and $\mathcal A$ for the base $0=0$; no transport is used. The two
unrestricted introductions are then \cref{disch:lem:repair}, whose hypotheses hold in
$\mathcal K_{\mathrm{Name}}$ by \cref{ctr:lem:nameelim}(1): in the only instance of
\rulename{Case} that \rulename{Close} does not provide, the compound matrix $L(E(x,y))$ is
introduced at the witness $0$ in case mode, opened at a fresh $c$, and $\D{x}{y}E$ is
introduced at the witness $c$, which is not literally $0$. Induction on literal equations is
\rulename{Close} in progress mode at the witness $0$ of an unused coordinate, followed by the
elimination of \cref{ctr:lem:nameelim}(2) with the supplied base.
\end{proof}

\begin{theorem}\label{ctr:thm:nameequiv}
On normal terms, with $\mathcal A$ in every context, $\mathcal K_{\mathrm{Name}}$ and
$\mathcal L$ derive the same sequents. In particular numerical substitution is admissible in
$\mathcal K_{\mathrm{Name}}$.
\end{theorem}
\begin{proof}
\rulename{Refl} is \rulename{Use} of $\mathcal A$, \rulename{Comp} is an instance of it on
normal terms, and \rulename{Tr}, \rulename{DE}, \rulename{Case} and \rulename{Prog} are
\cref{ctr:lem:nametransport,ctr:lem:nameelim,ctr:lem:envelope}. Conversely \rulename{Close}
is an instance of \rulename{Case} or \rulename{Prog}. For \eqref{ctr:eq:name}, apply
\rulename{DE}; in the opened context derive $F$ and $b=t$; transport in the template
obtained from $F$ by replacing the free occurrences of $b$ by a new variable
(\cref{ctr:lem:libderived}) gives $F\langle t/b\rangle$, which does not contain $b$; close
\rulename{DE}. This is the second half of the proof of \cref{disch:thm:name}.
\end{proof}

\subsection{Use, Introduce and Open with an inferred witness}

\begin{definition}[The calculus $\mathcal K_{\mathrm{Open}}$]\label{ctr:def:open}
Terms are normal and all displayed instances are normalized. The rules are \rulename{Use}
and the introduction \eqref{ctr:eq:CEDI}, as in \cref{ctr:def:name}, and the following
elimination. For a major premise $\D{x}{y}B(x,y)$ and a parameter $b$ fresh for $\Gamma$
and the parameters of $B$, let $B_b(y)$ be the normalized instance $B(b,y)$ and
\begin{equation}\label{ctr:eq:probe}
 \tau_B(b)\defeq
 \begin{cases}
  r, & B_b(0)\equiv(0=0) \text{ and } B_b(1)\equiv(b=r),\\
  b, & \text{otherwise.}
 \end{cases}
\end{equation}
This is the selector $\rho_B$ of \cref{disch:def:selector} with the neutral antecedent
$0=0$: the variable $b$ has to be the entire left side of the row at $1$, and no search for
a provable equality takes place. The rule is
\begin{equation}\label{ctr:eq:open}
 \infrule{\Gamma\proves\D{x}{y}B\qquad
          \Gamma,\ [B(\tau_B(b),0);\,y.\,B(\tau_B(b),y)]\proves A}
         {\Gamma\proves A}
 \qquad b\notin\FV(A).\tag{Open}
\end{equation}
The branch is checked in the scope that contains $b$, also when the inferred term contains
$b$. The replacement of $b$ by $\tau_B(b)$ is simultaneous and occurs once.
\end{definition}

The rule \eqref{ctr:eq:open} is the rule \rulename{Open} of \cref{disch:sec:inferred} on
the present normal terms. It is sound for the ordinary reading by the proof of
\cref{disch:prop:inferred-sound}, and weakening, injective renaming of parameters and cut of
an ordinary assumption are admissible by \cref{disch:lem:structural}; injective renaming
preserves the literal tests.

\begin{lemma}[Safe openings and the repair of cases]\label{ctr:lem:safe}
In $\mathcal K_{\mathrm{Open}}$:
\begin{enumerate}
\item if the witness coordinate is unused in $B$, or the root of $B$ is compound, then
$\tau_B(b)\equiv b$ and \eqref{ctr:eq:open} is the ordinary elimination \rulename{DE};
\item the closure rules and numerical cases of \cref{ctr:lem:nameelim}(2),(3) hold;
\item \rulename{Case} is derivable for every matrix and witness.
\end{enumerate}
None of these uses a substitution theorem or an equality rule.
\end{lemma}
\begin{proof}
(1) is \cref{disch:lem:selector}(b): if $x$ is unused, the row at $1$ does not contain the
fresh $b$; if the root of $B$ is compound, the row at $0$ is compound and is not $0=0$.
(2) and (3) follow as in \cref{disch:lem:fallback}: \cref{disch:lem:closure} and
\cref{disch:lem:repair}, in the form of \cref{ctr:lem:envelope}, require the ordinary
elimination only for the two kinds of matrices of (1).
\end{proof}

\begin{lemma}[Substitution]\label{ctr:lem:opensubst}
If $\Gamma\proves A$ in $\mathcal K_{\mathrm{Open}}$, then
$\Gamma\langle\sigma\rangle\proves A\langle\sigma\rangle$ for every simultaneous numerical
substitution $\sigma$.
\end{lemma}
\begin{proof}
The proof of \cref{disch:thm:subst} applies without change, with \eqref{ctr:eq:SC} in place
of the coherence of normalization with substitution on iterator terms, \cref{ctr:lem:safe}(3)
for a case introduction whose witness becomes $0$ under an equation matrix, and the
preservation of a successful probe under a substitution that fixes $b$
(\cref{disch:lem:selector}(c), whose proof uses only that coherence). As there, a failed
probe is not stable under substitution, and in that case the induction hypothesis is applied
to the original branch with the substitution that extends $\sigma$ by $b\mapsto\tau_{B'}(b)$,
$B'$ the normalized instance of $B$ under $\sigma$; the induction is on proper
subderivations of the original derivation only.
\end{proof}

\begin{lemma}[Ordinary elimination, transport and induction]\label{ctr:lem:openrules}
In $\mathcal K_{\mathrm{Open}}$:
\begin{enumerate}
\item \rulename{DE} is admissible;
\item for every term template $h(w)$, from $s=t$ and $h(s)=r$ one infers $h(t)=r$; hence
symmetry, transitivity, congruence and transport in every formula template are derivable;
\item \rulename{Prog} and induction on literal equations are derivable.
\end{enumerate}
\end{lemma}
\begin{proof}
(1) As in \cref{disch:cor:open-elim}: given a branch under the clause at $b$, apply
\cref{ctr:lem:opensubst} to it with $b\mapsto\tau_B(b)$ and then \eqref{ctr:eq:open}; the
context and the conclusion do not contain $b$.

(2) Use the matrix
\begin{equation}\label{ctr:eq:packet}
 B(x,y)\defeq\Bigl(C\bigl(y,0,C(P(y),x,h(x))\bigr)=C\bigl(y,0,C(P(y),t,r)\bigr)\Bigr).
\end{equation}
Its normalized rows are $B(x,0)\equiv(0=0)$, $B(x,1)\equiv(x=t)$ and
$B(x,\Suc 1)\equiv(h(x)=r)$, and its row at $\Suc z$ is literally
$C(z,x,h(x))=C(z,t,r)$, because $P(\Suc z)$ computes to $z$. At the witness $s$, numerical
cases on $z$ prove this row: the zero case is $s=t$ and the successor case is $h(s)=r$. The
derivation ignores the local assumption of the introduction, so it is valid in either mode
and proves $\D{x}{y}B$. Now \eqref{ctr:eq:open} infers the witness $t$: the row at $0$ is
$0=0$ and the row at $1$ is $b=t$. The opened clause has the base $0=0$, which
$\mathcal A$ proves, and \rulename{Use} at $\Suc 1$ gives $h(t)=r$, which is in the outer
scope. This is the derivation of \cref{disch:thm:open-transport} with the packet
\eqref{ctr:eq:packet}. The remaining statements follow by
\cref{disch:lem:equality,disch:lem:formula}, with (1) for the eliminations.

(3) is \cref{disch:lem:repair} (progress introduction) and \cref{claus:cor:induction}:
the introduction at the witness $0$ with an unused witness coordinate gives $q_yE(y)$ from
$\Gamma,E(z)\proves E(\Suc z)$, its opening is safe by \cref{ctr:lem:safe}(1), and
\rulename{Use} with the base gives $E(t)$; \rulename{Prog} at a witness $t$ follows with
\rulename{Case}.
\end{proof}

The decoder makes the packet \eqref{ctr:eq:packet} simpler than the packet of
\cref{ctr:lem:nametransport} and than the parity packet of \cref{disch:def:packet} on
iterator terms: one successor computation exposes the variable on which the numerical cases
are taken.

\begin{theorem}\label{ctr:thm:openequiv}
On normal terms, with $\mathcal A$ in every context, $\mathcal K_{\mathrm{Open}}$ and
$\mathcal L$ derive the same sequents.
\end{theorem}
\begin{proof}
The rules of $\mathcal L$ are derivable by \cref{ctr:lem:safe,ctr:lem:openrules}.
Conversely the introduction is an instance of \rulename{Case} or \rulename{Prog}. For
\eqref{ctr:eq:open} with a successful probe at $r$, apply \rulename{DE} at $b$; the opened
clause has the antecedent $0=0$, and \rulename{Use} at $1$ gives $b=r$. For a fresh $y$,
\rulename{Use} gives $B(b,y)$ and transport along $b=r$ gives $B(r,y)$; hence the clause
$[B(r,0);\,y.\,B(r,y)]$ of the branch can be cut (\cref{ctr:lem:libstruct}), and the branch
gives $A$. With a failed probe the rule is \rulename{DE}. The comparison of $\mathcal O$ with
the reference calculus in \cref{disch:thm:open} is the same argument with the clause cut
replaced by a closure.
\end{proof}

\subsection{The scalar normal forms embed exactly}

The calculus of \pkg{inferred\_witness} is $\mathcal K_{\mathrm{Open}}$ on iterator terms
$0,\Suc,I$ in full normal form: a normal iterator term $\It{n}{a}{v}{h}$ has a neutral
counter, and the normalizer $\nf_I$ applies both computation equations of
\cref{base:sec:terms}. Let $T_{\mathrm{raw}}$ translate variables and $0$ identically, the
successor by the defined successor, and $\It{n}{a}{v}{h}$ by
$\operatorname{Rec}(n,a;i,v.h)$ with $i$ fresh, and for a normal scalar term $m$ put
$T(m)\defeq\nf(T_{\mathrm{raw}}(m))$.

\begin{lemma}[Exact embedding of normal forms]\label{ctr:lem:embedding}
For every scalar term $t$, $\nf(T_{\mathrm{raw}}(t))=T(\nf_I(t))$. The raw translation
commutes with substitution and $T$ commutes with normalized substitution. On normal scalar
terms $T$ preserves and reflects the literal $0$, each variable and the property of being a
successor, and it preserves the set of free variables.
\end{lemma}
\begin{proof}
The two scalar computation equations translate to the computations of
\cref{ctr:lem:decoder}, so $T_{\mathrm{raw}}(t)$ and $T_{\mathrm{raw}}(\nf_I(t))$ have the
same normal form by \cref{ctr:lem:nf}. For a normal scalar term describe $T$ by induction.
Zero, variables and successors are immediate. A normal iterator has a neutral counter, and
its image is
\[
 P\bigl(R(Tn,Ta;i,u.\Suc((Th)[P(u)/v]))\bigr).
\]
Here $Tn$ is neutral, so the inner recursor is retained and neutral, and the outer $P$ is
retained. Substituting the neutral term $P(u)$ for the body variable does not turn a neutral
counter into $0$ or a successor, so the whole image is normal. The raw translation commutes
with substitution by induction on terms, with the two new variables fresh at every
translated iterator, and \eqref{ctr:eq:SC} gives the normalized statement. The four kinds of
normal source term have distinct kinds of image: $0$, the same variable, a canonical
successor, a neutral recursor. No free variable of a normal source term is erased.
\end{proof}

\begin{proposition}[Translation of the scalar calculus]\label{ctr:prop:scalartrans}
Extended to formulas without change of roots, $T$ maps every instance of \rulename{Use}, of
the introduction and of \eqref{ctr:eq:open} of the normalized scalar calculus to an instance
of the same rule of $\mathcal K_{\mathrm{Open}}$.
\end{proposition}
\begin{proof}
For \rulename{Use} this is commutation with normalized substitution. For the introduction,
reflection of the literal $0$ preserves the mode, and the translated successor row is the
required successor row. For \eqref{ctr:eq:open} the two tests on rows are preserved and
reflected, so $\tau_{T(B)}(b)=T(\tau_B(b))$, and normalized substitution gives the translated
branch clause. Preservation of free variables gives the scope conditions.
\end{proof}

The translation concerns derivations of the already normalized scalar calculus. If raw
scalar proofs are normalized first, a witness can become $0$; that separate step is handled
by \cref{ctr:lem:safe}(3). The embedding preserves more than standard numerical values: a
row becomes $0=0$, or has a given fresh variable as its entire left side, exactly when the
source row did.

\subsection{The main theorem for the two calculi}

\begin{theorem}\label{ctr:thm:threemain}
$\mathcal K_{\mathrm{Name}}$ and $\mathcal K_{\mathrm{Open}}$ are sound and faithful for
$\mathsf T_R^2$ and are equiconsistent and bi-interpretable with $\iPRA$, with whole
domains, actual equality, numerical composites of graph $x=y$, and every formula provably
equivalent to its double translation. Every fixed primitive recursive description compiles
to a term whose defining equations hold by computation
(\cref{ctr:thm:compiler}).
\end{theorem}
\begin{proof}
By \cref{ctr:thm:nameequiv,ctr:thm:openequiv} both calculi are $\mathcal L$ on normal terms,
and \cref{ctr:thm:adequacy,ctr:thm:transfer} apply. In particular
$\mathcal K\proves 0=1$ if and only if $\iPRA\proves\bot$; this is a metatheorem and not an
assertion that either system proves its own consistency.
\end{proof}

On normal terms the coverage theorem of \cref{ctr:lem:macros} strengthens to
$t=0\lor_*t=\Suc(p(t))$, since $p(\Suc z)\equiv z$; this is the form used by
\pkg{uniform\_recursion} and \pkg{single\_constructor}.

\section{Uniform binding and one traversal}\label{ctr:sec:binding}

This section draws on \pkg{clausal\_canonical}, \pkg{counter\_evaluation},
\pkg{uniform\_binding}, \pkg{uniform\_recursion} and \pkg{single\_constructor}, and for the
separate refinements on \pkg{canonical\_recursors} and \pkg{register\_graph}.

\subsection{The concrete syntax}

In the checkers of these packages the integer $0$ is the numerical constant, a negative
integer is a free parameter (a level), and a positive integer is a de Bruijn index of a
bound variable. Without a successor constructor and without a closure constructor, the
array arities distinguish the constructors:
\begin{center}
\small
\begin{tabular}{@{}llll@{}}
\toprule
Object & Array & Child depths, one binder in $R$ & Child depths, two binders in $R$\\
\midrule
$R(n,a;\vec u.h)$ & $[n,a,h]$ & $(d,d,d+1)$ & $(d,d,d+2)$\\
$s=t$ & $[s,t]$ & $(d,d)$ & $(d,d)$\\
$\D{x}{y}B$ & $[B]$ & $(d+2)$ & $(d+2)$\\
\bottomrule
\end{tabular}
\end{center}
In a paired matrix $y$ has index $1$ and $x$ index $2$; in the body of the counter-aware
recursor the state has index $1$ and the counter index $2$. With two binders in $R$, every
constructor other than the equation binds exactly two variables in its last child, so that
for child number $j$ of a node of arity $k$ the depth increment is
\begin{equation}\label{ctr:eq:uniform}
 2\,[\,k\neq 2 \text{ and } j=k-1\,].
\end{equation}
The sort is still checked: formula arities are $1$ and $2$, and the only arity of a
compound term is $3$. The arity convention does not identify a formula with a term. A
successor is $[0,a,0]$, and a compound canonical counter is a successor exactly when its
first child is $0$; the test that distinguishes a successor node from an iterator node on
the terms of \cref{base:sec:terms} is not needed.

A context clause is a pair $(A,B)$ whose second component has one more bound variable. It
is a separate data structure and is never passed to the opening operation as one object;
clause use opens $B$ alone.

\begin{definition}[Opening]
For an expression $e$ and a term $v$, the operation $O_d(e,v)$ removes the index $d+1$ and
preserves the $d$ nearer indices: on an integer it replaces $d+1$ by $v$, decrements an
integer greater than $d+1$ and leaves every other integer unchanged; on an array it
processes the children at the depths of the table; and at each resulting recursor node it
erases the body if the resulting counter is $0$. $O_d^{\mathrm{raw}}$ is the same operation
without the erasure.
\end{definition}

\begin{lemma}[Opening]\label{ctr:lem:opening}
Let $e$ have free scope $m$ and ambient depth $k>d$, and let $v$ be a term of free scope
$m$ with no free positive index.
\begin{enumerate}
\item $O_d(e,v)$ has the sort of $e$, ambient depth $k-1$ and free scope $m$; it contains
no parameter beyond those of $e$ and $v$; and it is zero-body canonical if $e$ and $v$ are.
\item For canonical $e,v$, $O_d(e,v)=\nu(O_d^{\mathrm{raw}}(e,v))$, and for arbitrary raw
$e,v$, $\nu(O_d^{\mathrm{raw}}(\nu e,\nu v))=\nu(O_d^{\mathrm{raw}}(e,v))$.
\item For a matrix $B$ and terms $t,w$ without free positive indices,
$O_0(O_1(B,t),w)=\nu(B[t/x,w/y])$, and the two openings may be performed in the opposite
order.
\end{enumerate}
\end{lemma}
\begin{proof}
(1) Induct on $e$. The integer cases remove the selected index, lower an outer index, or
keep a nearer index or a parameter. A replacement without free positive indices stays well
scoped under any additional binders. The array cases use the binder increments of the
table. Erasing a body removes occurrences only. (2) is \cref{ctr:lem:coherence} for the
substitution that opening performs: a body discarded before the opening had a counter that
was literally $0$ and stays $0$; a counter that becomes $0$ only by the opening is erased on
both sides by the outer $\nu$; and $\nu(t)$ is $0$ only if $t$ is. (3) follows from the
simultaneous-substitution lemma, since neither replacement contains a free positive index,
and from (2).
\end{proof}

The hypothesis on $v$ is essential. Because free parameters are negative levels and do not
depend on the binder depth, inserting such a $v$ under binders needs no lifting; inserting a
term with a free positive index without lifting is incorrect. For the same reason an old
context needs no shifting when a fresh parameter is allocated: the new negative integer is
absent from everything validated before. A branch can use it, but its discharged result
cannot contain it, because the outer goal and the witness are validated at the old scope.

\begin{remark}[Unfolding under binders]\label{ctr:rem:counterfirst}
With full normalization a successor contraction can occur under binders, and the inserted
terms can then have free positive indices. The simultaneous substitution of
\eqref{ctr:eq:r2succ} is implemented counter first: remove the counter slot at depth $1$,
under the retained state slot and with lifting of the inserted term, and then remove the
state slot. Substituting the previous state first without lifting captures an ambient index
under the counter binder that is still present. For example, under one ambient binder
$R(1,v;j,w.w)$ normalizes to $\Suc v$, that is $[[0,0,0],1,1]$ normalizes to $[0,1,0]$,
whereas the state-first order returns $\Suc 0$. This correction is recorded in the
mathematical review shipped with \pkg{uniform\_recursion}.
\end{remark}

\subsection{One traversal}

\begin{lemma}[The shared walker]\label{ctr:lem:walker}
Let $\mathrm{walk}$ visit a tree with the depth increments \eqref{ctr:eq:uniform}, check
the sort and arity of every node, apply a given operation to the numerical leaves, and
erase the body at every recursor node whose resulting counter is $0$.
\begin{enumerate}
\item With the leaf operation that checks $-n\le i\le d$ for an exact integer $i$ at scope
$n$ and depth $d$ and returns $i$, the test that the result equals the input holds exactly
for the members of the scoped zero-body canonical grammar.
\item With the leaf operation that maps $d+1$ to $v$, an index $i>d+1$ to $i-1$ and every
other leaf to itself, and with $v$ canonical without free positive indices, the result is
the canonical capture-avoiding opening $O_d$.
\end{enumerate}
\end{lemma}
\begin{proof}
Induct on the tree. The integer clause checks the sort and the scope interval. At an array
the admissible arity determines whether the children are terms or formulas, and the last
child of a recursor or a paired formula receives two more protected indices. Every child is
visited before the erasure, so a malformed body is rejected before it can be discarded. In
(1), equality with the returned tree excludes every zero-counter node with a nonzero body.
In (2) the three index cases are opening of one slot, and the final erasure is that of
$O_d$.
\end{proof}

Thus one traversal owns grammar validation, the canonical-form check and opening. For a
certificate that starts from the empty context every branch of an elimination or
introduction adds exactly one clause and one fresh parameter, and no other rule adds
either; by induction on checker calls the scope equals the length of the context, so no
separate scope counter is passed. This is an invariant of certificates built from the empty
root, not a statement about abstract clause contexts. In the rule \eqref{ctr:eq:CE} the
checker opens the state slot with the canonical term $R(t,a;i,u.h)$ and then the former
counter slot with $t$, which is the simultaneous substitution; the opening is computed
before the two premises are checked, and this does not bypass validation, because proving
the major premise validates both endpoints and proving the minor premise validates the seed
and the original body, whose outer counter is a successor and cannot have been erased.

\begin{proposition}[Checker correspondence]\label{ctr:prop:checkers}
Disregarding physical resource bounds:
\begin{enumerate}
\item A closed canonical formula has a certificate accepted by the goal-directed checker of
\pkg{uniform\_binding} if and only if it is a theorem of $\mathcal K_{\mathrm{CE}}$, and
accepted certificates translate effectively to and from the certificates of the
conclusion-synthesizing checker of \pkg{counter\_evaluation}, with literally identical final
formulas.
\item The checker of \pkg{uniform\_recursion} accepts a certificate for a closed normal goal
exactly when the certificate describes a derivation in $\mathcal K_{\mathrm{Name}}$, and
every finite derivation has an accepted certificate; the same holds for
\pkg{single\_constructor} and $\mathcal K_{\mathrm{Open}}$.
\end{enumerate}
\end{proposition}
\begin{proof}
(1) Use the scope invariant and \cref{ctr:lem:walker}. An accepted reflexive goal is a
\rulename{Refl} node, and every other accepted branch is the displayed instance of
\eqref{ctr:eq:CE}, \rulename{Use}, \rulename{DE} or \eqref{ctr:eq:CEDI} after translation of
its premises; this is soundness, by induction on the accepted call tree. Conversely induct
on a derivation: insert the derived endpoints and the seed and body of the minor premise as
the hints of a \eqref{ctr:eq:CE} node, copy the clause index and target of \rulename{Use},
the matrix of \rulename{DE} and the witness of the introduction. The opening invariant
identifies every required premise goal, and each construction preserves the end formula.

(2) Induct on the proof packets. The formula synthesized at each node is checked for
grammar, normal form and scope. In \rulename{Use} the term is checked before instantiation
and a proof of the antecedent has to synthesize exactly the retained antecedent. In the
introduction the matrix and the witness are checked, the generated assumption is the row
selected by the rule, and the result of the branch has to equal the generated successor
row. In \rulename{Name} the major premise has to be a paired formula, the equality has to
have the freshly allocated parameter as its left endpoint, and the final scope check
excludes that parameter from the normalized export. In \eqref{ctr:eq:open} the major premise
is checked before the context is extended, the probe \eqref{ctr:eq:probe} is computed, and
the conclusion of the branch is validated again in the outer scope; when the probe
succeeds the stored base is $0=0$ and coincides with the substituted base of the rule.
\Cref{ctr:lem:nf} and \cref{ctr:rem:counterfirst} identify the computed instances with the
normalized substitutions. Conversely a derivation is encoded by its rule data and the
certificates of its premises, with fresh negative levels in order of scope creation.
\end{proof}

\subsection{The checkers}

All checkers are standalone Python modules that read one JSON object with exactly the keys
\texttt{goal} and \texttt{proof}, validate a closed goal against the full grammar, and
accept by printing \texttt{ACCEPT} with exit status zero; none imports a proof producer or a
theorem table. Exact integer tests exclude Booleans, and recursion and memory limits of the
interpreter can reject a large valid certificate. Sizes are in \cref{ctr:sec:sizes}.

\emph{\pkg{register\_graph}.} The checker implements the natural deduction of
\cref{ctr:def:TJ} with named variables, simultaneous substitution, alpha equivalence, the
three arithmetic schemas and \eqref{ctr:eq:EI}. It is a readable reference and is much
larger than the compact checkers of the kernel. A separate ancillary module is a profile
of the kernel $\Nat$ on iterator terms, not a checker for $\mathsf T_J$: it replaces a
recursive continuation of its parent by a loop and uses immutable tuple contexts, which
lowers the node count of the inspected leader by one and raises its bytes. Its generic
opener inherits a qualification on malformed input: in an interpreter without limits
certain nonempty strings cause unbounded traversal, so soundness concerns finite successful
executions.

\emph{\pkg{canonical\_recursors}.} Three modules share one proof-tree format, context
handling, de Bruijn convention and driver: iterator terms with a successor and seven rules,
raw seeded terms with seven rules, and $\mathcal K_6$. The canonical module applies the
constructor $[n,a,\,0 \text{ if } n=0 \text{ else } h]$ at every newly built recursor node
after its children have been normalized, validates canonical input and has no
zero-computation rule. The scalar control uses crossed eliminations and the
witness-directed introduction and is an equivalent presentation of the formalized seven-rule
calculus named in \cref{ctr:sec:six}, not a transcription of it.

\emph{\pkg{clausal\_canonical}.} The checker for $\mathcal K_5$ has the instruction tags
\texttt{r}, \texttt{x}, \texttt{u}, \texttt{e}, \texttt{i} and one opening function for
terms and formulas. Its matched control is $\mathcal K_6$ ported to the same variable
convention and opening function, with the closure constructor, the restricted transport and
the root-reversing closure elimination retained, so that the comparison isolates the change
of rules and formulas from the change of representation.

\emph{\pkg{counter\_evaluation}.} The checker for $\mathcal K_{\mathrm{CE}}$ has the tags
\texttt{r} (reflexivity), \texttt{v} (counter evaluation), \texttt{u}, \texttt{e},
\texttt{i} and synthesizes the conclusion of every node. Its matched control is
$\mathcal K_{\mathrm{CT}}$ on the same syntax. An untrusted translator implements
\cref{ctr:thm:interface}: a computation node becomes \eqref{ctr:eq:CE} on two reflexive
leaves and a transport node becomes the construction of \cref{ctr:prop:cost}.

\emph{\pkg{uniform\_binding}.} The checker for $\mathcal K_{\mathrm{CE}}$ is goal directed:
a validated reflexive goal closes at once and its payload is ignored, and the other
instructions are \texttt{v}, \texttt{u}, \texttt{e}, \texttt{i} with hints. Validation and
opening are the two uses of \cref{ctr:lem:walker}, and the scope is the length of the
context. Its matched control applies the same organization to iterator terms.

\emph{\pkg{uniform\_recursion}.} The checker for $\mathcal K_{\mathrm{Name}}$ has the
packets \texttt{u}, \texttt{i}, \texttt{e}; the fixed clause \eqref{ctr:eq:reflclause} is
the initial context $([\,],[1,1])$, so that at context length $n$ the parameters are
$-1,\dots,-(n-1)$ and the next one is $-n$. The minimal proof is
\texttt{\{"goal":[0,0],"proof":["u",0,0,null]\}}. A variant carries the budget of
\cref{ctr:rem:fuel}; its control is the scalar control with the same budget instrumentation.

\emph{\pkg{single\_constructor}.} The checker for $\mathcal K_{\mathrm{Open}}$ has the
packets \texttt{u}, \texttt{i}, \texttt{e}, with the same three logical tests as its
control on iterator terms; the changes are confined to the arithmetic syntax, the shared
binder increment, the recursor computation and the successor abbreviation in the rules.

\subsection{Three refinements of checkers for the kernel on iterator terms}

Three packages of this chapter also contain refinements of existing checkers that do not
concern the calculi of this chapter; no size of these programs is a size of a calculus on
the new syntaxes.

\pkg{canonical\_recursors} refines the compact guarded codec: it pairs the comparisons of
the major premise and of the target that are already checked, uses a shorter discriminator
between integers and lists, and composes the administrative truth convention of the scalar
zero with the guarded checker. The byte and node improvements belong to different modules.
The compact modules do not validate the public grammar themselves; a strict front end is
counted separately. The eager paired comparison may inspect instruction fields that were
unused before, so the completeness statement is made for the specified well-scoped language
of integer and list certificates, and no equality of behavior on malformed JSON is claimed.

\pkg{uniform\_recursion} fuses the grammar validation of the strict native checker of
\pkg{deterministic\_opening} into its first scope traversal. That walker increases its
depth by $5-k$ at a node of arity $k$. At a formula position under $p$ protected slots the
depth is $2p$, and at a term position it is $3+2p$; hence even depth means formula and odd
depth means term, and the admissible arities at depth $d$ are exactly $1$ and
$2+(d\bmod 2)$. A numerical leaf is $0$ or an odd integer at least $3$ within the scope
interval. The grammar predicate is enabled for the initial traversal without substitution
image and disabled before any proof instruction is read. By induction on the tree the fused
traversal succeeds exactly when the old grammar test and the old scope traversal both
succeed, and it returns the same tree; every later transition of the checker is unchanged,
so the translation of certificates in both directions is the identity. Thresholds of
resource exhaustion and failure diagnostics are not equated. One variant states the
predicate explicitly; a smaller one folds it into the existing cases.

\pkg{register\_graph} supplies the ancillary profile described above.

\section{Boundary results}\label{ctr:sec:bounds}

This section draws on \pkg{canonical\_recursors} and, for \cref{ctr:prop:norule}, on
\pkg{clausal\_canonical}.

The constructions above move a computation into the syntax or change the representation of
the terms. The results of this section show that three nearby simplifications are not
available: deleting the initializer together with the counter, deleting one of the two
computation equations of the iterator, and restricting induction to equations with $0$ on
one side. A fourth shows that the single computation rule of $\mathcal K_5$ cannot be
deleted without replacement.

\subsection{A recursor with neither counter nor initializer}

Consider the terms generated by variables, $0$, $\Suc$ and a self-seeded constructor
$F(n;v.h)$ with the standard interpretation
\[
 F(0;v.h)=\Suc(h(0)),\qquad F(\Suc n;v.h)=h\bigl(F(n;v.h)\bigr),
\]
with arbitrary nesting and passive parameters, and no further operation. The successor may
be taken as the abbreviation $F(0;v.a)$, $v\notin\FV(a)$. This is the unary-update variant
$\mathsf F$ of the tied-seed language of \cref{obs:def:languages}(b), the instance
$\mathsf B(n;u,v.h)$ with $u$ not in $h$. Its terms are almost unary by
\cref{obs:cor:tied-seed}, which is no restriction on unary functions; the theorem below is a
lower bound for unary terms and uses the finite-control invariant of \cref{obs:thm:fc}.

\begin{theorem}[Unary terms of the self-seeded language]\label{ctr:thm:unary}
For every term $t(n)$ of this language with one free variable there is a natural number $C$
such that
\[
 \forall n\in\mathbb N\quad\bigl(t(n)\le C\ \lor\ t(n)\ge n\bigr).
\]
Consequently no term defines the predecessor $\max(n-1,0)$, integer halving, or any function
with unbounded values strictly below its argument.
\end{theorem}

Let $\mathcal U$ be the class of functions $f\colon\mathbb N\to\mathbb N$ for which there is
$C$ with $f(n)\le C$ or $f(n)\ge n$ for every $n$.

\begin{lemma}[Closure]\label{ctr:lem:Uclosure}
Constants, the identity and the successor belong to $\mathcal U$. The class is closed under
composition and under pointwise selection from a finite family of its members.
\end{lemma}
\begin{proof}
For a pointwise selection take the maximum of the finitely many bounds. For $f\circ g$ with
bounds $C_f,C_g$: if $g(n)\le C_g$ the value is at most $\max_{j\le C_g}f(j)$; otherwise
$g(n)\ge n$, and the value is at most $C_f$ or at least $g(n)\ge n$.
\end{proof}

\begin{lemma}[Orbits]\label{ctr:lem:orbit}
If $h\in\mathcal U$, then the sequence $r(0)=1+h(0)$, $r(n+1)=h(r(n))$ is a member of
$\mathcal U$ as a function of $n$.
\end{lemma}
\begin{proof}
If the orbit repeats a value it is eventually periodic and hence bounded. Otherwise it is
injective. Let $C$ be a bound for $h$. A step with $r(n+1)<r(n)$ ends in $\{0,\dots,C\}$,
so an injective orbit has finitely many such steps. After the last one the orbit increases
strictly, and it eventually exceeds every value of the earlier finite prefix. For all
sufficiently large $n$ the $n+1$ distinct values $r(0),\dots,r(n)$ lie in $[0,r(n)]$, so
$r(n)\ge n$. Increase the constant to cover the finitely many earlier arguments.
\end{proof}

\begin{proof}[Proof of \cref{ctr:thm:unary}]
% ed.: the package strengthens a finite-control invariant that it cites from a reference file; the invariant is that of obs:def:fc with profiles in U, and the steps that coincide with obs:lem:fc-comp and obs:thm:fc are cited.
Say that a function $g\colon\mathbb N^k\to\mathbb N$, $k\ge 1$, has \emph{$\mathcal U$-control}
if it has finite control in the sense of \cref{obs:def:fc} with every profile in
$\mathcal U$. A finite-valued feature of the active coordinate may be added to the index:
for fixed features of the other coordinates the function is then a pointwise selection from
finitely many profiles, which belongs to $\mathcal U$ by \cref{ctr:lem:Uclosure}. A unary
function with $\mathcal U$-control is its own profile and belongs to $\mathcal U$. It
therefore suffices to show that every term, as a function of any nonempty list of variables
that contains its free variables, has $\mathcal U$-control.

\emph{Composition.} Functions with $\mathcal U$-control are closed under composition. The
proof of \cref{obs:lem:fc-comp}(ii) shows that the composite has finite control, and that
for fixed combined controls its profile is $x_k\mapsto F_{a(x_k)}(G(x_k))$, where $G$ is a
profile of the inner function, the $F_a$ are profiles of the outer function, and the index
$a(x_k)$ ranges over finitely many values determined by finite features of $x_k$
(\cref{obs:lem:fc-comp}(i)). This profile is a pointwise selection from finitely many
compositions $F_a\circ G$ of members of $\mathcal U$, and belongs to $\mathcal U$ by
\cref{ctr:lem:Uclosure}.

\emph{Induction on terms.} Variables, $0$ and successors are immediate, the last by
composition with the successor. For $F(n;v.h)$ first replace the counter by a fresh
variable $m$. By the induction hypothesis $h(v,\vec x)$ has $\mathcal U$-control with an
active coordinate $y$; write $c=\beta(\vec x)$ for the features of the passive coordinates
other than $y$ and $\beta_v$ for the feature of $v$ when $v$ is not active.

If $y$ is the state $v$, fix $c$. The update $h_c(v)=H_c(v)$ is a profile and belongs to
$\mathcal U$, and by \cref{ctr:lem:orbit} so does its orbit $r_c$. Then
$F(m;v.h)=r_{\beta(\vec x)}(m)$ has $\mathcal U$-control with the counter active.

If $y$ is a passive parameter $p$, the catalogue argument of the third case of the proof of
\cref{obs:thm:fc} applies to $F(m;v.h)=\mathsf B(m;u,v.h)$, with the old counter $u$
absent, so that no feature of the counter enters the transitions. The catalogue consists of
the update profiles $H_{(b,c)}(p)$ and the initial profiles $1+H_{(\beta_v(0),c)}(p)$; it
yields
\[
 F(m;v.h)=V_{\eta(m)_{\beta(\vec x),\kappa(p)}}(p),\qquad
 \kappa(p)=\bigl(\beta_v(V_k(p))\bigr)_k,
\]
with a finite-valued feature $\eta$ of the counter and with every catalogue function $V_k$
in $\mathcal U$ by \cref{ctr:lem:Uclosure}. This is $\mathcal U$-control with $p$ active.
The argument uses only finite observations and no test of numerical equality.

Finally substitute the actual counter term $n$ for $m$ by the composition step. This
completes the induction and proves the bound. For the predecessor choose $n>C+1$: its
value exceeds $C$ and is less than $n$. Halving fails at every $n>2C+1$.
\end{proof}

\begin{remark}
The theorem concerns definitions by terms under the standard interpretation. It is not an
obstruction to interpretations by graphs or on changed domains. The counter-aware tied-seed
language of the package \pkg{local\_scopes} has the finite-control obstruction of
\cref{obs:thm:fc} but defines the predecessor (\cref{obs:lem:Bterms}); the unary bound is
specific to removing the access to the old counter together with the independent
initializer. The majorant obstructions of \cref{obs:thm:single-body,obs:cor:tied-seed}
bound terms from above by a function of one coordinate and say nothing about unary terms;
the theorem above is a dichotomy for unary terms, which they do not give.
\end{remark}

\subsection{The two computation equations of the iterator are independent}

\begin{proposition}\label{ctr:prop:indep}
Neither computation equation of the iterator of \cref{base:sec:terms} follows from the
other, even over the standard natural numbers with coverage, separation and every
first-order induction instance.
\end{proposition}
\begin{proof}
Interpret every fixed body compositionally in $\mathbb N$. If
$\It{n}{a}{u}{h}$ is interpreted as $h^n(0)$, the successor equation holds and
$\It{0}{1}{u}{u}=0\neq 1$. If it is interpreted as $a$, the zero equation holds and the
successor equation fails for $h(u)=\Suc u$.
\end{proof}

Thus the canonical calculi succeed by changing the representation; they do not derive a
deleted law from the remaining axioms.

\subsection{Induction on zero equations does not give induction on equations}

Let $\mathsf T_0$ be the theory $\Teq$ of \cref{base:sec:teq} with its induction rule
restricted to equations of the form $0=t(n,\vec p)$, or equivalently $t(n,\vec p)=0$, with
an arbitrary context. A calculus that uses the progress introduction only for matrix
equations with the literal $0$ on one side has this induction. The following theorem shows
that such a restriction does not preserve the consequence relation.

\begin{theorem}[A model separating the two inductions]\label{ctr:thm:zeroind}
There is a classical model of $\mathsf T_0$, in which the successor is injective, and in
which the sentence
\[
 \forall n\,\forall a\ \It{n}{a}{u}{u}=a
\]
is false, although its instance at $n=0$ and its induction step are provable from the
computation equations and equality logic. Consequently induction on arbitrary literal
equations is not admissible in $\mathsf T_0$. The same model satisfies induction for every
equation $\bar k=t(n,\vec p)$ with a fixed standard numeral $\bar k$ on one side.
\end{theorem}

A classical model is in particular a model of intuitionistic logic. Its construction uses
compactness and choices of compatible paths in the metatheory; these are not principles
granted to $\mathsf T_0$.

\begin{lemma}[Paths through finite fibers]\label{ctr:lem:path}
Let $(B_j)_{j\in\mathbb Z}$ be finite nonempty sets and $f_j\colon B_j\to B_{j+1}$ total
maps. There is a sequence $(b_j)_{j\in\mathbb Z}$ with $b_j\in B_j$ and
$b_{j+1}=f_j(b_j)$ for every $j$.
\end{lemma}
\begin{proof}
The product $\prod_jB_j$ of finite discrete spaces is compact, and each equation between
adjacent coordinates is a closed condition. Every finite set of these conditions is
satisfiable: choose a value at the least index involved, propagate it forward, and choose
arbitrary values elsewhere. The finite intersection property gives a common solution.
\end{proof}

\begin{proof}[Proof of \cref{ctr:thm:zeroind}]
\emph{The base.} Fix a nonstandard model $M$ of Peano arithmetic expanded by symbols for all
primitive recursive descriptions; it exists by compactness. For $m\in M$ let $A_m=\{0\}$ if
$m$ is a standard numeral and $A_m=\{0,1\}$ otherwise; the distinction is external. Put
\[
 D=\{(m,i): m\in M,\ i\in A_m\},\qquad \pi(m,i)=m,\qquad
 0_D=(0_M,0),\qquad \Suc_D(m,i)=(\Suc_Mm,i).
\]
The successor is well defined because it preserves standardness and nonstandardness. The
map $\pi$ is surjective and preserves zero and successor. Every standard numeral has a
one-element fiber and every nonstandard element has two lifts. On $D$ the successor is
injective, zero is not a successor, and coverage holds: a nonstandard $(m,i)$ has the
predecessor $(m-1,i)$. The nonstandard part of $M$ is a disjoint union of chains of the
order type of $\mathbb Z$ under successor, and the label $i$ is constant along each chain
of $D$.

\emph{The iterator on actual functions.} Let $\mathscr H$ be the set of functions
$h\colon D\to D$ for which there is a unary function $f\colon M\to M$, primitive recursive
with parameters, such that
\begin{equation}\label{ctr:eq:projection}
 \pi(h(x))=f(\pi(x))\qquad(x\in D).
\end{equation}
By surjectivity of $\pi$ the function $f$ is unique. Write
$\operatorname{It}_f(m,b)$ for its iteration in $M$ with initial value $b$. This does not
depend on the description chosen for $f$: two descriptions of the same function give
iterations with the same base and pointwise equal updates, and an induction in $M$ proves
them equal at every counter. This induction is used in the base model, not in
$\mathsf T_0$. Define an operation
$\mathcal I\colon D\times D\times\mathscr H\to D$ as follows. For a standard counter $k$
put $\mathcal I(k,a,h)=h^k(a)$, the finite iteration on $D$. For fixed $a,h$ and a
nonstandard chain $n_j=(m_j,i)$, $j\in\mathbb Z$, with $m_{j+1}=\Suc_Mm_j$, put
\[
 r_j=\operatorname{It}_f(m_j,\pi(a)),\qquad B_j=\pi^{-1}(r_j).
\]
The sets $B_j$ have one or two elements, $r_{j+1}=f(r_j)$ by the computation law in $M$,
and \eqref{ctr:eq:projection} makes $h$ a total map $B_j\to B_{j+1}$. Choose a path
$b_j\in B_j$ by \cref{ctr:lem:path} and put $\mathcal I(n_j,a,h)=b_j$. Make these choices
for all triples of a function, an initializer and a chain, with one specified family: for
$h=\mathrm{id}_D$ choose on every nonstandard chain the constant path
\begin{equation}\label{ctr:eq:idchoice}
 \mathcal I(n,a,\mathrm{id}_D)=(\pi(a),0).
\end{equation}
This is legitimate, since the projected iteration is constantly $\pi(a)$, the point with
label $0$ always exists, and the identity fixes it. The choices depend on the actual
function $h$, so every description of the identity body receives the same family. By
construction
\[
 \mathcal I(0_D,a,h)=a,\qquad \mathcal I(\Suc_Dn,a,h)=h(\mathcal I(n,a,h)),\qquad
 \pi(\mathcal I(n,a,h))=\operatorname{It}_f(\pi(n),\pi(a)),
\]
the standard and the nonstandard counters never meeting under a finite number of successor
steps.

\emph{Terms.} Interpret
$\It{n}{a}{u}{h}$ in $D$ under a valuation $\rho$ as
$\mathcal I(\llbracket n\rrbracket,\llbracket a\rrbracket,v\mapsto\llbracket h\rrbracket_{\rho[u:=v]})$.
By induction on the term, with a bound variable of a body treated as a further free
variable,
\begin{equation}\label{ctr:eq:termproj}
 \pi\bigl(\llbracket t\rrbracket_{D,\rho}\bigr)=\llbracket t\rrbracket_{M,\pi\circ\rho},
\end{equation}
where the right side uses the iterator of $M$; in particular every body function supplied
to $\mathcal I$ satisfies \eqref{ctr:eq:projection} with the interpretation of the body in
$M$ and belongs to $\mathscr H$. The operation $\mathcal I$ is chosen for actual functions
before this induction, so nested terms create no circular definition. Equal parameter
values determine the same body function and the same iterator; hence $D$ is a first-order
structure for the staged language, the substitution lemma holds, and both computation
schemas hold for every fixed body.

\emph{Zero equations.} Since the fiber over $0_M$ has one element,
$\llbracket t\rrbracket_D=0_D$ if and only if $\llbracket t\rrbracket_M=0_M$ at the
projected valuation. Fix parameters $\vec a$ in $D$ and suppose that in $D$
\[
 t(0,\vec a)=0,\qquad \forall n\,\bigl(t(n,\vec a)=0\to t(\Suc n,\vec a)=0\bigr).
\]
By \eqref{ctr:eq:termproj} and surjectivity of $\pi$ the same two statements hold in $M$
for the primitive recursive interpretation of $t$ at $\pi\vec a$. Induction in $M$ and
reflection of zero along $\pi$ give $\forall n\in D\ t(n,\vec a)=0$. The rule with a context
is valid as well: at a valuation that makes the context true the two premises supply the
displayed base and universal step, and no projection of the context is needed. The same
argument applies to an equation $\bar k=t(n,\vec a)$, because every standard numeral has a
one-element fiber.

\emph{Failure of the full rule.} Choose a nonstandard $b\in M$ and $a=(b,1)$, and let
$E(n)$ be $\It{n}{a}{u}{u}=a$. Its base is a computation equation, and its step is provable
from $\It{\Suc n}{a}{u}{u}=\It{n}{a}{u}{u}$ and transitivity. For every nonstandard counter
$n$, however, \eqref{ctr:eq:idchoice} gives
$\It{n}{a}{u}{u}=(b,0)\neq(b,1)=a$. The two elements have the same projection; tests
against zero cannot distinguish the labels.
\end{proof}

\begin{corollary}[No reflection of equality]\label{ctr:cor:noreflection}
$\mathsf T_0$ proves for no term $e(x,y)$ the law
$e(x,y)=0\leftrightarrow x=y$.
\end{corollary}
\begin{proof}
In the model take $x=(b,0)$ and $y=(b,1)$ over a nonstandard $b$. By
\eqref{ctr:eq:termproj} the terms $e(x,y)$ and $e(x,x)$ have the same projection. If the law
held, $e(x,x)=0$, and the one-element fiber over zero would force $e(x,y)=0$, contradicting
$x\neq y$.
\end{proof}

The usual conversion of an equation into a zero equation therefore cannot justify the
restriction: the reflection of equality of \cref{base:thm:reflect} is exactly the missing
conclusion. The theorem refutes the recovery of the induction rule of $\Teq$ in the
presence of full intuitionistic logic with equality, all numerical cases, both computation
schemas and injectivity of the successor. Under the reading of \cref{base:sec:reading}, the
rules \rulename{QE}, \rulename{DE}, introduction by cases and progress for zero equations
are valid in the same model, so the native rules do not evade the counterexample. The
construction does not show that $\mathsf T_0$ has no interpretation of arithmetic on a
quotient or on a changed domain (\cref{ctr:sec:questions}).

\begin{remark}
\Cref{ctr:thm:zeroind,ctr:cor:noreflection} are term-level model constructions for a
restricted induction rule; the obstructions of \cref{obs:sec:profiles} concern
term-definability, and the nonstandard model of \cref{obs:sec:progress-independence}
concerns the deletion of the progress rule, not its restriction to zero equations; neither
contains them. The model above refutes the rule of
successor uniqueness of the package \pkg{orbit\_closure}, as recorded in
\cref{kern:rem:sugap}, and therefore does not decide \cref{kern:q:su}.
\end{remark}

\subsection{The last computation rule}

\begin{proposition}\label{ctr:prop:norule}
The calculus $\mathcal K_5$ without rule (R) derives no equation from the empty context,
not even $0=0$.
\end{proposition}
\begin{proof}
Assign the value false to every equation and the value true to every paired formula, and
read a clause $[A;\,y.\,B]$ as the implication from the value of $A$ to the value of $B$.
This agrees with the reading of the paired binder as a conditional closure when all
equations are false, and it does not impose the equality axioms. By induction on
derivations, if all clauses of $\Gamma$ are true, then the conclusion of $\Gamma\proves A$
is true. Rule (T) has equations as premises, so it applies only when some clause is false.
Rule (U) is modus ponens. In rule (E) the opened clause $[B(b,0);\,y.\,B(b,y)]$ is true,
because its two components have the same root. Rule (I) concludes a paired formula. With
the empty context the conclusion is true and is therefore not an equation.
\end{proof}

Adding a reflexivity rule or another seed of equations repairs the loss, as in
\cref{ctr:sec:ce,ctr:sec:three}, but that is a replacement and not a deletion. The
proposition concerns this one deletion; it is not a theorem that a presentation of $\iPRA$
needs five rules.

\section{Checker sizes}\label{ctr:sec:sizes}

This section draws on the measurements stated by \pkg{register\_graph},
\pkg{canonical\_recursors}, \pkg{clausal\_canonical}, \pkg{counter\_evaluation},
\pkg{uniform\_binding}, \pkg{uniform\_recursion}, \pkg{single\_constructor},
\pkg{constantfree} and \pkg{guarded\_register}.

Each row is a complete standalone module in the sense of \cref{base:sec:measure}: imports,
JSON decoding, validation of syntax and scope, normalization, substitution, every rule and
the final comparison are counted. All measurements are stated by the packages for CPython
3.12.14, after removal of docstrings and \texttt{ast.unparse} with one final line feed. The
role column names the new program of a package, its matched control, and the other programs
that the package reports for orientation. The programs marked $\dagger$ implement the
kernel on iterator terms and no calculus of this chapter.

\begin{longtable}{@{}p{0.55\linewidth}rrrl@{}}
\toprule
Program & Tokens & AST & Bytes & Role\\
\midrule
\endhead
\bottomrule
\endfoot
\multicolumn{5}{@{}l}{\pkg{register\_graph}}\\
Reference checker for $\mathsf T_J$ & 2814 & 3060 & 13633 & new\\
Ancillary profile of the kernel$^\dagger$ & 261 & 288 & 604 & new\\
Direct parent of the ancillary profile$^\dagger$ & 261 & 292 & 576 & control\\
Earlier leader in nodes$^\dagger$ & 263 & 289 & 613 & other\\
Earlier leader in bytes$^\dagger$ & 266 & 310 & 529 & other\\
\midrule
\multicolumn{5}{@{}l}{\pkg{canonical\_recursors}, strict family}\\
$\mathcal K_6$ on canonical seeded terms, six rules & 1338 & 1485 & 4611 & new\\
Iterator terms with successor, seven rules & 1433 & 1595 & 4923 & control\\
Raw seeded terms, seven rules & 1378 & 1532 & 4753 & control\\
\midrule
\multicolumn{5}{@{}l}{\pkg{canonical\_recursors}, compact guarded codec$^\dagger$}\\
Refined compact byte profile & 265 & 306 & 516 & new\\
Frozen compact byte control & 265 & 303 & 522 & control\\
Refined compact node profile & 263 & 280 & 602 & new\\
Frozen compact node control & 262 & 281 & 606 & control\\
Refined strict guarded checker & 542 & 624 & 1376 & new\\
Matched strict guarded control & 541 & 620 & 1381 & control\\
\midrule
\multicolumn{5}{@{}l}{\pkg{clausal\_canonical}}\\
$\mathcal K_5$, five rules & 946 & 1072 & 3305 & new\\
$\mathcal K_6$ in the same variable convention & 1107 & 1266 & 3820 & control\\
$\mathcal K_6$ as shipped by \pkg{canonical\_recursors} & 1338 & 1485 & 4611 & other\\
Six rules on zero-normal iterator terms (\pkg{zero\_normal}) & 1013 & 1075 & 3095 & other\\
\midrule
\multicolumn{5}{@{}l}{\pkg{counter\_evaluation}}\\
$\mathcal K_{\mathrm{CE}}$ & 854 & 987 & 3004 & new\\
$\mathcal K_{\mathrm{CT}}$, same terms, computation and transport & 929 & 1063 & 3272 & control\\
Scalar zero-normal, evaluation with one case & 884 & 1025 & 3061 & other\\
Raw scalar clause checker with five rules, unchanged & 853 & 993 & 3046 & other\\
The same with the known deletion of the root case & 848 & 986 & 2999 & other\\
$\mathcal K_5$ of \pkg{clausal\_canonical}, unchanged & 946 & 1072 & 3305 & other\\
\midrule
\multicolumn{5}{@{}l}{\pkg{uniform\_binding}}\\
$\mathcal K_{\mathrm{CE}}$, goal directed, one traversal & 693 & 764 & 2238 & new\\
Scalar checker with the same shared traversal & 699 & 773 & 2284 & control\\
Incoming scalar checker, exact source & 741 & 826 & 2283 & control\\
Conclusion-synthesizing checker of \pkg{counter\_evaluation} & 854 & 987 & 3004 & other\\
\midrule
\multicolumn{5}{@{}l}{\pkg{uniform\_recursion}}\\
$\mathcal K_{\mathrm{Name}}$ & 662 & 741 & 1950 & new\\
Scalar control with \rulename{Name} & 663 & 751 & 1983 & control\\
$\mathcal K_{\mathrm{Name}}$ with a budget & 710 & 800 & 2136 & new\\
Scalar control with a budget & 711 & 810 & 2169 & control\\
Strict native checker, explicit fusion$^\dagger$ & 340 & 385 & 685 & new\\
Strict native checker, integrated cases$^\dagger$ & 336 & 381 & 661 & new\\
Strict native control$^\dagger$ & 360 & 409 & 780 & control\\
Compact control with a different contract$^\dagger$ & 267 & 302 & 508 & other\\
\midrule
\multicolumn{5}{@{}l}{\pkg{single\_constructor}}\\
$\mathcal K_{\mathrm{Open}}$ & 683 & 764 & 1965 & new\\
Iterator terms with inferred opening & 683 & 771 & 1990 & control\\
\midrule
\multicolumn{5}{@{}l}{\pkg{constantfree}}\\
Reference checker for $\mathsf T_J^{\varnothing}$ & 3072 & 3349 & 13723 & new\\
Reference checker for $\mathsf T_J$, as measured by \pkg{constantfree} & 2811 & 3054 & 12626 & control\\
\midrule
\multicolumn{5}{@{}l}{\pkg{guarded\_register}}\\
Strict checker for $\mathcal K_{\mathrm G}$ & 1041 & 1181 & 3467 & new\\
Strict checker for $\Nat$ on iterator terms, same interface$^\dagger$ & 1045 & 1193 & 3566 & control\\
\end{longtable}

The comparisons that isolate a change of calculus are those against a matched control. The
checker for $\mathcal K_6$ is $312$ bytes ($6.34\%$) below the scalar control and $142$
bytes below the raw control, with normalization counted. The checker for $\mathcal K_5$ is
$515$ bytes ($13.48\%$), $161$ tokens and $194$ nodes below its matched control; the
difference of $1306$ bytes to the shipped checker for $\mathcal K_6$ combines the change of
interface with the simpler variable representation and is not to be attributed to the
deleted rule. The checker for $\mathcal K_{\mathrm{CE}}$ is $268$ bytes ($8.19\%$), $75$
tokens and $76$ nodes below $\mathcal K_{\mathrm{CT}}$, and the node overheads of the
translation are those of \cref{ctr:prop:cost}. The goal-directed checker with one traversal
is $46$ bytes, $6$ tokens and $9$ nodes below the matched scalar control and $45$ bytes,
$48$ tokens and $62$ nodes below the exact incoming checker. The checker for
$\mathcal K_{\mathrm{Name}}$ is $33$ bytes, one token and $10$ nodes below its control, and
the checker for $\mathcal K_{\mathrm{Open}}$ is $25$ bytes and $7$ nodes below its control
with the same token count. The checker for $\mathsf T_J^{\varnothing}$ is $261$ tokens, $295$
nodes and $1097$ bytes above its control: removing the constant enlarges this
implementation, through the construction of the expanded terms, the checks of anchors and
the scoped rule \eqref{ctr:eq:param}. The control is the reference checker of
\pkg{register\_graph}, but the figures that \pkg{constantfree} measures for it differ from
the figures $2814$, $3060$, $13633$ that \pkg{register\_graph} states; the packages do not
reconcile the two measurements. The strict checker for $\mathcal K_{\mathrm G}$ is $4$
tokens, $12$ nodes and $99$ bytes ($2.78\%$) below its control, which has the same JSON
interface, variable representation, logical nodes, scope checks and acceptance convention;
the guard, the state-free test and the successor-head test are counted in it. Compressed sizes are reported by the packages as supplementary
data and are not an optimization target; they are omitted here.

\section{Further questions}\label{ctr:sec:questions}

This section draws on \pkg{register\_graph}, \pkg{canonical\_recursors},
\pkg{clausal\_canonical}, \pkg{counter\_evaluation}, \pkg{uniform\_binding},
\pkg{uniform\_recursion}, \pkg{single\_constructor}, \pkg{constantfree} and
\pkg{guarded\_register}. It collects their open questions, the claims that they state
without proof, and what they report as not achieved.

\paragraph{Status of formalization.}
None of the calculi and term languages of this chapter is formalized in Lean or Rocq. The
theorems of \cref{ctr:sec:terms,ctr:sec:counter,ctr:sec:bridge,ctr:sec:constfree,ctr:sec:guarded,ctr:sec:six,ctr:sec:five,ctr:sec:ce,ctr:sec:three,ctr:sec:bounds}
are written proofs, and the executable tests of the packages are finite evidence for the
implementations. The four certificates of \pkg{constantfree} for
\cref{ctr:lem:cfconst,ctr:cor:cfzero}, and the twelve certificates of
\pkg{guarded\_register} for \cref{ctr:lem:gboot}, are checked by the Python checkers of
the packages only; \cref{ctr:thm:gequiv} rests on the formalized \cref{base:thm:bridge},
whose Lean build \pkg{guarded\_register} did not rerun. The formalized
results named in \cref{ctr:sec:six,ctr:sec:three} concern
iterator terms.

\begin{question}[Formalization]
Each term language of this chapter needs its own syntax, its normalizer where it has one,
and the bi-interpretation with iterator terms of \cref{ctr:thm:transfer}. The packages
propose the following orders of work. For \cref{ctr:thm:canonical,ctr:thm:six}: body
extensionality before the soundness of normalization, and the tests on literal sources and
witnesses after the unrestricted equivalence. For $\mathcal K_5$: the invariant on context
occurrences and the eta expansion of the closure, proved directly. For
$\mathcal K_{\mathrm{CE}}$: a deep embedding with a checked equivalence to the kernel; the
arithmetic reconstruction is small and the main work is the syntax with two binders and its
canonicalization. For normal terms: extend normalization by evaluation from iterator terms
to the counter-aware recursor, prove the environment laws
\eqref{ctr:eq:envsub}--\eqref{ctr:eq:envcomp}, and lift the comparison of the inferred
opening through \cref{ctr:lem:embedding}. For the checkers: the invariant of the sorted
walker (\cref{ctr:lem:walker}) together with the interpretation of the counter-aware
recursor. For the register: the compiler and both round trips as translations of finite
descriptions, keeping fixed syntactic bodies apart from arbitrary functions of the
metatheory. Can these be carried out by extending the existing development?
\end{question}

\begin{question}[A native register kernel]
Can a native checker with $\Q{y}$ and $\D{x}{y}$ share its verified equality premise with
the increment and load steps \eqref{ctr:eq:inc}, \eqref{ctr:eq:load}, so that the register
grammar also reduces the complete trusted source, when all conditional premises, scope
checks and syntax checks are counted? \Cref{ctr:prop:regrefl} shows compatibility of the
basis with such a kernel; no checker for it is supplied. The guarded calculus
$\mathcal K_{\mathrm G}$ of \cref{ctr:sec:guarded} is a native register kernel of a different
form: its bodies are restricted, its computations are unconditional, and the conditional
rules are derived (\cref{ctr:prop:grules}). Its strict checker is below its matched control
(\cref{ctr:sec:sizes}). For the unrestricted register with conditional rules the question
remains open.
\end{question}

\begin{question}[Guarded registers]
\begin{enumerate}
\item Can $\mathcal K_{\mathrm G}$ and \cref{ctr:thm:gequiv} be formalized on top of the
existing development of $\Nat$? \pkg{guarded\_register} gives the equivalence as a written
proof only.
\item Is there a complete checker for $\mathcal K_{\mathrm G}$, or for a variant of it, below
the compact envelope of $256$ tokens, $276$ nodes and $506$ bytes? That envelope combines
three different programs, a token profile of \pkg{kernel\_refinement} ($256$, $283$,
$594$), a node profile of \pkg{deterministic\_opening} ($263$, $276$, $580$) and a byte
profile of \pkg{deterministic\_introduction} ($267$, $303$, $506$), which share a guarded
public-language contract that differs from the complete validation of the strict checkers.
The package states that a claim of this kind has to implement the whole certificate
contract, charge the guard and every scope operation, and compare complete programs under
the same normalization and acceptance conventions.
\item \pkg{guarded\_register} reports two shortcuts that it examined and did not use, and
proves neither claim. Restricting both the introduction of a witness and the specialization
of a clause at once is said to need a joint completeness argument, since the separate
completeness arguments may each depend on the unrestricted form of the other rule. Removing
the scope information before the phase that chooses fresh names is complete is said to
admit future witnesses.
\end{enumerate}
\end{question}

\begin{question}[The register selector and the cost of the counter]
Is there a substantially shorter numerical selector for the register basis than the term
$C$ of \cref{ctr:lem:regboot}, which has a short proof and a large expansion? Source
brevity, term size and certificate size are different objectives. Does the register with
two bound variables, or the one-variable register of \cref{ctr:prop:scalarreg}, yield the
smaller complete checker after comparable implementations?
\end{question}

\begin{question}[Weakening the graph axiom]
Which parts of \eqref{ctr:eq:STEP} can be weakened while actual equality and the formula
round trips on the whole domain are retained? \pkg{register\_graph} states without proof
that one-sided operational membership and agreement on the standard model are insufficient.
\pkg{constantfree} states in the same way, without proof, that replacing the biconditional
of $\mathsf T_J^{\varnothing}$ by an unconditional evaluation equation gives a different
theory, and asks whether fewer or cheaper inference patterns can supply both the
computation and the coverage that the graph schema provides.
\end{question}

\begin{question}[The constant-free presentation]
\begin{enumerate}
\item Can the anchored schemas of \cref{ctr:def:TJ0} be compiled into a native calculus with
$\Q{y}$ and $\D{x}{y}$, without a large selector, normalizer or second certificate
language, and what is the size of the complete standalone checker, including validation of
the public grammar?
\item Can a certificate protocol recover the anchor of each arithmetic instance from a
numerical parameter already in scope, instead of naming it? By
\cref{ctr:prop:cfparam} an empty scope still needs \eqref{ctr:eq:param} or a definable
constant.
\item What is the least complete cost of eliminating the constant? \Cref{ctr:prop:cfsize}
fixes the size of the zero term, but not the cost of the parser, the scoping discipline, the
expansion of the abbreviations, the rules and the interpretation.
\item Can \cref{ctr:thm:cfbridge}, in particular the equivalence with $\mathsf T_J$ and the
definition of zero by $\operatorname{Zero}$, be formalized in Lean or Rocq? The register
theory $\mathsf T_J$ itself is not among the formalized syntaxes, so that dependency has to
be formalized as well.
\item Is there one checker, for any presentation, below every coordinate of the currently
observed compact minima, with its input contract and the theory it recognizes stated?
\end{enumerate}
\end{question}

\begin{question}[Zero-equation induction]
Which principle beyond induction on zero equations excludes the two-point fibers of
\cref{ctr:thm:zeroind}? Reflection of actual equality suffices, at a cost that has to be
counted. Does a narrowly stated uniqueness or extensionality principle for the identity
iterator already give full induction on equations, and what does it cost in a complete
checker? Does $\mathsf T_0$ interpret arithmetic on a quotient, or admit a bi-interpretation
with $\iPRA$ on a changed domain? The countermodel decides neither.
\end{question}

\begin{question}[Successor rigidity]
The package on raw clause calculi with ordinary substitution states a restricted
successor-rigidity principle that yields induction on zero equations and leaves a faithful
bootstrap of equality and predecessor unresolved. Does that narrower principle suffice?
\pkg{counter\_evaluation} records the question and does not use the unproved conclusion.
\end{question}

\begin{question}[The self-seeded language]
Can \cref{ctr:thm:unary} be strengthened to a normal form or an effective classification of
the functions defined by self-seeded terms? The proof gives the existence of a bound and no
algorithm that extracts an optimal bound from a term. Is there an obstruction for
interpretations by graphs or on changed domains?
\end{question}

\begin{question}[Tests on literal syntax]
The witness-directed introduction is concise but not stable under substitution: a
substituted witness can become the literal $0$. Is there a substitution-stable presentation
with a comparable complete checker? In particular, guarded atomic matrices make the
introduction depend on the root of the formula only and guarantee a reflexive zero row for
the opening \eqref{ctr:eq:open}; is the resulting complete grammar with its checker simpler
than the unrestricted $\mathcal K_{\mathrm{Open}}$? Is there a uniform counter rule on raw
syntax, with the zero computation stated explicitly instead of placed in the canonical
substitution, whose checker is of comparable size?
\end{question}

\begin{question}[Compact encodings]
Can an encoding of canonical seeded terms absorb the erasure of zero-counter bodies into a
traversal that is needed anyway and improve the compact source minima, while validating the
complete public language? Does the grammar with disjoint arities, without a successor
constructor and without a closure constructor, yield a smaller fully validated compact
program once normalization and every case of malformed input are counted? Can the
combination of the counter-aware arithmetic with the three families be ported to the compact
certificate architecture below the existing profiles? Can a uniform representation remove
the partial opening of binders, given that the recursor and the paired binder bind two
variables and a clause one, while reducing the whole checker? Can the two instantiations of
the counter body in \eqref{ctr:eq:CE} be given a simpler complete scope discipline? Can a
pointfree representation of proofs reduce the substitution machinery without introducing
pairing structure, when the typing and conversion costs of external tuple contexts are
counted?
\end{question}

\begin{question}[Proof size]
What are the optimal translations of proof graphs between $\mathcal K_6$ and
$\mathcal K_5$, and how much sharing avoids repeated neutral premises? Can a comparison of
proof complexity between counter evaluation and the evaluation rule on the state of an
iterator quantify certificate lengths and label sizes in addition to the inference nodes of
\cref{ctr:prop:cost}? Should checkers with an explicit budget be compared only with
controls that carry the same budget?
\end{question}

\begin{question}[Measures and lower bounds]
The coordinates used in \cref{ctr:sec:sizes} and the counts of constructors, rules and
proof expansion disagree with one another. Is there a robust measure that compares
different proof interfaces? Can a lower bound be proved for a restricted class of checkers,
with the term grammar, the class of decidable operations on syntax, the permitted runtime
primitives and the certificate interface fixed in advance? A count of named rules alone
changes too easily, by merging cases or by moving work into the syntax.
\end{question}

\begin{question}[A new arithmetic reduction]
Eliminating the initializer, changing the domain or replacing the recursor may still give a
simpler basis. A candidate has to respect the self-seeded constructions and the obstruction
of \cref{ctr:thm:unary}, and it has to supply both interpretation composites, not only
expressiveness on the standard model.
\end{question}

\paragraph{Not achieved.}
The packages state the following limits.
\begin{itemize}
\item No program of \cref{ctr:sec:sizes} for a calculus of this chapter is below the compact
codecs of the kernel ($516$ and $508$ bytes), and no minimality of an axiomatization or of
a checker is claimed. An observed minimum among programs is an upper bound on the optimum.
\item The register theory does not improve the size of a complete checker: its graph axiom
is richer, its body has two bound variables, and its reference checker has $13633$ bytes.
The ancillary profile lowers the node count by one and is larger in bytes.
\item The refined compact byte profile does not improve the token count of its control, and
the token frontier of $258$ reported by \pkg{canonical\_recursors} is not beaten.
\item The checker for $\mathcal K_5$ has fewer rules and tokens but more bytes than the
six-rule checker on zero-normal iterator terms of \pkg{zero\_normal}; its positive claim in
bytes is the matched comparison.
\item The checker for $\mathcal K_{\mathrm{CE}}$ is smaller in all three coordinates than
the scalar control with one evaluation case, but the raw scalar clause checker with the
root case deleted is five bytes smaller.
\item For the checker of \pkg{uniform\_binding} the compressed size is larger than that of
the exact incoming scalar checker.
\item Two directions were examined in \pkg{uniform\_binding} and discarded. A formula of the
clause syntax is a chain of paired binders ending in an equation and can be flattened to a
triple of a depth and two terms; the change is exact, but the implementation grew, because
checking of ranks and opening of endpoints outweighed the saved traversal. A rule of
parameter substitution used only in the reverse direction is sound, but the intended
completeness argument by a fresh alias fails: once a goal no longer mentions the alias, the
rule can require that same goal as its premise.
\item The translator of test fixtures of \pkg{uniform\_recursion} declines an inherited
fixture whose witness becomes the literal $0$ under normalization; a direct proof of the
same formula is supplied, and no universal executable proof compiler is claimed.
\item Removing the constant reduces the numerical syntax and not the checker: the reference
checker for $\mathsf T_J^{\varnothing}$ ($3072$ tokens, $3349$ nodes, $13723$ bytes) is
larger in every coordinate than its control. No minimality of the axiomatization is
claimed, \cref{ctr:prop:cfsize} is not a bound on checker size, and the measurement of the
control differs from the one stated by \pkg{register\_graph} (\cref{ctr:sec:sizes}).
\item The strict checker for $\mathcal K_{\mathrm G}$ ($1041$ tokens, $1181$ nodes, $3467$
bytes) improves only its matched control; it is larger in every coordinate than each of the
three compact programs named in the question on guarded registers, and \pkg{guarded\_register} does not
reach its stated target of a program below their envelope. Its survey of the existing
programs measured the $442$ distinct Python sources of at most $1500$ source bytes in the
ninety package trees; the threshold limits the audit, and it is not a lower bound for
larger sources that normalize to smaller programs. No minimality of the axiomatization or of
the checker is claimed, and a smaller number of constructors or rule names is not by itself
a smaller checker.
\item Concrete limits of the interpreter on recursion depth and memory can reject a valid
large certificate for every checker of the chapter; completeness is a statement about the
calculi and about the checkers without such limits.
\end{itemize}
